% concatenated from appleby_lawless_pantograph7.tex
%------------------------------------------------------------------------------
% filename: 
% date: 
% last revised: 
% authors: 
%------------------------------------------------------------------------------
\documentclass[a4paper,10pt,reqno]{amsart}
\usepackage{amssymb,latexsym,bbm,amsmath,amsfonts}
\usepackage{amsthm}
\usepackage[mathscr]{eucal}
\usepackage{rotating}
\usepackage{multirow}
%\usepackage{slashbox}
%\usepackage[notref,notcite]{showkeys}

\numberwithin{equation}{section}
\newtheorem{theorem}{Theorem}
\newtheorem{lemma}{Lemma}
\newtheorem{proposition}{Proposition}
\newtheorem{corollary}{Corollary}
\newtheorem{conjecture}{Conjecture}

\theoremstyle{definition}
\newtheorem{definition}[theorem]{Definition}
\newtheorem{assumption}[theorem]{Assumption}
\newtheorem{example}[theorem]{Example}

%\theoremstyle{definition}
%\newtheorem{definition}{Definition}
%%\newtheorem{example}[theorem]{Example}
%%\newtheorem{xca}[theorem]{Exercise}
\theoremstyle{remark}
\newtheorem{remark}{Remark}
%\newtheorem{example}{Example}

\newcommand {\Cov}{\mbox{Cov}}
\newcommand {\Var}{\mbox{Var}}
\newcommand {\Corr}{\mbox{Corr}}
\newcommand {\diag}{\mbox{diag}}
\newcommand {\sgn}{\mbox{sgn}}
\renewcommand{\o}{{o}}                               % small Landau

\newcommand{\norm}[1]{\left\lVert #1 \right\rVert}
\newcommand{\abs}[1]{\left\lvert #1 \right\rvert}    % absolut value

\DeclareMathOperator{\1}{\mathbbm 1}
\DeclareMathOperator{\F}{{\cal F}}                   % sigma-fiel
\DeclareMathOperator{\R}{{\mathbb R}}                % reals
\DeclareMathOperator{\Rp}{{\mathbb R}_+}             % positive reals
\DeclareMathOperator{\C}{{\mathbb C}}                % complexe values
\DeclareMathOperator{\N}{{\mathbb N}}                % integer
\DeclareMathOperator{\Id}{Id}                        % identity
\DeclareMathOperator{\I}{I}                          %
\renewcommand{\P}{P}


\begin{document}
	
	\title[Asymptotic behaviour of pantograph equations]
	{Characterisation of asymptotic behaviour of perturbed deterministic and stochastic pantograph equations}
	
	\author{John A. D. Appleby}
	\address{School of Mathematical
		Sciences, Dublin City University, Glasnevin, Dublin 9, Ireland}
	\email{john.appleby@dcu.ie} 
	
	\author{Emmet Lawless}
	\address{Department of Mathematics, 
		University of Michigan, 
		Ann Arbor, MI 48109-1043, U.S.A.}
	\email{elawless@umich.edu}
	
	\thanks{JA is supported by the RSE Saltire Facilitation Network on Stochastic Differential Equations: Theory, Numerics and Applications (RSE1832). EL was supported by Science Foundation Ireland (16/IA/4443).} 
		\subjclass{34K06, 34K12, 34K20, 34K25, 34K50, 60H10}
%	\subjclass{34D05; 34D20; 93D20; 93D09}
%	\keywords{asymptotic stability,
%		global asymptotic stability, fading perturbation, asymptotic preserving functions, regular variation, dominated variation, decay rates, polynomial asymptotic stability}
	\date{24 September 2025}
	
	\begin{abstract}
This paper considers the asymptotic behaviour of  deterministically and stochastically forced linear pantograph equations. The asymptotic behaviour is studied in the case when all solutions of the  pantograph equation without forcing tend to a trivial equilibrium. 
In all  cases, we  give necessary and sufficient conditions on the forcing terms which enable all solutions to converge to the equilibrium of the unforced equation, and which enable solutions to remain bounded. In the deterministic case, we give sharp conditions on forcing terms which enable the solutions of the forced equations to inherit the power law behaviour of the unforced equation, as well as slower rates of decay or growth not present in the unforced equation. Extensions to equations with general unbounded delay, and to finite--dimensional equations are also presented.
	\end{abstract}
	
	\maketitle
	
	\section{Motivation and Scope of the Work}
	\label{sec:2}
	
	Let $a\neq 0$, $b$ be real numbers, and let $q\in (0,1)$. The (scalar valued solution of the) delay differential equation  
	\begin{equation}\label{eq.unforced}
		u'(t)=bu(t)+au(qt), \quad t\geq 0; \quad u(0)\neq 0,
	\end{equation}
	has attracted a lot of attention. The is a so--called pantograph equation, first studied in depth by Kato and McLeod~\cite{KatoMc} and Fox et al and Ockendon and Tayler~\cite{Foxetal,OckTay}, with finite dimensional analysis in \cite{CarrD1,CarrDyson2}. Apart from interest in applications, it attracts mathematical interest as it is an example of a delay differential equation with unbounded delay, and, despite the constant multipliers $a$ and $b$, the equation is also non--autonomous, by virtue of the ``proportional'' delay $qt$. One interesting feature of linear differential equations with unbounded delay is that solutions that tend to an equilibrium need not do so at an exponentially fast rate (see e.g., \cite{AppBuck10,Kris88,MakTerj,Diblik98,DibRuz}). Pantograph equations are no exception, and Kato and McLeod \cite{KatoMc} showed that solutions which tend to zero do so at a power law rate of decay. Further treatments of deterministic equations with unbounded delay include \cite{cermak, HaddKris, Iser, Liu}, and work continues to date e.g. \cite{cer2022}. Fractional and integrodifferential equations of pantograph type have also been studied \cite{MR4534554}.
	
	There has been a lot of activity in extending results to stochastically perturbed pantograph equations and equations with unbounded delay, especially in the direction of asymptotic behaviour and numerical methods (see e.g.\cite{App2005, BakBuck2000,BakBuck2005, Fan,Huang,Huetal,Jiang,LiuChen,Meng1,Meng2,PJ2012, YanHuang,Yangetc}). The first such work on asymptotic behaviour is \cite{AppBuck2016}. In this field of stochastic asymptotic analysis, papers concerning power--law behaviour and numerical methods \cite{MR4439350,MR3927238,Milo}, large deviations \cite{MR4585696}, and neutral equations \cite{MR4278069} and L\'evy noise \cite{MR4053553} have also been recently considered. 
	%, and many other works on stochastic equations with unbounded delay exist. 
	Many \textit{sufficient conditions} have been obtained with respect to which solutions of stochastic equations exhibit convergence to a limit, either in $p$--th moment or pathwise. However, what appears to be absent in this literature are results which \textit{characterise} certain types of asymptotic behaviour in terms of forcing terms, such as, for instance, convergence to the equilibrium of the unforced equation, or asymptotic rates of growth or decay of forced solutions; it is one of the chief contributions of this paper to supply such asymptotic characterisations. But as we indicate now, we believe that the techniques in this paper have much wider application than to pantograph equations. 
	
	In this paper, we consider perturbations of \eqref{eq.unforced} where the intensity of the forcing terms is state--independent. Both deterministic and stochastic equations are considered, and in the case of stochastic equations, we allow the solution to be forced both deterministically and stochastically. The main contribution of the work is to give a  \textit{characterisation} of certain convergence, boundedness, and growth estimates in terms of the forcing terms. By a  characterisation of an asymptotic property \textbf{A} of solutions, we mean that solutions possess property \textbf{A} holds if and only if some conditions \textbf{C} are satisfied by the forcing terms; moreover, the conditions \textbf{C} should be expressible without recourse to properties of the solution of the forced equation, or other auxiliary objects that cannot be directly computed from the problem data. Ideally, such conditions should be easy to check. The characterisations we give satisfy these natural requirements.   
	
	The techniques used in this paper are especially suited to considering such ``additive'' rather than ``multiplicative'' perturbations (the terminology and analysis is more commonplace for stochastic equations, but is appropriate for deterministic equations too). Stochastic equations with such state--dependent noise have been studied in the papers listed above. Our analysis of additive noise equations relies on  a decomposition approach to solutions of stochastic delay differential equations, which allows such solutions to be built from solutions of stochastic ordinary differential equations with non--differentiable paths, and solutions of random functional differential equations which are continuously differentiable. A similar approach can be applied to equations with multiplicative noise, and a sketch of possible research in that direction is given in Section 6.  
	
	Other generalisations, to scalar equations with a general unbounded delay, or to multidimensional equations, are considered in Section 6, and sketch proofs are given which show that at least basic convergence and boundedness results still hold, giving reassurance that the theory is not confined to the narrow class of scalar, proportional delay equations. However, the main thrust of our paper concerns pantograph equations, partly due to their celebrated role in the historical development of the theory of unbounded delay equations, and partly because, since the asymptotic behaviour of the unperturbed pantograph equation is essentially perfectly understood (in contrast to many other unbounded delay equations) we can exploit this knowledge to obtain equally pristine results in the perturbed case, thereby demonstrating the sharpness of our approach. 
	
	In doing so, we are able to improve decay estimates on perturbed equations, show that these estimates are optimal in certain cases, accurately characterising when perturbation sizes cause solutions to transition from preserving the asymptotic behaviour of the unperturbed equation to slower rates of decay (or even growth) when the perturbation are larger than a critical size. In fact, by proving both big-$O$ and little-$o$ growth estimates on solutions throughout, we are able to characterise in many cases (in terms of an easily estimated functional of the perturbation, which is entirely independent of the structure of the unperturbed equation) when the solution of the perturbed equation has a given order of growth (or decay). The combination of (a) sharp estimation of growth rates (b) corresponding  correct determination of the transition point from the small to the large perturbation regime, and (c) characterisation of the growth conditions on the forcing terms which give rise to certain rates of growth or decay in solutions, are to the best of our knowledge new in the asymptotic study of unbounded delay differential equations. 
	
	Another novel feature in our analysis is to be able, for functional differential equations which are simultaneously deterministically and stochastically perturbed, to characterise certain important asymptotic behaviours. In particular, we are able to characterise exactly when solutions of perturbed equations are convergent or bounded, and the characterising conditions are very easy to check. Such a characterisation, with only continuity side conditions on forcing terms being required, seems new to the literature, \textit{even for non--delay stochastic differential equations}. Better yet, we sketch at the end of Section 5 an argument which suggests strongly that \textit{such a characterisation might well apply to a very broad class of functional differential equations} which can have finite or unbounded delay, and which does not possess any of the special structure of the pantograph equations which we study here. 
	
	We note that a great deal of the sharpness achieved in our results comes from a very careful asymptotic characterisation of solutions of scalar linear \textit{ordinary} and stochastic differential equations. Remarkably, given the simplicity of the equations involved, such a characterisation does not seem to be present in the literature, and part of the purpose of this paper is to show that a thorough \textit{development of this theory for linear ordinary differential equations will unlock a characterising perturbation theory for a very wide class of functional differential equations}. 
	
	We have chosen perturbed pantograph equations to do this, because the perturbation theory of non--autonomous delay differential equations (of which the pantograph equation is a well--studied example) has traditionally found comparatively little use for resolvents or fundamental solutions, at least compared with autonomous equations. Indeed, the analysis in this paper likewise makes no use whatsoever of resolvents or variation of constants formulae. In itself, this is neither noteworthy nor surprising. However, in other works, the authors have given asymptotic characterisation results for perturbations of autonomous deterministic and stochastic functional and Volterra equations, in which deterministic and stochastic variation of constants formulae are in constant use. This might leave one to believe that such formulae play a critical role in making the perturbation theory both achievable and sharp. \textit{However, this paper shows that while these resolvent--based representations are certainly useful when available, they do not appear to be essential for the development of a sharp perturbation theory}. As a result, the pantograph equation represents an interesting test case for a sharp perturbation theory that applies to other non--autonomous or even non--linear functional differential equations.         
	
	\subsection{Technical introduction}
	To make some of this general discussion precise, we let $f$ and $\sigma$ be continuous deterministic functions, $B$ be a standard Brownian motion, and $a$ and $b$ be real numbers obeying 
	\begin{equation} \label{eq.unforcedasstable}
		b<0,  \quad |b|>|a|,
	\end{equation}
	and consider 
	\begin{gather*}
		x'(t)=ax(qt)+bx(t)+f(t),\\
		dX(t)=(aX(qt)+bX(t)+f(t))\,dt+\sigma(t)\,dB(t).
	\end{gather*}
	\eqref{eq.unforcedasstable} is a  necessary and sufficient condition for all solutions of \eqref{eq.unforced} to tend to zero as $t\to\infty$ \cite{KatoMc}. Throughout we take $a\neq 0$, since this case deals with non--delayed equations. In the case of the deterministic equation, we give necessary and sufficient conditions on $f$ for which solutions will tend to zero, or for which solutions are bounded. 
	%	It is notable, perhaps even remarkable, that the conditions on $f$ and $\sigma$ which characterise when solutions of perturbed equations are bounded or convergent are completely independent of the structure of the unperturbed equation \eqref{eq.unforced}. 
	We also characterise when solutions $x$ are unbounded with a certain order of growth, or decay to zero more slowly than those of \eqref{eq.unforced}, and give sharp conditions under which the decay rate of $x$ is the same as that of \eqref{eq.unforced}. For the stochastic equation, we give conditions on $f$ and $\sigma$ which characterise when $X(t)\to 0$ as $t\to\infty$ a.s. and when $X$ is almost surely bounded. Results for multi--dimensional equations are also sketched. In later works, we will develop quantitative pathwise estimates for $X$, and relax the second inequality in \eqref{eq.unforcedasstable}, so as to study the situation when the underlying unforced equation \eqref{eq.unforced} has unbounded solutions. 
	
	In short, therefore this paper shows that provided all solutions of \eqref{eq.unforced} tend to zero, we can find coincident necessary and sufficient conditions on the forcing terms $f$ and $\sigma$ (which do not depend on the parameters $b$, $a$ and $q$) under which the solutions are convergent to zero; in particular, it is not necessary to strengthen conditions on the underlying equation in order to preserve stability, merely to impose appropriate smallness conditions on $f$ and $\sigma$. 
	
	The results in this paper develop work of the authors for certain types of \textit{uniformly asymptotically stable} and \textit{autonomous} linear differential and Volterra equations which are forced both deterministically and stochastically by time--independent forcing terms.  The nature of conditions we impose on $f$ and $\sigma$ for pantograph equations are identical to those needed in the autonomous case, and we avail of existing results from \cite{AppLaw2024ODE} for ordinary differential equations in our proofs. 
	%In \cite{AL:2024Lp}, if $\nu$ is a finite measure on $[0,\infty)$, it is shown that if the underlying differential resolvent of 
	%\begin{equation} \label{eq.VolterraSFDE}
	%	dX(t)=\left(\int_{[0,t]} \nu(ds)X(t-s)+f(t)\right)\,dt + \sigma(t)\,dB(t)
	%\end{equation}
	%is absolutely integrable, then for $p\geq 2$, $X\in L^p(\mathbb{R}^+)$ a.s. if and only if 
	%\[
	%\int_{(t-\theta)^+}^t f(s)\,ds \in L^p(\mathbb{R}^+), \quad \text{for all $\theta>0$} \quad \int_{t-1}^{t}\sigma^2(s)\,ds \in L^{p/2}(\mathbb{R}^+),
	%\]
	%with a related, but more complicated, characterisation for $p\in[1,2]$. 
	%This develops an $L^p$ characterisation for $p\geq 1$ in the deterministic case (i.e. \eqref{eq.VolterraSFDE} with $\sigma=0$) from  \cite{AL:2023(AppliedMathLetters)}. 
	Related results for the mean square asymptotic behaviour of affine equations, where there is also state--dependence in the diffusion coefficient, are given in  \cite{AL:2023(AppliedNumMath)}. For the stochastic differential equation 
	$
	dY_0(t)=-Y_0(t)\,dt +\sigma(t)\,dB(t)$ it is shown in \cite{ACR:2011(DCDS)} that  $Y_0(t)\to 0$ as $t\to\infty$ a.s.  if and only if $S(\epsilon)<+\infty$ for all $\epsilon>0$, where 
	\begin{equation} \label{eq.S}
		S(\epsilon):= \sum_{n=1}^\infty \sqrt{\int_{n-1}^n\sigma^2(s)\,ds} \exp\left(-\epsilon\int_{n-1}^{n}\sigma^2(s)\,ds\right).
	\end{equation}
	In \cite{AppLaw2024ODE} (as discussed in Section 4), it is shown that solutions of the ordinary differential equation $y'(t)=-y(t)+f(t)$ obey $y(t)=O(\gamma(t))$ as $t\to\infty$ for non--decreasing and so--called subexponential weight functions $\gamma$ if and only if $f_\theta$, defined by  
	$f_\theta(t):=\int_{(t-\theta)^+}^t f(s)\,ds$ for $(t,\theta)\in [0,\infty)\times[0,1]$, obeys
	\begin{equation} 
		\label{eq.fthetagammabdd}
		\sup_{0\leq \theta\leq 1} |f_\theta(t)| \leq K\gamma(t), \quad t\geq 0, \quad \text{for some $K>0$}.
	\end{equation}
	(In the above, and throughout the paper, we make extensive use of the standard ``big O'' and ``little o'' Landau notation, as well as the notation $a(t)\sim b(t)$, $t\to\infty$ for real--valued functions $a$ and $b$ to signify that $a(t)/b(t)\to 1$ as $t\to\infty$).
	Finally, in this volume, we also show that the asymptotic behaviour of $f_\theta$  characterises the situation when solutions of deterministic functional and  Volterra differential equations have time averages (see \cite{AL:2024Cesaro}).  
	
	Therefore, we see that the asymptotic behaviour of $f_\theta$ for $\theta\in [0,1]$ and $\sigma^2_1(t):=\int_{(t-1)^+}^t \sigma^2(s)\,ds$ appear to characterise many interesting types of asymptotic behaviour of solutions of forced autonomous equations. In this work, we show that this is true for perturbed pantograph equations also. \textit{Therefore, the evidence accumulates in favour of the overall viewpoint that situations in which important qualitative and quantitative asymptotic properties of solutions of perturbed differential systems (whether deterministic or stochastic, memory--dependent or not, autonomous or not, and perhaps even linear or not), may be characterised by in terms of perturbations through $\{f_\theta:\theta\in [0,1]\}$ and $\sigma^2_1$}.  
	
	%	\section{Fundamental Estimate}
	%	
	%	In this section, we establish an estimate on the size of solutions of perturbed pantograph equations which enables us to obtain not only convergence and stability results, but also (after further analysis later in the paper) very sharp asymptotic bounds, especially when perturbations are, in a certain sense, small.  
	
	We recall some facts about the asymptotic behaviour of solutions of  \eqref{eq.unforced}; all results appear in \cite{KatoMc}. There is no loss of generality in taking the initial condition to be unity; denoting the solution of the differential equation with initial condition $\xi\in \mathbb{R}$ by $u(\cdot;\xi)$, we have that 
	$u(t;\xi)=u(t;1)\xi$ for all $t\geq 0$. Notice that if $\xi=0$, then $u(t;\xi)=0$ for all $t\geq 0$, so that the zero function is an equilibrium solution for all choices of the parameters (when $a+b=0$, every constant is a solution, but we do not study this phenomenon). 
	
	If $b>0$, $u(t)/e^{bt}$ tends to a finite limit, so solutions exhibit real exponential growth: we will not study this case here. When $b<0$, however, power--law asymptotic behaviour occurs. Specifically, since $a\neq 0$ and $b<0$ we may define 
	\begin{equation} \label{eq.kappa}
		\kappa:=-\frac{\log|b/a|}{\log(1/q)},
	\end{equation}
	and $u(t)=O(t^\kappa)$ as $t\to\infty$. Indeed, there are no non--trivial solutions that are $o(t^\kappa)$ as $t\to\infty$, so this estimate is very sharp \cite[Theorem 3]{KatoMc}. More precise information regarding oscillation of solutions are known, but since we focus on sharp upper estimates of solutions of perturbed equations, we do not seek such refined results. 
	
	This discussion explains in part why we wish to confine attention in this paper to the case when \eqref{eq.unforcedasstable} holds, but does not mean that it is not interesting, or impossible, to give a perturbation theory in the other parameter regimes mentioned above, and it is a task that we will return to in other works. 
	
	%	If $b<0$, $\kappa<0$, $\kappa=0$, or $\kappa>0$ according as $|b|>|a|$, $|b|=|a|$ or $|b|<|a|$. Thus in the first case, $u(t)\to 0$ as $t\to\infty$. In fact, the conditions $b<0$, $|b|>|a|$ \textit{characterise} when $u(t)\to 0$ as $t\to\infty$ (as pointed out, if $b>0$, $u(t)$ grows exponentially, and is also unbounded if $b=0$). Thus, we confine attention to the case when 
	%	\begin{equation} \label{eq.unforcedasstable}
		%		b<0,  \quad |b|>|a|. 
		%	\end{equation}
	%	
	%	In this paper we consider both deterministic and stochastic perturbations of \eqref{eq.unforced} (sometimes deterministic and stochastic forcing occurs simultaneously). We focus on deterministic case first, partly for simplicity, and partly because stochastic proofs case can be reduced to deterministic ones by arguing pathwise in the sample space. 
	
	%	We study the deterministic pantograph equation first. Suppose that 
	%	\begin{equation} \label{eq.fcns}
		%		f\in C([0,\infty);\mathbb{R}),
		%	\end{equation} 
	%	and consider the forced pantograph equation
	%	\begin{equation} \label{eq.forced}
		%		x'(t)=ax(qt)+bx(t)+f(t), \quad t\geq 0; \quad x(0)=\zeta\in \mathbb{R}.
		%	\end{equation}
	%	There is a unique continuous solution to this initial value problem (see e.g.~\cite{AppBuck10}). If \eqref{eq.unforcedasstable} holds, so $u(t)\to 0$ as $t\to\infty$, we ask  under what conditions on the forcing function $f$ certain asymptotic properties (such as convergence to zero, or rates of decay) of $u$ are preserved. 
	%	%On the other hand, when an underlying deterministic differential system is ``stable'', and it is perturbed by ``large'' perturbations, solutions of the forced system can inherit properties of the forcing terms. 
	%	%Conversely, if solutions of the forced equation exhibit certain ``large'' asymptotic behaviour, it it would be agreeable to be able to trace this behaviour back to an asymptotic property of the forcing term. 
	%	On the other hand, we also prove results which give necessary and sufficient conditions on forcing terms for the solution to exhibit decay or even growth bounds which differs from those of \eqref{eq.unforced}. 
	%	%certain types of ``large'' asymptotic behaviour, including explicit growth and decay bounds. 
	
	\section{A Fundamental Estimate and Large Perturbation Estimates}
	
	We study the deterministic pantograph equation first, partly for reasons of simplicity, but also because solutions of the stochastic equation can be written in terms of the deterministic equation, when viewed path--by--path in the sample space. Therefore the analysis of the deterministic equation is fundamental. Suppose that 
	\begin{equation} \label{eq.fcns}
		f\in C([0,\infty);\mathbb{R}),
	\end{equation} 
	and consider the forced pantograph equation
	\begin{equation} \label{eq.forced}
		x'(t)=ax(qt)+bx(t)+f(t), \quad t\geq 0; \quad x(0)=\zeta\in \mathbb{R}.
	\end{equation}
	There is a unique continuous solution to this initial value problem (see e.g.~\cite{AppBuck10}). If \eqref{eq.unforcedasstable} holds, then  $u(t)\to 0$ as $t\to\infty$. It is therefore reasonable to ask under what ``smallness'' conditions on the forcing function $f$ certain asymptotic properties (such as convergence to zero, or rates of decay) of $u$ are preserved; on the other hand, we may be interested in understanding the properties of solutions when $f$ does not obey smallness   properties which preserve properties of $u$.  
	
	To this end, in this section, we establish a ``fundamental'' estimate on the size of solutions of perturbed pantograph equations which enables us to obtain not only convergence and stability results, but also (after further analysis later in the paper) very sharp asymptotic bounds, especially when perturbations are, in a certain sense, small.
	The results are of a ``sufficiency'' type, which is to say, if sufficiently strong asymptotic conditions are imposed on the forcing term, this will generate a rate of decay on the solution. A full characterisation of asymptotic behaviour of solutions will require ``necessity'' results, as well as some transformations and decomposition results; this is because our fundamental lemma is usually applied to a perturbed pantograph equation which results from such a decomposition, rather than to the original pantograph equation itself. We will also need estimates which deal with large perturbations, and these will be proved separately. Therefore, our fundamental estimate, although a key tool in the asymptotic analysis, provides only part of the picture. 
	
	This approach differs from our analysis of autonomous differential and Volterra integrodifferential equations, where we exploit variation of constants formulae and known (integrable) behaviour of the differential resolvent. Although perturbed pantograph equations admit such variation of constants formulae, the non--autonomous character of the equation makes it difficult to get good asymptotic information about the underlying resolvent, and this motivates an approach avoiding this problem. In this sense, our approach bears relation to that taken in \cite{AppBuck10}.      
	
	In order to state our fundamental estimate, we introduce some notation. For $a\neq 0$, there are (in general, complex valued) solutions $k$ to the equation $aq^k+b=0$. If $k_0\in \mathbb{C}$ is such a solution,  $\kappa:=\text{Re}(k_0)$ is uniquely defined, and given by \eqref{eq.kappa}.
	Let 
	\begin{align} \label{eq.c}
		c&:=\log q<0,\\
		\label{eq.alpha}
		\alpha&:=|a/b|,
	\end{align}
	so that $e^c=q$, and under \eqref{eq.unforcedasstable}, $\alpha\in (0,1)$. 
	
	\begin{lemma} \label{lemma.fundamental}
		Suppose $q\in (0,1)$ and \eqref{eq.unforcedasstable} holds. Let $\varphi\in C([0,\infty);\mathbb{R})$ and $z$ be the unique continuous solution to 
		\begin{equation} \label{eq.z}
			z'(t)=az(qt)+bz(t)+\varphi(t),\quad t\geq 0; \quad z(0)=\xi\in\mathbb{R}.
		\end{equation}
		Let $\alpha$, $c$ be defined by \eqref{eq.alpha}, \eqref{eq.c} and define $I_n:=[s_0-nc,s_0-(n+1)c]$. Let  
		\begin{equation} \label{eq.Kn}
			K_n:=\sup_{s\in I_n} e^{-\kappa s}|z(e^s)|,
		\end{equation}
		and let $\epsilon_n>0$ be a sequence such that  
		\begin{equation} \label{eq.epsn}
			\sup_{s\in I_{n+1}} |\varphi(e^s)|\leq \epsilon_{n+1}.
		\end{equation}
		\begin{enumerate}
			\item[(i)] If $(\epsilon_n/\alpha^n)$ is summable, $K_n\leq K^\ast$, and  $z(t) = \text{O}(t^\kappa)$ as $t\to\infty$.  
			%	\[
			%\sup_{t\in [e^{s_0-nc},e^{s_0-nc}e^{-c}]} |z(t)|\leq K^\ast [e^{(s_0-nc)}]^\kappa,
			%\] 
			%			Suppose $\lambda_n$ is summable. Then $K_n\leq K^\ast$, so as $\kappa<0$, we have 
			%		\[
			%		e^{-\kappa(s_0-nc)}\sup_{s\in [s_0-nc,s_0-(n+1)c]} |z(e^s)|
			%		\leq \sup_{s\in [s_0-nc,s_0-(n+1)c]]} e^{-\kappa s} |z(e^s)|\leq K^\ast,
			%		\]
			%		so 
			%		\[
			%		\sup_{t\in [e^{s_0-nc},e^{s_0-nc}e^{-c}]} |z(t)|\leq K^\ast [e^{(s_0-nc)}]^\kappa,
			%		\]  
			\item[(ii)] If $(\epsilon_n/\alpha^n)$ is not summable, then $
			%\begin{equation}  \label{eq.fundamentalestimate}
			K_n\leq K^\ast \sum_{j=1}^n \frac{\epsilon_j}{\alpha^j}$
			%\end{equation}
			and 
			\begin{equation} \label{eq.fundest2}
				\sup_{t\in [e^{s_0-nc},e^{s_0-nc}e^{-c}]} |z(t)|\leq K^\ast  \alpha^n \sum_{j=1}^n \frac{\epsilon_j}{\alpha^j}.
			\end{equation}
		\end{enumerate}
		%	\begin{equation} \label{eq.fundamentalestimate}
			%	K_{n+1}\leq \alpha K_n+\frac{1}{|b|}\epsilon_{n+1}, \quad n.
			%	\end{equation}
	\end{lemma}
	%	The utility of this lemma is immediate: since $K_n$ satisfies \eqref{eq.fundamentalestimate} (which is an autonomous first order linear difference  inequality with a forcing term, and therefore easy to solve) it is easy to get an estimate for $K_n$ (and hence an upper bound on $z$ ) in terms of $\epsilon_n$ (which can be inferred from the forcing term $\varphi$). Note for any choice of $N\in\mathbb{N}$ that  
	%	$K_n\leq k_n$ for $n\geq N$, where $k$ satisfies the first--order linear difference equation 
	%	\[
	%	k_{n+1} =\alpha k_n + \frac{1}{|b|}\epsilon_{n+1}, \quad n\geq N; \quad k_{N}\in (K_N,\infty),
	%	\]
	%	for which an explicit solution can be found by variation of constants:
	
	%	\subsection{Convergence and boundedness}	
	The proof of the result is inspired especially by that of Kato and McLeod~\cite[Theorem 3(i)]{KatoMc} and also of Cermak~\cite[Lemma 3.3]{cermak}. It differs from the former in considering perturbed equations. Concerning the latter paper and result, we note that it  considers more general perturbed delay equations, and yields estimates on their solutions using \textit{pointwise} conditions on forcing terms. However, the summability conditions in Lemma~\ref{lemma.fundamental} hint to the development of \textit{integrability} conditions on forcing terms in our analysis, which will ultimately yield extra sharpness to our results. In addition, our decomposition approach to solutions of the perturbed pantograph equations, shown in the next section, has the effect of turning pointwise conditions on the forcing terms into average conditions too, which will further sharpen our deterministic results. The decomposition approach will even enable us  exploit the fundamental lemma for stochastic equations, where (poor) pointwise asymptotic behaviour and regularity (i.e. non--differentiability)  of the forcing terms would confound the analysis in \cite{cermak}.    
	
	At this instant, the reader must take these claims on trust, since it is not immediately apparent how sharp or useful the estimate in part (ii) is, nor that it can be applied  to study stochastic equations. In fact, quite a lot of further analysis is needed so that  full value can be extracted from this lemma. Nevertheless, we will eventually show that this fundamental lemma (a) develops essentially unimprovable size estimates on the size of solutions of deterministic equations, which improve all existing asymptotic estimates; (b) allows us to give characterisations of the almost sure convergence and boundedness of stochastically perturbed equations, results hitherto absent from the literature. 
	
	We turn to these matters later. 
	However, we can immediately prove two direct corollaries of Lemma~\ref{lemma.fundamental}, under the assumption \eqref{eq.unforcedasstable}, giving some immediate reassurance that the Lemma~\ref{lemma.fundamental} will be of some utility. 
	
	\begin{theorem} \label{thm.1}
		Suppose $q\in (0,1)$ and \eqref{eq.unforcedasstable} holds. Let $\varphi\in C([0,\infty);\mathbb{R})$ and $z$ be the unique continuous solution to \eqref{eq.z}. 
		\begin{enumerate}
			\item[(a)] If $\varphi(t)\to 0$ as $t\to\infty$, then $z(t)\to 0$ as $t\to\infty$. 
			\item[(b)] If $\varphi$ is bounded, then $z$ is bounded. 
		\end{enumerate}		
	\end{theorem}
	\begin{proof}
		If $\varphi(t)\to 0$ as $t\to\infty$, we may take equality in \eqref{eq.epsn}, and note the hypothesis forces $\epsilon_n\to 0$ as $n\to\infty$. By \eqref{eq.unforcedasstable}, $\alpha\in (0,1)$. Either $\lambda_n=\epsilon_n/\alpha^n$ is summable or not. 
		
		Suppose $\lambda_n$ is summable. Then, by Lemma~\ref{lemma.fundamental} (i),  
		$\sup_{t\in [e^{s_0-nc},e^{s_0-nc}e^{-c}]} |z(t)|\leq K^\ast [e^{(s_0-nc)}]^\kappa$,
		so $z(t)\to 0$ as $t\to\infty$, because $\kappa<0$ and $c<0$. 
		
		Suppose $\lambda_n$ is not summable. Since $(1/\alpha^n)$ is divergent, by Toeplitz lemma (see \cite[p.390]{Shiryaev}), 	$\sum_{j=1}^n \epsilon_j \alpha^{-j} = o\left(\sum_{j=1}^n \alpha^{-j}\right)$ as $n\to\infty$, since $\epsilon_n\to 0$ as $n\to\infty$. On the other hand $\sum_{j=1}^n \alpha^{-j} \sim \alpha^{-n} (1-\alpha^{-1})$ as $n\to\infty$, so $\alpha^n\sum_{j=1}^n \epsilon_j \alpha^{-j}\to 0$ as $n\to\infty$.
		%	so $K_n\alpha^n\to 0$ as $n\to\infty$. Thus 
		%	\[
		%	e^{-\kappa(s_0-nc)}\sup_{s\in [s_0-nc,s_0-(n+1)c]} |z(e^s)|
		%	\leq \sup_{s\in [s_0-nc,s_0-(n+1)c]]} e^{-\kappa s} |z(e^s)|= K_n,
		%	\]
		%	so 
		%	\[
		%	\sup_{t\in [e^{s_0-nc},e^{s_0-nc}e^{-c}]} |z(t)|\leq K_n \alpha^n 
		%	\left(\frac{e^{-\kappa c}}{\alpha}\right)^n e^{\kappa s_0} 
		%	= K_n \alpha^n  e^{\kappa s_0}.
		%	\]
		Since the righthand side of \eqref{eq.fundest2} tends to $0$ as $n\to\infty$, $z(t)\to 0$ as $t\to\infty$.  
		
		%Notice also that if $\lambda_n$ is not summable, we have the estimate $K_n\leq K^\ast \sum_{j=1}^n \lambda_j$. Taking this into account, we arrive at 
		%	\[
		%	\sup_{t\in [e^{s_0-nc},e^{s_0-nc}e^{-c}]} |z(t)|\leq K^{\ast\ast} \alpha^n \sum_{j=1}^n \frac{\epsilon_j}{\alpha_j}.
		%	\]	
		In the case when $\varphi$ is bounded on $[0,\infty)]$ by $\overline{\varphi}$, say, we may fix $\epsilon_n=\overline{\varphi}$. By part (i) of the lemma, which holds whether $(\epsilon_n/\alpha^n)$ is summable or not, we have that 
		$		K_n\leq K^\ast \overline{\varphi} \sum_{j=1}^n \alpha^{-j} \leq \bar{K} \alpha^{-n}$ 
		since $\alpha<1$, for some $\bar{K}>0$. Thus $K_n \alpha^n\leq \bar{K}$. Revisiting the above estimates gives  
		\[
		\sup_{t\in [e^{s_0-nc},e^{s_0-nc}e^{-c}]} |z(t)|\leq K_n \alpha^n  e^{\kappa s_0}
		\leq \bar{K}e^{\kappa s_0}=:\bar{z},
		\]
		and so $z$ is bounded, as required. 
	\end{proof}
	As we see in Section 6, this qualitative result generalises  to equations with unbounded delay, using methods from \cite{AppBuck10} (see also  Theorem~\ref{thm.monotonepert} below). But our primary interest in the fundamental Lemma derives from quantitative estimates that it generates, stated below. For these estimates, we recall that a (measurable) function $\gamma:[0,\infty)\to (0,\infty)$ is regularly varying at infinity with index $\alpha$ if $\lim_{t\to\infty} \gamma(\lambda t)/\gamma(t)=\lambda^\alpha$ for all $\lambda>0$. The notation $\gamma\in \text{RV}_\infty(\alpha)$ is used. We will exploit properties of regularly varying functions throughout this work; see Bingham, Goldie and Teugels \cite{BGT}, especially Chapter 1, for relevant results.    
	%	\begin{theorem} \label{thm.growthbounds}
		%	Suppose that $q\in (0,1)$ and \eqref{eq.unforcedasstable} holds. Let $\varphi\in C([0,\infty);\mathbb{R})$ and $z$ be the unique continuous solution to \eqref{eq.z},  Let $\varphi(t)=\text{O}(\psi(t)t^{\kappa+1})$.  
		%	\begin{enumerate}
			%		\item[(A)] If $\psi\in \text{RV}_{\infty}(-1)$, $\psi\in L^1(0,\infty)$, then 
			%		$z(t)=\text{O}(t^\kappa)$ as $t\to\infty$.
			%		\item[(B)] If $\psi\in \text{RV}_{\infty}(-1)$, $\psi\not\in L^1(0,\infty)$, then 
			%		\[
			%		z(t)=\text{O}\left(t^\kappa\int_0^t \psi(s)\,ds\right), \quad t\to\infty.
			%		\] 
			%		\item[(C)] If  $\psi\in \text{RV}_{\infty}(\beta)$ for $\beta>-1$, then 
			%		the estimate in (B) holds. Equivalently, if $\varphi=\text{O}(\gamma)$ where $\gamma\in \text{RV}_\infty(\eta)$ and $\eta>\kappa$, then $z(t)=\text{O}(\gamma(t))$ as $t\to\infty$.
			%		\item[(D)] If  $\psi\in \text{RV}_{\infty}(\beta)$ for $\beta<-1$, then 
			%		the estimate in (A) holds. Equivalently, if $\varphi=\text{O}(\gamma)$ where $\gamma\in \text{RV}_\infty(\eta)$ and $\eta<\kappa$, then $z(t)=\text{O}(t^\kappa)$ as $t\to\infty$.  
			%	\end{enumerate}	
		%	\end{theorem}
	%The proof is deferred to the end. 
	%
	%The reader will notice that the conclusions in the different parts of the theorem can be combined. To start, parts (A) and (B) can be combined: let $\varphi=\text{O}(\gamma)$ where $\gamma \in \text{RV}_\infty(\kappa)$. Then $\psi(t):=\gamma(t)/t^{\kappa+1}$ for $t\geq 1$ is in $\text{RV}_\infty(-1)$, and $\varphi(t)=\text{O}(t^{\kappa+1}\psi(t))$. Then, in both part (A) and (B), we have 
	%\[
	%z(t)=\text{O}\left(t^\kappa\int_1^t \psi(s)\,ds\right) = \text{O}\left(t^\kappa\int_1^t \frac{\gamma(s)}{s^{\kappa+1}}\psi(s)\,ds\right), \quad t\to\infty.
	%\] 
	%Likewise, for part (C), when $\eta>\kappa$; since $t\mapsto \gamma(t)/t^{\kappa+1} \in \text{RV}_\infty(\eta-\kappa-1)$, and $\eta-\kappa-1>0$, we have that 
	%\[
	%t^\kappa\int_1^t \frac{\gamma(s)}{s^{\kappa+1}}\,ds \sim t^\kappa\cdot \frac{1}{\eta-\kappa}t\frac{\gamma(t)}{t^{\kappa+1}} \sim \frac{1}{\eta-\kappa}\gamma(t),
	%\]
	%so by part (C)
	%\[
	%z(t)=\text{O}(\gamma(t)) = \text{O}\left(t^\kappa\int_1^t \frac{\gamma(s)}{s^{\kappa+1}}\,ds\right), \quad t\to\infty,
	%\]
	%agreeing once again with parts (A) and (B). Finally, for part (D), where $\eta<\kappa$, we have $t\mapsto \gamma(t)/t^{\kappa+1} \in \text{RV}_\infty(\eta-\kappa-1)$. Since $\eta-\kappa-1<-1$, $t\mapsto \gamma(t)/t^{\kappa+1} \in L^1(1,\infty)$, and so by part (D)
	%\[
	%z(t)=\text{O}(t^\kappa) = \text{O}\left(t^\kappa\int_1^t \frac{\gamma(s)}{s^{\kappa+1}}\,ds\right), \quad t\to\infty.
	%\] 
	%Thus we may restate Theorem~\ref{thm.growthbounds} in a much neater manner. 
	\begin{theorem} \label{thm.growthbounds2}
		Suppose that $q\in (0,1)$ and \eqref{eq.unforcedasstable} holds. Let $\varphi\in C([0,\infty);\mathbb{R})$ and $z$ be the unique continuous solution to \eqref{eq.z},  Let $\varphi(t)=\text{O}(\gamma(t))$ as $t\to\infty$, where $\gamma$ is regularly varying at infinity. Then   
		\[
		z(t)=\text{O}\left(t^\kappa\int_1^t \frac{\gamma(s)}{s^{\kappa+1}}\,ds\right), \quad t\to\infty.
		\] 
	\end{theorem} 
	The proof is deferred to the end. Since the theorem gives only bounds on solutions, it is reasonable to ask whether (i) the bounds are sharp; (ii) the transition from the rate of decay of the  solutions of \eqref{eq.unforced} to slower decay rates, do indeed occur when $t\mapsto\varphi(t)/t^{\kappa+1}$ changes from being an integrable to a non--integrable function. The following example shows that Theorem \ref{thm.growthbounds2} is sharp in both regards. Furthermore, in Section 4, we give another example which shows that the ``big O'' estimate cannot be improved in general to results of the form $z(t)\sim c t^\kappa$ as $t\to\infty$, at least in the case of small perturbations (i.e. when $t\mapsto \gamma(t)/t^{\kappa+1}$ is integrable).     
	\begin{theorem} \label{thm.growthbdssharp}
		Let $a>0$, $b<0$, $\varphi\in C(\mathbb{R}^+)$ and $z$ the continuous solution to \eqref{eq.z}.
		\begin{enumerate}
			\item[(a)] For every $C>0$ and every  $\gamma\in \text{RV}_\infty(\kappa)$ such that $t\mapsto \gamma(t)/t^{\kappa+1}$ is not in $L^1(0,\infty)$, there exists $\varphi(t)\sim C\gamma(t)$ such that for some $D>0$,
			%t^{\kappa+1}\psi(t)$, $t\to\infty$ with $\psi\in \text{RV}_\infty(-1)$ and $\psi\not\in L^1(0,\infty)$, such that 
			\[
			z(t)\sim Dt^\kappa\int_1^t \frac{\gamma(s)}{s^{\kappa+1}}\,ds, \quad t\to\infty.
			\] 
			\item[(b)] For every $C>0$ and every  $\gamma\in \text{RV}_\infty(\kappa)$ such that $t\mapsto \gamma(t)/t^{\kappa+1}$ is in $L^1(0,\infty)$, there exists $\varphi(t)\sim C\gamma(t)$ such that $z(t)\sim Dt^\kappa$ as $t\to\infty$ 
			for some $D>0$.
		\end{enumerate}
	\end{theorem}  
	\begin{proof}
		For part (a), write $\psi(t):=\gamma(t)/t^{\kappa+1}$. Then $\psi\in \text{RV}_\infty(-1)$ and $\psi$ is not integrable. We may take $\psi$ to be a smooth regularly varying function. Let $\Psi(t):=\int_{1/q}^t \psi(s)\,ds$.
		% where we may take $\psi$ to be a smooth regularly varying function with index $-1$. 
		Then $\Psi$ is regularly varying with index zero, and $\Psi(t)\to \infty$ as $t\to\infty$. Moreover, $\Psi$ is smoothly regularly varying, with $t\Psi'(t)/\Psi(t)\to 0$ as $t\to\infty$. Set $z(t):=Dt^\kappa \Psi(t)\,t\geq q$ and $\varphi(t):=z'(t)-bz(t)-az(qt), \,t\geq 1$. Then $z$ satisfies \eqref{eq.forced} on $[1,\infty)$. Moreover  
		$z'(t)=Dt^{\kappa-1}\Psi(t)(\kappa+t\Psi'(t)/\Psi(t))\sim \kappa Dt^{\kappa-1}\Psi(t)$ as $t\to\infty$. 	For $t\geq 1$, using the fact that $b+aq^\kappa=0$, we have that
		$bz(t)+az(qt) = bDt^\kappa(\Psi(t)-\Psi(qt))$. Compare the asymptotic behaviour of $z'$ and the last term, to get the dominant asymptotic behaviour in $\varphi$.  $\Psi'=\psi$ is regularly varying with index $-1$, the uniform convergence theorem for regularly varying functions gives 
		\[
		\frac{\Psi(t)-\Psi(qt)}{t\Psi'(t)}=\int_q^1 \frac{\Psi'(\lambda t)}{\Psi'(t)}\,d\lambda
		\to \int_q^1 \lambda^{-1}\,d\lambda = \log(1/q), \quad t\to\infty.
		\]
		Hence 
		$bz(t)+az(qt) \sim b \log(1/q) D t^{\kappa+1} \psi(t)$ as $t\to\infty$, 
		so this expression is in $RV_{\infty}(\kappa)$, while $z'$ is in $RV_{\infty}(\kappa-1)$. Therefore $z'(t)=o(bz(t)+az(qt))$, and so  we have
		$\varphi(t)\sim -(bz(t)+az(qt)) \sim C t^{\kappa+1} \psi(t)$, where $C=-b \log(1/q) D>0$.
		
		For part (b), let $\psi$ be a smooth regularly varying function with index $-1$. Then $\Psi$ is regularly varying with index zero, but now since $\psi$ is integrable $\Psi(t)\to L\in (0,\infty)$ as $t\to\infty$. Moreover, $\Psi$ is also smoothly regularly varying, and so $t\Psi'(t)/\Psi(t)\to 0$ as $t\to\infty$. Next let $D>0$ be arbitrary and set
		\[
		z(t)=\frac{D}{L+1}\left(t^\kappa + t^\kappa\int_0^t \psi(s)\,ds\right), \quad t\geq q.
		\]
		Set $\varphi(t)=z'(t)-(bz(t)+az(qt))$ for $t\geq 1$. Then $z$ is the unique continuous solution of \eqref{eq.forced} on $[1,\infty)$ and $z(t)\sim Dt^\kappa$ as $t\to\infty$. We determine the asymptotic behaviour of $f$. Write $\tilde{D}=D/(L+1)$. To start, note $z'(t)=\tilde{D}\kappa t^{\kappa-1}\left(\kappa+ \kappa \Psi(t) +  t\psi(t)\right)$. 
		Since $t\psi(t)=t\Psi'(t)/\Psi(t)\cdot \Psi(t)\to 0\cdot L$=0 as $t\to\infty$, $z'(t)$ is
		order $t^{\kappa-1}$ as $t\to\infty$. Next
		\[
		bz(t)+az(qt) = b\tilde{D} t^\kappa \int_{qt}^t \psi(s)\,ds,
		\]
		so, proceeding as before we get 
		%and, as before, by the uniform convergence theorem, we have 
		%\[
		%\int_{qt}^t \psi(s)\,ds \sim t\psi(t)\log(1/q), \quad t\to\infty.
		%\]
		$bz(t)+az(qt) \sim b \tilde{D}t^{\kappa+1} \psi(t)\log(1/q)$ as $t\to\infty$, which is regularly varying with index $\kappa$, and dominates $z'(t)$. Letting  $C:=-bD(L+1)^{-1}\log(1/q)>0$, gives  $\varphi(t)\sim C t^{\kappa+1} \psi(t)=C\gamma(t)$ as $t\to\infty$, as needed.  
	\end{proof}
	Perturbations with decay faster than a power law  are covered by Theorem~\ref{thm.growthbounds2}: if $\varphi(t)=o(t^{-M})$ for all $M>0$, then $z(t)=O(t^\kappa)$, by taking $\gamma(t)=t^{\kappa-2}$. 
	% ] ince all that is needed is that $\varphi(t)=\text{O}(\gamma(t))$ where $\gamma\in \text{RV}(\kappa)$ and $\gamma/t^{\kappa+1}$ is integrable. Thus, if $\varphi(t)=\text{o}(t^{-M})$ for all $M>0$, we have $z(t)=\text{O}(r^\kappa)$. 
	
	For larger perturbations that may not have a power--law type form, the following bound can be established directly, without recourse to the fundamental lemma. 
	
	\begin{theorem} \label{thm.monotonepert}
		Suppose that $q\in (0,1)$ and \eqref{eq.unforcedasstable} holds. Let $\varphi\in C([0,\infty);\mathbb{R})$ and $z$ be the unique continuous solution to \eqref{eq.z}. Let $\gamma$ be  $C^1$ and non--decreasing.
		\begin{enumerate}
			\item[(A)]
			If $\varphi(t)=\text{O}(\gamma(t))$ as $t\to\infty$, then $z(t)=\text{O}(\gamma(t))$, as $t\to\infty$. 
			\item[(B)] If $\varphi(t)=\text{o}(\gamma(t))$ and $\gamma(t)\to \infty$ as $t\to\infty$, then $z(t)=\text{o}(\gamma(t))$, as $t\to\infty$. 
		\end{enumerate}
	\end{theorem}
	\begin{proof} The proof is inspired by those in \cite{AppBuck10,AppBuck2016}. 
		Denote by 
		$D_+ u(t) = \limsup_{h\to 0^+} (u(t+h)-u(t))/h$ 
		the Dini derivative of any continuous function $u$. Since $b<0$, by considering the cases where $z(t)$ is positive, negative or zero, we arrive at the consolidated estimate 
		$D_+|z(t)|\leq b|z(t)|+ |a||z(qt)|+|\varphi(t)|$ for $t\geq 0$.
		We note that if $z_+$ is any positive $C^1(0,\infty)$ function satisfying $z_+(0)>|z(0)|$ and $z_+'(t)> bz+(t)+ |a|z_+(qt)+|\varphi(t)|$ for all $t\geq 0$, then $z_+(t)>|z(t)|$ for all $t\geq 0$, by standard comparison arguments. By hypothesis, there exists $C>0$ such that $|\varphi(t)|\leq C\gamma(t)$ for $t\geq 0$. Take $K>\max\left(C/(|b|-|a|), |z(0)|/\gamma(0)\right)$,
		and set $z_+(t)=K\gamma(t)$ for $t\geq 0$. Then $z_+(0)=K\gamma(0)> |z(0)|$ and for $t\geq 0$ we have 
		\begin{align*}
			bz_+(t)+|a|z_+(qt)+|\varphi(t)|&\leq bK\gamma(t)+|a|K\gamma(qt)+C\gamma(t)\\
			&\leq (bK+|a|K+C)\gamma(t)<0<K\gamma'(t)=z_+'(t),
		\end{align*}
		using the monotonicity of $\gamma$. Thus $|z(t)|< z_+(t)=K\gamma(t)$ for all $t\geq 0$, as needed. 
		
		For part (B), by hypothesis, we have $|\varphi(t)|<\epsilon \gamma(t), \, t\geq T_1(\epsilon)$. For $t\geq qT_1(\epsilon)$, define $z_\epsilon(t)=\epsilon \gamma(t)/(|b|-|a|)+C(\epsilon)$ 
		where $C(\epsilon):=1+\max_{t\in [qT_1,T_1]}|z(s)|>0$. Then $z_\epsilon(t)>|z(t)|$ for $t\in [qT_1,T_1]$. For $t\geq T_1$ we have
		\begin{align*}
			bz_\epsilon(t)+|a|z_\epsilon(qt)+|\varphi(t)|
			&=(b+|a|)C	+ \epsilon\left(\frac{b}{|b|-|a|}\gamma(t)+\frac{|a|}{|b|-|a|}\gamma(qt)\right)+|\varphi(t)|\\
			&\leq (b+|a|)C	+ \epsilon\frac{b+|a|}{|b|-|a|}\gamma(t)+\epsilon \gamma(t)\\
			&= (b+|a|)C<0<\frac{1}{|b|-|a|}\epsilon \gamma'(t)=z_\epsilon'(t).
		\end{align*}
		Thus $|z(t)|<z_\epsilon(t) = \epsilon \gamma(t)/(|b|-|a|)+C(\epsilon)$ for $t\geq T_1(\epsilon)$. Dividing by $\gamma(t)$ on both sides, letting $t\to\infty$, and then $\epsilon\to 0$, we get the desired result. 
	\end{proof}
	It is easy to show that these bounds are sharp in at least some cases by considering an example. Suppose that $\gamma'(t)/\gamma(t)\to 0$ as $t\to\infty$ and $t\gamma'(t)/\gamma(t)\to \infty$ as $t\to\infty$. Then $\gamma$ is increasing subexponentially (by the first limit, see \eqref{eq.psisub} and the discussion thereafter), but faster than any power, by the second. Thus $\gamma(qt)/\gamma(t)\to 0$ as $t\to\infty$. 
	Let $z(t)=D\gamma(t)$ for $t\geq 0$, and define 
	$\varphi(t)=D[\gamma'(t)-b\gamma(t)-a\gamma(qt)]$ for $t\geq 0$. Then $z$ obeys \eqref{eq.z}, $\varphi(t)\sim -bD\gamma(t)$ as $t\to\infty$, and $z(t)\sim D\gamma(t)$ as $t\to\infty$, confirming the sharpness in part (a). 
	
	The ``little--o'' result in part (B) of Theorem~\ref{thm.monotonepert} is less traditional, and the reader may query its  value. However, we will later show that such results can be used in conjunction with the conventional ``big--O'' results to characterise when solutions obey $\limsup_{t\to\infty}|z(t)|/\gamma(t)\in (0,\infty)$, thereby establishing the exact asymptotic order of the running maxima of solutions.  
	
	We note that Theorem~\ref{thm.monotonepert} \textit{does not} give results where an exact rate of growth of the solution is given e.g., $z(t)\sim c\gamma(t)$ as $t\to\infty$, where $\gamma(t)\to \infty$ as $t\to\infty$. To give conditions on the forcing terms under which this would occur is also of interest, and will form part of a future investigation.    
	%\vspace{-0.25in}
	
	%	Finally, we compare the results of these Theorems with perturbation theorems of Cermak; he supposes that $f(t)=O(t^\beta)$ and $f'(t)=O(t^{\beta-1})$. Then 
	%	\begin{align*}
		%	x(t)&=O(t^\kappa), \text{if $\beta<\kappa$}\\
		%	x(t)&=O(t^\kappa\log t), \text{if $\beta=\kappa$}\\
		%	x(t)&=O(t^\beta), \text{if $\beta>\kappa$}.
		%	\end{align*}
	%	First, we observe that our results need neither need pointwise estimates on $f$, nor on $f'$, but merely on $y$, where $y'=-y+f$. It can readily be the case that $f$ has very bad pointwise behaviour, yet $y$ has good behaviour. For example, suppose that 
	%	\[
	%	f(t)=t^{\beta}+e^{t}\sin(e^{2t})
	%	\] 
	%	Clearly, $f(t)=O(e^t)$ and $f'(t)=O(e^{3t})$, yet $y(t)=O(t^\beta)$. Thus, we are able to conclude that 
	%	\[
	%	x(t)=O(t^{\max(\kappa,\beta)}, \text{for $\beta\neq \kappa$} 
	%	\]
	%	and $x(t)=O(t^\kappa\log t)$ as $t\to\infty$ if $\beta=\kappa$, while Cermak's theorem cannot give conclusions.
	%	
	%	Furthermore, the existing results require an explicit parametric comparison between the solution of the unperturbed equation and perturbation. For example, if $f(t)=O(t^\beta\log t)$, we cannot apply Cermak's theorem directly; instead, we would note that $f(t)=O(t^{\beta+\epsilon})$ for every $\epsilon>0$. Thus, if $\beta<\kappa$, we get $x(t)=O(t^\kappa)$ and if $\beta>\kappa$ we get $x(t)=O(t^\beta)$, but no prediction can be made in the case $\beta=\kappa$. Using our results however, we get 
	%	\[
	%	x(t)=O\left(t^\kappa\int_1^t s^{\beta-\kappa-1}\log s\,ds\right), \quad t\to\infty.
	%	\]   
	%	The estimates concur with Cermak in the cases when $\beta\neq \kappa$; but when $\beta=\kappa$, we get 
	%	\[
	%	x(t)=O(t^\kappa \log^2(t)), \quad t\to\infty,
	%	\]
	%	which is not recovered by his theorem. 
	%	
	%	We also have a converse theorem. Supposing that $x(t)=O(t^\beta)$ and $\beta>\kappa$; then $y(t)=O(t^\beta)$, while if $x(t)=O(t^\kappa)$ .   
	
	\section{Deterministically forced equation: first results, auxiliary ODEs}
	Theorem \ref{thm.1} shows that $f(t)\to 0$ as $t\to\infty$ is a \textit{sufficient} condition in order for the solution $x$ of  \eqref{eq.forced} to obey $x(t)\to 0$ as $t\to\infty$. With a little extra work, it can be used to develop a necessary and sufficient condition. Consider  
	\begin{equation} \label{eq.y}
		y'(t)=-y(t)+f(t), \quad t\geq 0; \quad y(0)=0.
	\end{equation}  
	\eqref{eq.y} has a unique continuous solution, given by 
	$y(t)=\int_0^t e^{-(t-s)}f(s)\,ds$ for $t\geq 0$. 
	
	We now show that the asymptotic behaviour of \eqref{eq.forced} is determined solely by  $y$.
	\begin{theorem} \label{thm.detasymptoticstab}
		Suppose that $q\in (0,1)$ and \eqref{eq.unforcedasstable} holds. Let $f\in C([0,\infty);\mathbb{R})$ and $x$ be the unique continuous solution to  \eqref{eq.forced}. Then the following are equivalent
		\begin{enumerate}
			\item[(A)] $y(t)\to 0$ as $t\to\infty$;
			\item[(B)] For every $\theta>0$, $\lim_{t\to\infty} \int_t^{t+\theta} f(s)\,ds =0$;
			\item[(C)] $x(t)\to 0$ as $t\to\infty$. 
		\end{enumerate}
	\end{theorem}
	\begin{proof}
		The equivalence of (A) and (B) is known; we show (A) and (C) are equivalent.
		Suppose (A) holds. Set $z(t)=x(t)-y(t)$ for $t\geq 0$, $z(0)=x(0)=\zeta$. Define 
		\begin{equation} \label{def.varphi1}
			\varphi(t)=ay(qt)+(b+1)y(t), \quad t\geq 0. 
		\end{equation}
		Note that $\varphi$ is continuous. Since $z'=x'-y'$, we have that $z$ obeys \eqref{eq.z}. 
		If (A) holds, $\varphi(t)\to 0$ as $t\to\infty$. Since $z$ obeys \eqref{eq.z}, we may apply Theorem \ref{thm.1} part (a) to it, so $z(t)\to 0$ as $t\to\infty$. But $y(t)\to 0$ as $t\to\infty$ and $x=y+z$, so (C) follows. 
		
		To show (C) implies (A), re--write \eqref{eq.forced} according to  
		$x'(t) = -x(t) + \{x(t)+bx(t)+ax(qt)\} + f(t)$ for $t\geq 0$. Let $g(t)$ be the sum of the last two terms; then $x'=-x+g$ and thus integrating by parts and rearranging we get 
		%\[
		%x(t)=\xi e^{-t} + \int_0^t e^{-(t-s)}g(s)\,ds = \xi e^{-t} + y(t) +  \int_0^t e^{-(t-s)} \{(1+b)x(s)+ax(qs)\}\,ds.
		%\] 
		%Hence
		\begin{equation} \label{eq.yrepx}
			y(t)=x(t)-\xi e^{-t} - \int_0^t e^{-(t-s)} \{(1+b)x(s)+ax(qs)\}\,ds, \quad t\geq 0.
		\end{equation}
		Since (C) holds, the first two terms on the righthand side of \eqref{eq.yrepx} tend to zero as $t\to\infty$. By (C) the term in curly brackets in the integrand tends to zero as $s\to\infty$. The integral in \eqref{eq.yrepx} is  the convolution of an $L^1(0,\infty)$ function and a function in $BC_0(0,\infty)$, so it tends to $0$ as $t\to\infty$. Thus $y(t)\to 0$ as $t\to\infty$, proving (A).    
	\end{proof}
	We note that this appears to be the first instance in the literature of a result which \textit{characterises} conditions on perturbations which preserve the asymptotic stability of non--autonomous functional differential equations. The same remark can be made concerning boundedness; see Theorem~\ref{thm.detasymptoticbound} later. Sufficient conditions for which convergence  to the equilibrium are preserved, which already appear in the literature, include $f(t)\to 0$ as $t\to\infty$ or $f\in L^1(0,\infty)$ (we note that both imply $y(t)\to 0$ as $t\to\infty$), but Theorem~\ref{thm.detasymptoticstab} shows that neither are necessary. 
	
	This leads us to ask what functions $f$ allow $y(t)\to 0$ as $t\to\infty$ which do not satisfy the familiar stability conditions above. This matter has been explored somewhat in \cite{AL:2023(AppliedNumMath)}. Two classes of functions with bad ``pointwise'' behaviour, but good ``average'' behaviour naturally emerge, but we do not claim that these characterise all the functions with stability preserving properties. 
	
	The first class are functions which exhibit high frequency oscillation (even with unbounded amplitude, provided the frequency grows faster than the amplitude). A prototype for such a function is 
	\[
	f(t)= e^{\beta t}\sin(e^{\theta t}), \quad t\geq 0,
	\]  
	where $\theta>\beta>0$. This function, by most measures, is ill--behaved as $t\to\infty$, oscillating between exponentially growing upper and lower envelopes, with increasingly higher frequency. Obviously, $f(t)=O(e^{\beta t})$ as $t\to\infty$, yet $y(t)=O(\max(e^{-t},e^{-(\theta-\beta)t}))$ as $t\to\infty$. Therefore, as $y(t)\to 0$ as $t\to\infty$, we have that $x(t)\to 0$ as $t\to\infty$. 
	
	Another class of functions $f$ which give rise to convergence of $y$, and which do not satisfy $f(t)\to 0$ as $t\to\infty$, or $f\in L^1(\mathbb{R}^+)$, are nonnegative functions with unboundedly large ``spikes'' of very short duration. Notice that if $f$ is a non--negative function, the condition (B) in Theorem~\ref{thm.detasymptoticstab} is equivalent to 
	\[
	\int_n^{n+1} f(s)\,ds \to 0, \quad n\to\infty,
	\]  
	where $n$ tends to infinity through the integers. We construct such a function now; let let $n\in \mathbb{N}$, and suppose $0<w_n<1$ and $w_n\to 0$ as $n\to\infty$. Let $f$ be zero on $[n,n+1]$ except on the interval $[n+1/2-w_n/2,n+1/2+w_n/2]$. Let $f$ be linear on $[n+1/2-w_n/2,n+1/2]$, passing through the points $(n+1/2-w_n/2,0)$ and $(n+1/2,h_n)$ where $h_n>0$ and $h_n\to\infty$ as $n\to\infty$, and likewise let $f$ be linear on  $[n+1/2,n+1/2+w_n/2]$, passing through the points  $(n+1/2,h_n)$ and $(n+1/2+w_n/2,0)$. Then $f$ is continuous and nonnegative and obeys 
	\[
	\int_{n}^{n+1} f(s)\,ds  = \frac{1}{2}w_n h_n.
	\] 
	Clearly, $\max_{n\leq s\leq n+1} f(s) = h_n$, so if we arrange that $h_n\to \infty$ as $n\to\infty$, while $w_n h_n\to 0$ as $n\to\infty$, we have a function which satisfies condition (B), and obeys $\limsup_{t\to\infty} f(t)=+\infty$. Indeed, since $h_n$ can be chosen to grow arbitrarily quickly, the running maximum of $f$ can grow arbitrarily quickly in time. It is also easy to arrange that $f$ is not in $L^1(\mathbb{R}^+)$, since for any
	$N\in \mathbb{N}$ we have 
	\[
	\int_0^{N+1} |f(s)|\,ds = \frac{1}{2} \sum_{n=0}^N w_n h_n,
	\]
	so it is merely necessary to arrange for $w_n h_n$ to be divergent. Thus, we may have $f$ having partial maxima which grow arbitrarily fast (let us take $h_n\to \infty$ as $n\to\infty$ in any way we choose such that  $h_n>1$ and $h$ increasing); picking $w_n=\frac{1}{nh_n}$, we have that  $0<w_n<1$, $w_n\to 0$ as $n\to\infty$, condition (B) holds, but $\limsup_{t\to\infty} f(t)=+\infty$ and $f\not\in L^1(\mathbb{R}^+)$. Thus, for such $f$, we have that $y(t)\to 0$ as $t\to\infty$, and hence $x(t)\to 0$ as $t\to\infty$.   
	
	Theorem~\ref{thm.detasymptoticstab} motivates a characterisation of the asymptotic behaviour of the linear differential equation \eqref{eq.y} given in \cite{AppLaw2024ODE}, from which we will presently draw. The conditions and arguments in \cite{AppLaw2024ODE} are inspired by analysis of Strauss and Yorke \cite{SY67a,SY67b} and by Grippenberg, London and Staffans~\cite[Lemma 15.9.2]{GLS}.  
	
	We now discuss the connection between (A) and (B) in Theorem~\ref{thm.detasymptoticstab}. The fact that (B) implies (A) is due to \cite{GLS}, while \cite{SY67a,SY67b} established the necessity and sufficiency of (B), albeit with extra uniform (in $\theta$) convergence required. As pointed out in \cite{AL:2023(AppliedNumMath)}, it is easy to see (A) implies (B) by integrating across \eqref{eq.y} and re--arranging: 
	\begin{equation}\label{eq.aveRepy}
		\int_{t-\theta}^t f(s)\,ds = y(t) - y(t-\theta) +\int_{t-\theta}^t y(s)\,ds,
	\end{equation}
	since the righthand side tends to zero as $t\to\infty$ if $y(t)$ does. 
	%	We define $f=0$ for $t<0$ and set $f_\theta(t)=\int_{t-\theta}^t f(s)\,ds$ for $t\geq 0$.
	%	
	%Important identities connecting $f_\theta$ and $y$ were implicitly considered in \cite{GLS} and also in \cite{AL:2023(AppliedMathLetters)}. They are exploited to tackle pointwise estimates on the solution of the differential equation \eqref{eq.y} in Appleby and Lawless. 
	%This identity is important, and we wish to highlight another crucial set of identities connecting the average of $f$ and $y$. It appears in \cite{GLS} and also in \cite{AL:2023(AppliedMathLetters)}.
	%\begin{lemma} \label{lemma.yintermsofftheta}
	%Let $f=0$ for $t<0$. Set $f_\theta(t)=\int_{t-\theta}^t f(s)\,ds$ for $t\geq 0$. Define $\delta=f-f_1$. Then %with 
	%\begin{equation} \label{eq.I}
	%I(t)=\int_0^t \delta(s)\,ds, \quad t\geq 0,
	%\end{equation}
	%we have
	%\begin{eqnarray} \label{eq.intdelta}
	%I(t)&=& \int_0^1 f_\theta(t) \,d\theta, \quad t\geq 1,\\
	%I(t)&=& \int_0^t f(s)\,ds -\int_0^t \int_0^s f(u)\,du\,ds
	%=f_1(t)-\int_0^t f_1(s)\,ds, \quad t\in [0,1], \nonumber
	%\end{eqnarray}
	%and if $y$ obeys \eqref{eq.y}, then 
	%\begin{equation} \label{eq.yrepave}
	%y(t)=\int_0^t e^{-(t-s)}f_1(s)\,ds + I(t) - \int_0^t e^{-(t-s)} I(s)\,ds.
	%\end{equation}
	%\end{lemma} 
	%We use \eqref{eq.aveRepy}, \eqref{eq.intdelta} and \eqref{eq.yrepave} extensively. 
	
	We state results characterising the rate of growth or decay of solutions of \eqref{eq.y}. We allow either for an increasing upper bound, or one which grows or decays to zero subexponentially 
	(exponential or faster--than--exponential decay is not interesting, because solutions of \eqref{eq.unforced}  decay at a power law rate and swamp exponentially small terms).  
	%An analogous result applies if we want the asymptotic behaviour of $y$ to be subexponential (which can allow for decay of $y$ to zero). 
	%Before stating such a result, %, this can only arise from subexponentially bounded $f_\theta$. 
	Recall that a function is $\gamma\in C((0,\infty);(0,\infty))$ is subexponential if 
	\begin{equation} \label{eq.psisub}
		\lim_{t\to\infty} \frac{\gamma(t-\theta)}{\gamma(t)}=1, \quad\theta>0.
	\end{equation}
	This is part of a definition of subexponential functions (see e.g., \cite{BGT}). For a comprehensive discussion of properties of this class of functions in the context of differential equations of the type studied in this paper, see \cite{AppLaw2024ODE}. We list some relevant properties to our investigation now. 
	The convergence in \eqref{eq.psisub} is in fact uniform on $\theta$--compacts, so if $\Gamma(t)=\int_{t}^{t+1} \gamma(s)\,ds$, then $\Gamma(t)\sim \gamma(t)$ as $t\to\infty$ and moreover $\Gamma'(t)/\Gamma(t)\to 0$ as $t\to\infty$. By asymptotic integration of $\Gamma'(t)/\Gamma(t)\to 0$ as $t\to\infty$ we get $\log \Gamma(t)/t\to 0$ as $t\to\infty$, whence $\log \gamma(t)/t\to 0$ as $t\to\infty$; thus we get
	$\lim_{t\to\infty} e^{\epsilon t} \gamma(t)=+\infty$ and 
	$\lim_{t\to\infty} e^{-\epsilon t} \gamma(t)=0$ for all $\epsilon>0$. These limits justify the terminology subexponential. 
	Lastly, by using l'H\^opital's rule, the above remarks show,  for any $\alpha>0$, that  
	\begin{equation} \label{eq.subconvexp}
		\int_0^t e^{-\alpha(t-s)}\gamma(s)\,ds \sim \frac{1}{\alpha}\Gamma(t)\sim \frac{1}{\alpha}\gamma(t) \quad t\to \infty.
	\end{equation}
	%which can be proven by applying L'H\^opital's rule to first factor on the righthand side of   
	%\[
	%\frac{\int_0^t e^{\alpha s}\psi(s)\,ds}{e^{\alpha t} \psi(t)}
	%=\frac{\int_0^t e^{\alpha s}\psi(s)\,ds}{e^{\alpha t} \Psi(t)} \cdot \frac{\Psi(t)}{\psi(t)}.
	%\] 
	The next result from \cite{AppLaw2024ODE} characterises of growth or decay rates of the solution \eqref{eq.y}, if the upper bound is monotone or subexponential. The crucial object to study is the function $f_\theta(t):=\int_{(t-\theta)^+}^t f(s)\,ds$, $t\geq 0$, $\theta\in [0,1]$; we compare it to a weight function $\gamma$ via hypotheses such as \eqref{eq.fthetagammabdd},  
	\begin{gather} 
		\label{eq.fthetagamma0}
		\text{
			$f_\theta(t)=o(\gamma(t))$ as $t\to\infty$}; \\
		\label{eq.fthetaequivgamma}
		\text{There exists $\theta'>0$ such that $\limsup_{t\to\infty} |f_\theta(t)|/\gamma(t)>0$}.
	\end{gather} 
	\begin{theorem} \label{lemma.fthetaymonotone}
		Let $f\in C([0,\infty);\mathbb{R})$, $y$ be the unique continuous solution to \eqref{eq.y}. 
		Let either $\gamma$ be (a)  positive, non--decreasing and in $C(0,\infty)$ or (b) subexponential. Then 
		\begin{enumerate}
			\item[(i)] \eqref{eq.fthetagammabdd} $\Longleftrightarrow$ 
			$y(t)=\text{O}(\gamma(t))$, as $t\to\infty$;	
			\item[(ii)] \eqref{eq.fthetagamma0} $\Longleftrightarrow$
			$y(t)=\text{o}(\gamma(t))$, as $t\to\infty$;	
			\item[(iii)] \eqref{eq.fthetagammabdd} and \eqref{eq.fthetaequivgamma}  $\Longleftrightarrow$
			$\limsup_{t\to\infty} |y(t)|/\gamma(t) \in (0,\infty)$.
		\end{enumerate} 
	\end{theorem}
	%\begin{proof}
	%(B) implies $|y(t)|\leq K \gamma(t)$ for some $K>0$. Apply the triangle inequality to \eqref{eq.aveRepy}, and using the monotonicity of $\gamma$, we get for all $t\geq \theta$ 
	%\[
	%\left|
	%\int_{t-\theta}^t f(s)\,ds\right| \leq K\gamma(t) +K\gamma(t-\theta) +\int_{t-\theta}^t K\gamma(s)\,ds\leq 2K\gamma(t)+\theta\gamma(t)\leq 3K\gamma(t),
	%\]  
	%which is (A) with $C=3K$. For $t\in [0,\theta]$ 
	%\[
	%\int_{(t-\theta)^+}^t f(s)\,ds=\int_0^t f(s)\,ds = y(t)+\int_0^t y(s)\,ds,
	%\]
	%The righthand side has absolute value less that $K\gamma(t)+tK\gamma(t)$, which in turn is less than $K(1+\theta)\gamma(t)\leq 2K\gamma(t)<3K\gamma(t)$, since $\theta\leq 1$. Hence (A) holds once more with $C=3K$.   
	%
	%Assuming (A), we have $|f_\theta(t)|\leq K\gamma(t)$ for all $t\geq 0$. For $t\geq 1$, 
	%we therefore have by \eqref{eq.intdelta} the estimate 
	%\[
	%|I(t)|\leq \int_0^1 |f_\theta(t)|\,d\theta\leq K\gamma(t).
	%\]
	%For $t\in [0,1]$, we have $|f_1(t)|\leq K\gamma(t)$, and by the equation after \eqref{eq.intdelta}, we get 
	%\[
	%|I(t)|\leq |f_1(t)| + \int_0^t |f_1(s)| \,ds \leq K\gamma(t)+\int_0^t K\gamma(s)\,ds \leq 2K\gamma(t).
	%\]
	%Thus, $|I(t)|\leq 2K\gamma(t)$ for all $t\geq 0$. Now, take the triangle inequality in \eqref{eq.yrepave}, apply the above estimates for $I$ and $f_1$, and use the monotonicity of $\gamma$ to get  
	%\begin{align*} 
	%|y(t)|&\leq \int_0^t e^{-(t-s)}|f_1(s)|\,ds + |I(t)| + \int_0^t e^{-(t-s)} |I(s)|\,ds\\
	%&\leq \int_0^t e^{-(t-s)}K\gamma(s) + 2K\gamma(t) + \int_0^t e^{-(t-s)} 2K\gamma(s)\,ds\\
	%&\leq K\gamma(t)+2K\gamma(t)+2K\gamma(t)=5K\gamma(t),
	%\end{align*}
	%as needed. 
	%\end{proof}
	The hypothesis (A) in the above Theorem requires uniform control of $f_\theta$ for all $\theta\in [0,1]$. The choice of $[0,1]$ is arbitrary; control on an interval of the form $[0,\delta]$ for any choice of $\delta>0$ is equivalent. This holds for all hypotheses of this type, whether the weight function $\gamma$ is subexponential, regularly varying or monotone.  
	
	%An analogous result applies if we want the asymptotic behaviour of $y$ to be subexponential (which can allow for decay of $y$ to zero). In advance of stating such a result, %, this can only arise from subexponentially bounded $f_\theta$. 
	%recall that a function is $\psi\in C((0,\infty);(0,\infty))$ is subexponential if 
	%\begin{equation} \label{eq.psisub}
	%\lim_{t\to\infty} \frac{\psi(t-\theta)}{\psi(t)}=1, \quad\theta>0.
	%\end{equation}
	%We make some brief comments about this class of functions. First, the convergence in \eqref{eq.psisub} is in fact uniform on $\theta$--compacts. Second, if we take the function $\Psi(t)=\int_{t}^{t+1} \psi(s)\,ds$, then $\Psi(t)\sim \psi(t)$ as $t\to\infty$ and moreover $\Psi'(t)/\Psi(t)\to 0$ as $t\to\infty$, so it is easy to manufacture subexponential functions with the same asymptotic properties, but with additional control and regularity features. Third, we notice that asymptotic integration of $\Psi'(t)/\Psi(t)\to 0$ as $t\to\infty$ leads to $\log \Psi(t)/t\to 0$ as $t\to\infty$, and in turn to $\log \psi(t)/t\to 0$ as $t\to\infty$. This implies 
	%\[
	%\lim_{t\to\infty} e^{\epsilon t} \psi(t)=+\infty, \quad 
	%\lim_{t\to\infty} e^{-\epsilon t} \psi(t)=0,
	%\]
	%so $\psi$ asymptotically dominates every negative exponential function, but is dominated by every positive exponential function, justifying the terminology subexponential. 
	%Fourth, a consequence of the second remark is that for any $\alpha>0$ we have 
	%\begin{equation} \label{eq.subconvexp}
	%\int_0^t e^{-\alpha(t-s)}\psi(s)\,ds \sim \frac{1}{\alpha}\psi(t), \quad t\to \infty
	%\end{equation}
	%which can be proven by applying L'H\^opital's rule to first factor on the righthand side of   
	%\[
	%\frac{\int_0^t e^{\alpha s}\psi(s)\,ds}{e^{\alpha t} \psi(t)}
	%=\frac{\int_0^t e^{\alpha s}\psi(s)\,ds}{e^{\alpha t} \Psi(t)} \cdot \frac{\Psi(t)}{\psi(t)}.
	%\] 
	%We have all the ingredients to prove the next result. 
	%\begin{theorem} \label{lemma.fthetaysubexponential}
	%		Let $f\in C([0,\infty);\mathbb{R})$, and $y$ be the unique continuous solution to \eqref{eq.y}. 
	%	Let $\psi$ be subexponential. 
	%	\begin{itemize}
		%		\item[(i)] The following are equivalent:
		%	\begin{enumerate}
			%		\item[(A)] There is $C>0$ such that $|\int_{(t-\theta)^+}^t f(s)\,ds|\leq C\psi(t)$ for all $\theta\in [0,1]$;
			%		\item[(B)] $y(t)=\text{O}(\psi(t))$, as $t\to\infty$;	
			%	\end{enumerate}
		%	\item[(ii)] The following are equivalent:
		%		\begin{enumerate}
			%		\item[(A)] $\int_{t-\theta}^t f(s)\,ds/\psi(t)\to 0$ for all $\theta\in [0,1]$;
			%		\item[(B)] $y(t)=\text{o}(\psi(t))$, as $t\to\infty$;	
			%	\end{enumerate}
		%		\end{itemize}
	%\end{theorem}
	%\begin{proof}
	%Assume (B). Then $|y(t)|\leq K\psi(t)$ for all $t\geq 0$. Let $\theta\in [0,1]$. For $t\geq 1\geq \theta$, by applying the triangle inequality to \eqref{eq.aveRepy}, we have 
	%\[
	%\frac{1}{\psi(t)}\left|\int_{(t-\theta)^+}^t f(s)\,ds\right| \leq K+K\frac{\psi(t-\theta)}{\psi(t)}+K\int_{u=0}^\theta \frac{\psi(t-u)}{\psi(t)}\,du.
	%\] 
	%By the uniform convergence in the limit, there is a fixed $T>1$ such that for all $u\in [0,1]$, $\psi(t-u)/\psi(t)\leq 2$ for $t\geq T$. Thus, for $t\geq T$, the righthand side is bounded by $5K$. Also, on the compact $(t,\theta)\in [1,T]\times [0,1]$ the continuous function 
	%\[
	%(t,\theta)\to \frac{1}{\psi(t)}\left|\int_{(t-\theta)^+}^t f(s)\,ds\right| =: F(t;\theta)
	%\]  
	%is uniformly bounded. Hence there is a $t,\theta$--independent positive constant $K'$ such that 
	%\[
	%\left|\int_{(t-\theta)^+}^t f(s)\,ds\right| \leq K'\psi(t), \quad \theta\in [0,1], \quad t\geq 1.
	%\]
	%To get a uniform estimate on the compact $[0,1]\times [0,1]$, simply appeal to the continuity of the function $F$ once again. This proves (A). 
	%
	%To show (A) implies (B), by assumption we have $|f_\theta(t)|\leq K\psi(t)$ for all $t\geq 0$ and $\theta\in [0,1]$. Thus, by \eqref{eq.intdelta}, we have $|I(t)|\leq K\psi(t)$ for all $t\geq 1$. For $t\in [0,1]$, by the equation after \eqref{eq.intdelta} we have 
	%\[
	%|I(t)|\leq K\psi(t)+K\int_0^t \psi(s)\,ds \leq K'\psi(t),
	%\]
	%for some $K'>K$, using the positivity and continuity of $\psi$ to bound the integral. Thus, we have the overall estimate $|I(t)|\leq K'\psi(t)$ for $t\geq 0$. Now, take the triangle inequality in \eqref{eq.yrepave}, scale by $\psi$, and apply the above estimates for $I$ and $f_1$ to get 
	%\begin{align*} 
	%	\frac{|y(t)|}{\psi(t)}&\leq \frac{1}{\psi(t)}\int_0^t e^{-(t-s)}|f_1(s)|\,ds + \frac{|I(t)|}{\psi(t)} + \frac{1}{\psi(t)}\int_0^t e^{-(t-s)} |I(s)|\,ds\\
	%	&\leq (K+K')\frac{1}{\psi(t)}\int_0^t e^{-(t-s)}\psi(s)\,ds + K'. 
	%	\end{align*}  
%By \eqref{eq.subconvexp} with $\alpha=1$, the second factor in the first term on the right tends to 1 as $t\to\infty$, so by continuity is bounded (by $B$ say). Hence 
%$|y(t)|\leq ((K+K')B+K')\psi(t)$ for all $t\geq 0$, proving (B).
%\end{proof}
%
%The following result helps to determine the exact size of subexponential solutions.
%\begin{lemma} \label{lemma.fthetaysubexponentiallittleo}
%		Let $f\in C([0,\infty);\mathbb{R})$, and $y$ be the unique continuous solution to \eqref{eq.y}. 
%	Let $\psi$ be subexponential. Then the following are equivalent:
%	\begin{enumerate}
	%		\item[(A)] $\int_{t-\theta}^t f(s)\,ds/\psi(t)\to 0$ for all $\theta\in [0,1]$;
	%		\item[(B)] $y(t)=\text{o}(\psi(t))$, as $t\to\infty$;	
	%	\end{enumerate}
%\end{lemma} 
%\begin{proof}
%	To show that (B) implies (A), note for every $\epsilon>0$ that $|y(t)|\leq \epsilon \psi(t)$ for all $t\geq T_1(\epsilon)$, which we can take greater than unity, without loss of generality. Thus, for all $t\geq T_1(\epsilon)+1$, and $\theta\in [0,1]$, we have 
%	\[
%	\left|\int_{t-\theta}^t f(s)\,ds\right| \leq \epsilon \psi(t)+\epsilon \psi(t-\theta) +\epsilon \int_{u=0}^{\theta} \psi(t-u)\,du
%	\] 
%	Next, by the uniform convergence property, $\psi(t-u)\leq 2\psi(t)$ for all $t\geq T$ and all $u\in [0,1]$. Take $T_2(\epsilon)=\max(T_1+1,T)$ and thus for $t\geq T_2(\epsilon)$ we have 
%	\[
%		\left|\int_{t-\theta}^t f(s)\,ds\right| \leq \epsilon \psi(t)+2\epsilon \psi(t) +\epsilon \int_{u=0}^{1} 2\psi(t)\,du=5\epsilon \psi(t),
%	\]
%	and since $\epsilon>0$ is arbitrary, we have (A). 
%	
%	To show that (A) implies (B), we have also that $\int_t^{t+\theta} f(s)\,ds = o(\psi(t))$. Hence we may define 
%	\[
%	Q_m=\left\{\theta\in [0,1]: \left|\int_{t}^{t+\theta} f(s)\,ds \right|\leq \psi(t)\right\}.
%	\]  
%	By continuity of the functions, $Q_m$ is closed and measurable, with $Q_{m+1}\subseteq Q_m$. Also $\cup_{m=1}^\infty Q_m=[0,1]$. Therefore, we have that there exists an $m=m'$ such that $\text{Leb}(Q_m)>0$. Fix this $m$. Using the notation  
%	\[
%	Q_m-Q_m:=\{\theta-\theta':\theta,\theta'\in Q_m\}
%	\]
%	by Lemma 15.9.2 in \cite{GLS}, $Q_m-Q_m$ contains an interval $(-\delta,\delta)$ for some $\delta>0$. Let $\theta,\theta'\in [0,1]$ be in $Q_m$, and let $\theta>\theta'$ without loss. Then for $t\geq m$, we have 
%	\[
%	\int_{t}^{t+\theta-\theta'} f(s)\,ds = \int_t^{t+\theta} f(s)\,ds - \int_{t+\theta-\theta'}^{t+\theta} f(s)\,ds.
%	\] 
%	Hence
%	\[
%	\left|	\int_{t}^{t+\theta-\theta'} f(s)\,ds\right|\leq \psi(t)+\psi(t+\theta-\theta').
%	\]
%	Now, for all $\vartheta\in [0,1]$, there is a $T>0$ such that $\psi(t+\vartheta)\leq 2\psi(t)$. Hence, with $T'=\max(T,m)$, we have
%	\[
%	\left|	\int_{t}^{t+\theta-\theta'} f(s)\,ds\right|\leq 3\psi(t), \quad t\geq T'.
%	\]
%	Hence
%	\[
%	\left|	\int_{t}^{t+\vartheta} f(s)\,ds\right|\leq 3\psi(t), \quad \vartheta\in (0,\delta), \quad t\geq T.
%	\]
%	From this, and the uniform convergence property of $\psi$, we get the estimate 
%	\[
%	|f_\theta(t)|/\psi(t) \leq K, \quad t\geq T^\ast, \quad \theta\in [0,1],
%	\]
%	for some $K>0$ which is $t$-- and $\theta$--independent. 
%	
%	Now, for $t\geq 1$, we have $I(t)/\psi(t)=\int_0^1 f_\theta(t)/\psi(t)\,d\theta$. Since $f_\theta(t)/\psi(t)\to 0$ as $t\to\infty$ for each $\theta\in [0,1]$, by dominated convergence, it follows that $I(t)/\psi(t)\to 0$ as $t\to\infty$. Since $f_1(t)/\psi(t)\to 0$ by hypothesis, it follows from \eqref{eq.yrepave} that $y(t)=\text{o}(\psi(t))$ as $t\to\infty$ as required in (B), provided 
%	\begin{equation} \label{eq.littleosubexpconv}
	%	g(t)=\text{o}(\psi(t)), \quad  t\to\infty\quad \text{ implies }
	%	\int_0^t e^{-(t-s)}g(s)\,ds = \text{o}(\psi(t)), \quad t\to\infty,
	%	\end{equation}
%	where we take $g$ to be $f_1$ or $I$. %Take $\Psi(t)=\int_t^{t+1} \psi(s)\,ds$ so that $\Psi'(t)/\Psi(t)\to 0$ as $t\to\infty$. 
%	To prove \eqref{eq.littleosubexpconv}, by hypothesis for every $\epsilon>0$ there is a $T(\epsilon)>0$ such that $|g(t)|\leq \epsilon\psi(t)$ for $t\geq T(\epsilon)$. Then for $t\geq T(\epsilon)$, we have 
%	\[
%	\frac{1}{\psi(t)e^t}\left|
%	\int_0^t e^s g(s)\,ds\right| \leq \frac{1}{\psi(t)e^t}\left|\int_0^T e^s g(s)\,ds\right|+\epsilon \frac{1}{\psi(t)e^t}\int_T^t e^s\psi(s)\,ds 
%	\]
%	Since $\psi$ is subexponential, $\psi(t)e^t \to \infty$ as $t\to\infty$, and using \eqref{eq.subconvexp} with 	$\alpha=1$ shows that the second term has unit limit as $t\to\infty$. Therefore  
%	\[
%	\limsup_{t\to\infty} \frac{1}{\psi(t)e^t}\left|
%	\int_0^t e^s g(s)\,ds\right| \leq \epsilon,
%	\]
%	and since $\epsilon>0$ is arbitrary, we have established  \eqref{eq.littleosubexpconv} as required. 
%\end{proof}
%
%We deliver on the promise made before the proof of this result, namely, that taken together, 

We see that the last part Theorem~\ref{lemma.fthetaymonotone} can be used to characterise the size of solutions of \eqref{eq.y} \textit{exactly} to leading order. Roughly, $y(t)$ is $O(\gamma(t))$ but not $o(\gamma(t))$ as $t\to\infty$ if and only if $f_\theta(t)$ is $O(\gamma(t))$ for all $\theta$, but not $o(\gamma(t))$ for some $\theta$, as $t\to\infty$. 
%\begin{lemma} \label{lemma.fthetaysubexponentialexact}
%	Let $f\in C([0,\infty);\mathbb{R})$, and $y$ be the unique continuous solution to \eqref{eq.y}. 
%	Let $\psi$ be subexponential. Then the following are equivalent:
%\begin{enumerate}
%	\item[(A)] There exists $\theta'>0$ such that 
%	\[
%	\limsup_{t\to\infty} \frac{|\int_{t-\theta'}^t f(s)\,ds|}{\psi(t)}>0
%	\]
%	while $\left|\int_{(t-\theta)^+}^t f(s)\,ds \right|\leq K \psi(t)$ for all $\theta\in [0,1]$, $t\geq 0$;
%	\item[(B)] 
%	\[
%	\limsup_{t\to\infty} \frac{|y(t)|}{\psi(t)} \in (0,\infty).
%	\]
%\end{enumerate}
%\end{lemma}
%\begin{proof}
%Assume (B). Then $|y(t)|\leq L\psi(t)$ for all $t\geq 0$, so by Lemma~\ref{lemma.fthetaysubexponential}, we have the first part of (A). Suppose, by way of contradiction, that the first part of (A) is false, so that 
%\begin{equation} \label{eq.contra1}
%\limsup_{t\to\infty} \frac{|\int_{t-\theta'}^t f(s)\,ds|}{\psi(t)}=0, \text{ for all $\theta'\in [0,1]$.}
%\end{equation}
%But then, by Lemma~\ref{lemma.fthetaysubexponentiallittleo}, we must have that $y(t)=\text{o}(\psi(t))$ as $t\to\infty$, contradicting the hypothesis (B). Therefore, the first part of (A) must also be true, and we have that (B) implies (A).
%
%Assume (A). By the second part of (A), we have $y(t)=\text{O}(\psi(t))$, by Lemma~\ref{lemma.fthetaysubexponential}, so the limsup in (B) is finite. Suppose, by way of contradiction, that it is zero, so that $y(t)=\text{o}(\psi(t))$ as $t\to\infty$. But then, by Lemma~\ref{lemma.fthetaysubexponentiallittleo} we have that \eqref{eq.contra1} holds. But this contradicts the first part of (A), which is assumed true by hypothesis. Hence, the supposition that the limsup in (B) is zero is false, and therefore the limsup in (B) is finite and positive, proving (B). This secures the desired equivalence. 
%\end{proof}
%
%It is reasonable to ask if in the case that $\gamma$ is monotone whether a result of the above type can be shown. The answer is positive. 
%
%The first part of  Lemma~\ref{lemma.fthetaysubexponentiallittleo} goes through more or less verbatim. With cosmetic changes to the beginning of second part of Lemma~\ref{lemma.fthetaysubexponentiallittleo}, it can be shown that when 
%\[
%\int_{(t-\theta)^+}^t f(s)\,ds = \text{o}(\gamma(t)), \quad t\to\infty,
%\]
%then $|f_\theta(t)|\leq C\gamma(t)$ for all $t\geq T$ and all $\theta\in [0,\delta]$ for some $\delta\in (0,1)$, where $C$ is $\theta--$ and $t$--independent. We show how this bound can be extended to all $\theta\in [0,1]$, since the proof was glossed over above. To start, note we have a minimal $N^\ast\geq 1$ such that $(N^\ast+1) \delta>1$. Also there is an $N(\theta)\leq N^\ast$ such that $N\delta\leq \theta$ and $(N+1)\delta>\theta$. Write 
%\[
%\int_{t-\theta}^t f(s)\,ds = \sum_{j=1}^N \int_{t-\theta+(j-1)\delta}^{t-\theta+j\delta} f(s)\,ds + \int_{t-(\theta-N\delta)}^t f(s)\,ds.
%\]
%Taking the triangle inequality, each term in the sum is bounded by $C\gamma(t-\theta+j\delta)$, which in turn is bounded by $C\gamma(t)$. Likewise, the last term is bounded by $C\gamma(t)$, since the interval is of length less than $\delta$. Thus
%for any $\theta\in [0,1]$ and $t\geq T+2$, we have 
%\[
%\left|\int_{t-\theta}^t f(s)\,ds\right|\leq (N+1)C\gamma(t)\leq (N^\ast+1)C\gamma(t)
%\leq \left(1/\delta+1\right)C\gamma(t)=:K\gamma(t).
%\]
%Therefore, by dominated convergence, we have as before that $I(t)/\gamma(t)\to 0$ as $t\to\infty$, and $f_1(t)=\text{o}(\gamma(t))$ by hypothesis. Hence for $t\geq T(\epsilon)$ we have $|I(t)|\leq \epsilon \gamma(t)$ and $|f_1(t)|\leq \epsilon \gamma(t)$. Set $J(T)= \int_0^T e^{s}(|f_1(s)|+|I(s)|)\,ds$. Hence by \eqref{eq.yrepave} for $t\geq T$ we have 
%\begin{align*}
%|y(t)|&\leq \int_0^T e^{-(t-s)}|f_1(s)|\,ds + \int_T^t  e^{-(t-s)}|f_1(s)|\,ds + |I(t)| \\ &\qquad+\int_0^T e^{-(t-s)} |I(s)|\,ds + \int_T^t e^{-(t-s)} |I(s)|\,ds\\
%& \leq e^{-t} J(T) 
%+ 2\epsilon \int_T^t  e^{-(t-s)}\gamma(s)\,ds + \epsilon \gamma(t)\\
%&\leq  e^{-t} J(T)
%+ 2\epsilon \int_0^{t-T}  e^{-u}\,du \cdot \gamma(t) + \epsilon \gamma(t)\\
%&\leq 3\epsilon \gamma(t)+ e^{-t}J(T). 
%\end{align*}
%Hence $\limsup_{t\to\infty} |y(t)|/\gamma(t)\leq 3\epsilon$, so as $\epsilon>0$ is arbitrary, we have $y(t)=o(\gamma(t))$ as $t\to\infty$, as claimed.  
%
%We have hence sketched the important details of the following result.
%\begin{lemma} \label{lemma.fthetaymonotonelittleo}
%	Let $f\in C([0,\infty);\mathbb{R})$, and $y$ be the unique continuous solution to \eqref{eq.y}. 
%	Let $\gamma$ be positive, non--decreasing and in $C(0,\infty)$. Then the following are equivalent:
%	\begin{enumerate}
	%		\item[(A)]  $\int_{t-\theta}^t f(s)\,ds=\text{o}(\gamma(t))$ as $t\to\infty$ for all $\theta\in [0,1]$;
	%		\item[(B)] $y(t)=\text{o}(\gamma(t))$, as $t\to\infty$;	
	%	\end{enumerate}
%\end{lemma}
%From this result, an analogue to Lemma~\ref{lemma.fthetaysubexponentialexact} can be proven, with Lemma~\ref{lemma.fthetaymonotone} also in support. The proof is essentially identical to that of Lemma~\ref{lemma.fthetaysubexponentialexact} so is omitted.

%\begin{lemma} \label{lemma.fthetaymonotoneexact}
%	Let $f\in C([0,\infty);\mathbb{R})$, and $y$ be the unique continuous solution to \eqref{eq.y}. 
%	Let $\gamma$ be positive, non--decreasing and in $C(0,\infty)$. Then the following are equivalent:
%		\item[(A)] There exists $\theta'>0$ such that 
%	\[
%	\limsup_{t\to\infty} \frac{|\int_{t-\theta'}^t f(s)\,ds|}{\gamma(t)}>0
%	\]
%	while $\left|\int_{(t-\theta)^+}^t f(s)\,ds \right|\leq K \gamma(t)$ for all $\theta\in [0,1]$, $t\geq 0$;
%	\item[(B)] 
%	\[
%	\limsup_{t\to\infty} \frac{|y(t)|}{\gamma(t)} \in (0,\infty).
%	\]
%\end{lemma}
\vspace{-0.25in}

\section{Asymptotics of deterministic pantograph equation}
It is now easy to characterise boundedness of solutions of the pantograph equation. 
\begin{theorem}\label{thm.detasymptoticbound}
	Suppose that $q\in (0,1)$ and \eqref{eq.unforcedasstable} holds. Let $f\in C([0,\infty);\mathbb{R})$ and $x$ be the unique continuous solution to  \eqref{eq.forced}. Then the following are equivalent
	\begin{enumerate}
		\item[(A)] $y$ is bounded;
		\item[(B)] There exists $F\geq 0$ such that $|\int_t^{t+\theta} f(s)\,ds|\leq F$ for all $\theta\in [0,1]$, $t\geq 0$;
		\item[(C)] $x$ is bounded. 
	\end{enumerate}
\end{theorem}
To prove this, apply part (A) of Theorem~\ref{thm.monotonepert} to $z=x-y$. Since $y$ is bounded, $\varphi$ given by \eqref{def.varphi1} is bounded. Thus $z$ is bounded and hence  $x=z+y$ is bounded, so (A) implies (C). (C) implies (A) by virtue of \eqref{eq.yrepx}. Applying part (i) of Theorem~\ref{lemma.fthetaymonotone} with $\gamma(t)=1$ shows that (A) and (B) are equivalent, completing the proof. 

We note of course that taking the results of Theorems~\ref{thm.detasymptoticstab} and~\ref{thm.detasymptoticbound} together characterise when $x$ is bounded but non--convergent. This will happen if and only if condition (B) of Theorem~\ref{thm.detasymptoticbound} holds, but condition (B) of Theorem~\ref{thm.detasymptoticstab} does not hold. 

A simple situation where this happens is if $f(t)\to f^\ast\neq 0$ as $t\to\infty$, or if $f(t)=\sin(t)$ for $t\geq 0$. A more complicated situation, where $f$ is unbounded, is furnished by a simple modification of the second example after Theorem~\ref{thm.detasymptoticstab}. To give boundedness (but non--convergence) of solutions of the perturbed pantograph equation, we now choose $w_n=1/h_n$, while still taking $h_n>1$ and $h_n\uparrow \infty$ as $n\to\infty$ (once again, therefore, the rate of growth of the running maximum of $|f|$ can be arbitrarily fast in time). In this case, for all $\theta\in [0,1]$ and all $t\geq 0$ we have  
\[
\int_{t}^{t+\theta} f(s)\,ds \leq \int_{t}^{t+1} f(s)\,ds \leq \int_{n}^{n+2} f(s)\,ds = 1,
\]   
where $n$ is the largest integer less than or equal to $t$. Thus condition (B) of Theorem~\ref{thm.detasymptoticbound} holds. On the other hand
\[
\int_{n}^{n+1} f(s)\,ds = \frac{1}{2},
\]
for every integer $n$, so condition (B) of Theorem~\ref{thm.detasymptoticstab} does not hold.

An interesting class of bounded solutions are those which tend to non--trivial limits. The situation when this happens can also be characterised. Once again, we believe this is the first such characterisation of convergent solutions for nonautonomous functional differential equations.  We leave to the interested reader the proof of the following theorem, whose proof is similar to those of Theorems~\ref{thm.detasymptoticstab} and \ref{thm.detasymptoticbound}. Consult the proof of~\cite[Lemma 15.9.2]{GLS} to prove the equivalence of (A) and (B); the equivalence of (A) and (C) can be proven with sole recourse to arguments given here.

\begin{theorem}\label{thm.detasymptoticlimit}
	Suppose that $q\in (0,1)$ and \eqref{eq.unforcedasstable} holds. Let $f\in C([0,\infty);\mathbb{R})$ and $x$ be the unique continuous solution to  \eqref{eq.forced}. Then the following are equivalent
	\begin{enumerate}
		\item[(A)] $y(t)$ tends to a non--trivial finite limit as $t\to\infty$;
		\item[(B)] There exists $F\neq 0$ such that for every $\theta\in [0,1]$ that $\frac{1}{\theta}\int_t^{t+\theta} f(s)\,ds\to F$ as $t\to\infty$ for all $\theta\in [0,1]$;
		\item[(C)] $x(t)$ tends to a non--trivial finite limit as $t\to\infty$.  
	\end{enumerate}
\end{theorem}
It can be easily shown that the limits at infinity of $x$ and $y$ are related to $F$ in condition (B) according to 
\[
\lim_{t\to\infty} x(t)=\frac{F}{1-(b+a)}, \quad  \lim_{t\to\infty} y(t)=F.
\] 
%\begin{proof}
%The equivalence of (A) and (B) follows by taking $\gamma(t)=1$ in Lemma~\ref{lemma.fthetaymonotone}.  
%, 
%Suppose (A) is true. Define $z(t)=x(t)-y(t)$ for $t\geq 0$. With $\varphi$ given by \eqref{def.varphi1}, we see that $z$ satisfies \eqref{eq.z}. Moreover, by (A), $\varphi$ is uniformly bounded and continuous. Hence, by part (b) of Theorem \ref{thm.1}, $z$ is bounded. Hence $x$ is bounded, proving (A).  
%
%Assume (C) and proceed as above to deduce the representation \eqref{eq.yrepx}. This time, (C) ensures the boundedness of the terms in curly brackets in the integral, and hence the integral itself, while the first two terms on the right of \eqref{eq.yrepx} are bounded too. Thus $y$ is bounded, proving (A).    
%%To see that (A) implies (B), use the uniform bound on $y$, and take the triangle inequality in \eqref{eq.aveRepy}. To show that (B) implies (A), use the fact that $I$ defined in \eqref{eq.intdelta} is uniformly bounded, by (B), and that $f_1$ is uniformly bounded by hypothesis. This means that all the terms on the righthand side of \eqref{eq.yrepave} are uniformly bounded, proving (A).  
%\end{proof}

We now use Theorem~\ref{thm.growthbounds2} to give rather sharp conditions on the asymptotic behaviour of solutions, without imposing pointwise conditions on the forcing function $f$. We start by giving a theorem for large perturbations. 
\begin{theorem}\label{thm.detlargepert}
	Suppose that $q\in (0,1)$ and \eqref{eq.unforcedasstable} holds. Let $f\in C([0,\infty);\mathbb{R})$ and $x$ be the unique continuous solution to \eqref{eq.forced}. Suppose $\gamma$ obeys either (i) $\gamma\in \text{RV}_\infty(\eta)$ for some $\eta>\kappa$ or (ii) $\gamma$ is positive, increasing and in $C^1$, with $\gamma(t)\to\infty$ as $t\to\infty$. 
	\begin{enumerate}
		\item[(I)]
		%Let $\eta>\kappa$ and $\gamma\in \text{RV}_\infty(\eta)$. 
		\eqref{eq.fthetagammabdd} $\Longleftrightarrow$ 
		$y(t)=\text{O}(\gamma(t))$, as $t\to\infty$
		$\Longleftrightarrow$ 
		$x(t)=\text{O}(\gamma(t))$, as $t\to\infty$. 
		\item[(II)] 
		\eqref{eq.fthetagamma0}  $\Longleftrightarrow$  
		$y(t)=\text{o}(\gamma(t))$, as $t\to\infty$ 
		$\Longleftrightarrow$  
		$x(t)=\text{o}(\gamma(t))$, as $t\to\infty$. 
		\item[(III)] \eqref{eq.fthetagammabdd}, \eqref{eq.fthetaequivgamma} 
		$\Leftrightarrow$ 
		$\limsup_{t\to\infty} |y(t)|/\gamma(t)\in (0,\infty)$
		$\Leftrightarrow$ 
		$\limsup_{t\to\infty} |x(t)|/\gamma(t)\in (0,\infty)$.
	\end{enumerate}
	%	\item[(III)] \hspace{0.1in}
	%	Let $\gamma$ be positive, increasing and in $C^1(0,\infty)$. Then the following are equivalent:
	%	\begin{enumerate}
		%		\item[(A)] There is $K>0$ such that $\left|\int_{(t-\theta)^+}^t f(s)\,ds\right|\leq K\gamma(t)$ for all $\theta\in [0,1]$, $t\geq 0$;
		%		\item[(B)] $y(t)=\text{O}(\gamma(t))$, as $t\to\infty$;
		%		\item[(C)] $x(t)=\text{O}(\gamma(t))$, as $t\to\infty$. 
		%	\end{enumerate}
	%	\item[(IV)] \hspace{0.1in} 
	%	Let $\gamma$ be positive, increasing and in $C^1(0,\infty)$ with $\gamma(t)\to\infty$ as $t\to\infty$. Then the following are equivalent:
	%	\begin{enumerate}
		%		\item[(A)]  $\int_{t-\theta}^t f(s)\,ds=\gamma(t)$ for all $\theta\in [0,1]$;
		%		\item[(B)] $y(t)=\text{o}(\gamma(t))$, as $t\to\infty$;
		%		\item[(C)] $x(t)=\text{o}(\gamma(t))$, as $t\to\infty$. 
		%	\end{enumerate}
	%	\end{itemize}
	\end{theorem}
	\begin{proof}
Note part (III) follows from parts (I) and (II). 
For part (I), note that if $\gamma$ is regularly varying, it is also subexponential, so (A) and (B) are equivalent by Theorem~\ref{lemma.fthetaymonotone}. Thus is suffices to prove the equivalence of (B) and (C).

Next, let $z=x-y$ and $\varphi(t)=(1+b)y(t)+ay(qt)$. Then $z$ obeys \eqref{eq.z}. If (B) holds, $|y(t)|\leq C_0\gamma(t)$. Then, as $\gamma$ is regularly varying, $|\varphi(t)|\leq C\gamma(t)$ for some $C>0$. Since $\eta>\kappa$, part (C) of Theorem \ref{thm.growthbounds2} applies, so $z(t)=O(\gamma(t))$ as $t\to\infty$. But since $y(t)=O(\gamma(t))$ and $x=y+z$,  $x(t)=O(\gamma(t))$ as $t\to\infty$, which is (C). 

To show (C) implies (B), note that (C) yields $|x(t)|\leq K\gamma(t)$ for $t\geq 0$. This also gives $|x(qt)|\leq K\gamma(qt)\leq L\gamma(t)$, using the regular variation of $\gamma$, for constants $K$ and $L$. Applying the triangle inequality to \eqref{eq.yrepx}, and using these estimates, we get 
\begin{equation} \label{eq.yrvbound}
	|y(t)|\leq |\xi|e^{-t}+K\gamma(t)+\int_0^t e^{-(t-s)}\left(|1+b|K\gamma(s)+|a|L\gamma(s)\right)\,ds.
\end{equation}
Now, note that there is a $\Gamma$ such that $\gamma\sim \Gamma$ and $t\Gamma'(t)\sim \eta\Gamma(t)$ as $t\to\infty$. Thus, $\Gamma'(t)=o(\Gamma(t))$ as $t\to\infty$ and $e^t\Gamma(t)\to\infty$ as $t\to\infty$. Thus the integral in \eqref{eq.yrvbound} is majorised by some positive constant times $e^{-t}\int_0^t e^s \Gamma(s)\,ds$. 
By L'H\^opital's rule, 
\[
\lim_{t\to\infty} \frac{\int_0^t e^s\gamma(s)\,ds}{e^t \Gamma(t)}
=\lim_{t\to\infty} \frac{e^t \Gamma(t)}{e^t\Gamma(t)+e^t\Gamma'(t)}
=\lim_{t\to\infty} \frac{1}{1+\Gamma'(t)/\Gamma(t)}=1.
\]
Since $\gamma\sim \Gamma$, the integral in \eqref{eq.yrvbound} is $O(\gamma(t))$, and thus $y(t)=O(\gamma(t))$, which is (B).

For part (II) in the regularly varying case, note that if $\gamma$ is regularly varying, it is also subexponential, so (A) and (B) are equivalent by Theorem \ref{lemma.fthetaymonotone}. Thus it suffices to prove the equivalence of (B) and (C). Once again, let $z=x-y$ and $\varphi(t)=(1+b)y(t)+ay(qt)$. Then $z$ obeys \eqref{eq.z} and $\varphi(t)=o(\gamma(t))$ as $t\to\infty$, by hypothesis. Revisiting the proof of part (C) of Theorem~\ref{thm.growthbounds2}, it can be shown in a similar way that  $z(t)=o(\gamma(t))$ as $t\to\infty$. Hence $x(t)=o(\gamma(t))$ as $t\to\infty$, so (B) implies (C). To show that (C) implies (B), note by hypothesis that $|x(t)| <\epsilon \gamma(t)$ as $t\to\infty$ for all $t\geq T$. Using the regular variation, this also gives $|x(qt)|\leq C\epsilon \gamma(t)$ for all $t\geq T/q=T_1$ and some $t$ and $\epsilon$ independent $C>0$. Now take $t\geq T_1>T$, and applying the triangle inequality to \eqref{eq.yrepx}, we have 
\begin{multline*}
	|y(t)|\leq |\xi|e^{-t}+\epsilon \gamma(t)+e^{-t}\left|\int_{0}^{T_1} e^s((1+b)x(s)+ax(qs))\,ds\right| \\
	+ \epsilon \left(|1+b| + |a|C\right) \int_{T_1}^t e^{-(t-s)} \gamma(s)\,ds.
\end{multline*}
Dividing by $\gamma(t)$ and taking the limit as $t\to\infty$ on the righthand side, we get limits $\epsilon$, $0$ and $\epsilon \left(|1+b| + |a|C\right)$ respectively for the three terms, using the fact that $\gamma$ is subexponential in the last case. Thus $\limsup_{t\to\infty} |y(t)|/\gamma(t)\leq \epsilon(1+|1+b| + |a|C)$,
so since $\epsilon$ is arbitrary, we have $y(t)=o(\gamma(t))$ as $t\to\infty$, as needed. 

For part (I) in the monotone case, (A) and (B) are equivalent by Theorem~\ref{lemma.fthetaymonotone}, so if (B) and (C) are equivalent, we are done. Let $z=x-y$ and $\varphi(t)=(1+b)y(t)+ay(qt)$. Then $z$ obeys \eqref{eq.z}. Suppose (B) holds, so $|y(t)|\leq C\gamma(t)$. Then, as $\gamma$ is increasing, $|\varphi(t)|\leq C(|1+b|+|a|)\gamma(t)$. Hence by Theorem \ref{thm.monotonepert}, $z(t)=O(\gamma(t))$. But since $y(t)=O(\gamma(t))$ and $x=y+z$, we have $x(t)=O(\gamma(t))$ as $t\to\infty$, or (B).  
To show (C) implies (B), note from (C) that $|x(t)|\leq K\gamma(t)$ for $t\geq 0$. Applying the triangle inequality to \eqref{eq.yrepx}, using this estimate, and using the monotonicity of $\gamma$, we get 
\begin{align*}
	|y(t)|&\leq |\xi|e^{-t}+K\gamma(t)+\int_0^t e^{-(t-s)}\left(|1+b|K\gamma(s)+|a|K\gamma(qs)\right)\,ds\\
	&\leq |\xi|e^{-t}+K\gamma(t)+\int_0^t e^{-(t-s)}\left(|1+b|K+|a|K\right)\,ds \cdot \gamma(t)\\
	&\leq (|\xi|/\gamma(0)+K+ |1+b|K+|a|K)\gamma(t),
\end{align*}
which is (B). The proof of (II) in the monotone case uses  ideas from parts (I) in the monotone and (II) in the regularly varying case, %, except that we use Theorem~\ref{thm.littleomonotonepert} and Lemma~\ref{lemma.fthetaymonotonelittleo}. 
which is left to the reader.
\end{proof}

%Taking parts (III) and (IV) together, and parts (I) and (II), together we can identify to leading order the maximal growth rate, or slowest decay rate, of the solution exactly. 
%\begin{theorem}
%	Suppose that $q\in (0,1)$ and \eqref{eq.unforcedasstable} holds. Let $f\in C([0,\infty);\mathbb{R})$ and $x$ be the unique continuous solution to \eqref{eq.forced}. 
%	Let $\gamma$ obey one of  
%	\begin{gather*}
%		\text{$\gamma$ be positive, nondecreasing and in $C^1(0,\infty)$, with $\gamma(t)\to\infty$ as $t\to\infty$};\\
%		\gamma\in \text{RV}_\infty(\eta), \quad \eta>\kappa.
%	\end{gather*}
%	Then the following are equivalent:
%	\begin{enumerate}
%		\item[(A)] There exists $\theta'>0$ such that $\limsup_{t\to\infty} |\int_{t-\theta'}^t f(s)\,ds|/\gamma(t)>0$, while $\left|\int_{(t-\theta)^+}^t f(s)\,ds \right|\leq K \gamma(t)$ for all $\theta\in [0,1]$, $t\geq 0$;
%		\item[(B)] $\limsup_{t\to\infty} |y(t)|/\gamma(t)\in (0,\infty)$;
%		\item[(C)] $\limsup_{t\to\infty} |x(t)|/\gamma(t)\in (0,\infty)$.
%	\end{enumerate}
%\end{theorem}
For smaller perturbations, it is hard to prove as clean a result. If $\gamma$ is regularly varying, and $x(t)=O(\gamma(t))$,  using \eqref{eq.yrepx} yields $y(t)=O(\gamma(t))$ as $t\to\infty$. In particular, take $\gamma(t)=t^\kappa$. However, this over--estimates the size of the smallest perturbation needed to preserve the rate of decay $t^\kappa$ of \eqref{eq.unforced}, since we know that if $y(t)=O(\gamma(t))$ where $\psi(t):=\gamma(t)/t^{\kappa+1}$ is regularly varying with index $-1$, and is integrable, then $x(t)=O(t^\kappa)$ as $t\to\infty$. But the integrability of $\psi$ means that $t\psi(t)\to 0$ as $t\to\infty$, and so the largest perturbations which preserve the rate of decay are $o(t^\kappa)$ (without further restrictions being placed on $y$). 

\begin{theorem} \label{thm.smallstochpert}
Suppose that $q\in (0,1)$ and \eqref{eq.unforcedasstable} holds. Let $f\in C([0,\infty);\mathbb{R})$ and $x$ be the unique continuous solution to \eqref{eq.forced}. Let $\gamma\in \text{RV}_\infty(\eta)$. 
%$\eta\leq \kappa$. 
%Suppose $\tilde{S}$ is given by \eqref{eq.Stilde}.
%Let $\gamma$ be regularly varying with index $\eta>\kappa$. 
\begin{itemize}
	\item[(i)] Let $y(t)=\text{O}(\gamma(t))$ as $t\to\infty$.
	\begin{enumerate}
		\item[(A)] 
		If $\eta< \kappa$, then $x(t)=O(t^\kappa)$ as $t\to\infty$.
		\item[(B)] If $\eta=\kappa$ and  $t\mapsto \gamma(t)/t^{\kappa+1}$ is in $L^1$, then 
		$x(t)=O(t^\kappa)$ as $t\to\infty$.
		Moreover, if $\limsup_{t\to\infty} |y(t)|/\gamma(t)\in (0,\infty)$, then $\limsup_{t\to\infty} |x(t)|/\gamma(t)>0$.
		\item[(C)] If $\eta=\kappa$ and  $t\mapsto\gamma(t)/t^{\kappa+1}$ is not in $L^1$, then 
		\[
		x(t)=\text{O}\left(t^\kappa\int_1^t \frac{\gamma(s)}{s^{\kappa+1}}\,ds\right), \,t\to\infty.
		\]
		Moreover $\limsup_{t\to\infty} |y(t)|/\gamma(t)>0$ implies  $\limsup_{t\to\infty} |x(t)|/\gamma(t)>0$.  
		%	\item[(D)] If $\eta>\kappa$ then the estimate in part (C) holds, which is equivalent to $x(t)=O(\gamma(t))$, as $t\to\infty$.
	\end{enumerate}  
	\item[(ii)] \hspace{0.1in} If $x(t)=\text{O}(\gamma(t))$ as $t\to\infty$, then $y(t)=\text{O}(\gamma(t))$ as $t\to\infty$. Likewise, if $x(t)=\text{o}(\gamma(t))$ as $t\to\infty$, then $y(t)=\text{o}(\gamma(t))$ as $t\to\infty$.   
	%	\item[(C)] We have 
	%	\[
	%\limsup_{t\to\infty} \frac{|Y(t)|}{t^{\kappa+\epsilon}}=0, \,\text{a.s.}\,
	%\Longleftrightarrow 
	%	\limsup_{t\to\infty} \frac{\log |Y(t)|}{\log t}\leq \kappa, \,\text{a.s.}\,
	%	\Longleftrightarrow \limsup_{t\to\infty} \frac{\log |X(t)|}{\log t}\leq \kappa,\, \text{a.s.}
	%	\]
\end{itemize}
\end{theorem}
Notice that the hypotheses on $y$ required to classify the asymptotic behaviour of $x$ are all equivalent to hypotheses on  $f$, as laid out in Theorem~\ref{thm.detlargepert}. Proofs of upper bounds in part (i) follow by considering $z=x-y$, noting that $y=O(\gamma)$ implies $\varphi=O(\gamma)$ and then applying Theorem \ref{thm.growthbounds2} to $z$. The proofs of lower bounds are by contradiction. Note first that (ii) can be shown by using \eqref{eq.yrepx}. Then, contradiction of the desired conclusion leads to supposing that $x(t)=o(\gamma(t))$. But then $y(t)=o(\gamma(t))$ by \eqref{eq.yrepx}, contradicting the hypothesis that $\limsup_{t\to\infty} |y(t)|/\gamma(t)>0$.

We now make some remarks. First, observe that, taken together, Theorems~\ref{thm.smallstochpert} and \ref{thm.detlargepert} cover a very wide class of perturbation sizes, and are in fact very   comprehensive for practical purposes. 

Second, note that part (A) and (B) of the theorem are in the form of a Hartman--Wintner type--result (see e.g.,  \cite{HW55} for the first results by Hartman and Wintner, \cite[Cor X.16.4]{Hart2002} for a typical result on ordinary equations and for linear functional equations, see Pituk~\cite{Pituk2002}). A typical feature of Hartman--Wintner results is that if perturbation terms obeys some integral smallness condition, then the solution of the perturbed equation $x$ can be related to the solution of the unperturbed problem in the form $x=O(u)$, or perhaps $x/u$ tends to a finite limit; we see therefore that parts (A) and (B) of the theorem have this character. Clearly, Hartman--Wintner results are a strong type of ``asymptotic preservation'' result. 

A weaker, but still very useful asymptotic preservation result, is a Perron-type result (for a typical result for linear, or linearisable, autonomous functional differential equations, see e.g., Pituk~\cite{Pituk2006}). In that case,  we show that some nonlinear function of the solution behaves like a known function of $t$. This allows us ignore dependence on initial data in bounds or limits (which can be hard to quantify explicitly anyway), and enables us to admit more irregularity (or to get away with weaker estimation) in the perturbed terms. For linear equations (or equations that can be linearised) Perron--type results often concentrate on (top) Liapunov exponents i.e. $\limsup_{t\to\infty} \log|x(t)|/t$ or $\lim_{t\to\infty} \log |x(t)|/t$, if the limit exists. We see that the existence of these limits is weaker than knowing that $x(t)\sim ce^{\lambda t}$ as $t\to\infty$ for some $c\neq 0$, and this is the price that must typically be paid in Perron--type results. In our case, because the solution has power--law asymptotic behaviour in the unperturbed case, we consider $\limsup_{t\to\infty} \log |x(t)|/\log t$, noting that for almost all solutions of the unperturbed equation we have 
\[
\limsup_{t\to\infty} \frac{\log |u(t)|}{\log t} =\kappa.
\]
We also notice that this is consistent with classic ideas of polynomial asymptotic stability in stochastic equations, due to Mao~\cite{Mao1992a,Mao1992b} and generalisations to other (in general, non--exponential) decay rates; see e.g.,  
Wu and Hu~\cite{WuHu2012}. In addition, we notice that considering solutions on this logarithmic scale is implicit in the bounds obtained for general delay differential equations in Cermak~\cite{cermak}. 

%\bibitem{HW55}
%P.~Hartman and A.~Wintner,  Asymptotic integration of ordinary nonlinear differential equations, {\em Amer. J. Math.} 77, 692--724, 1955.
%
%\bibitem{Hart2002} Ph.~Hartman, Ordinary differential equations, second ed., SIAM, Philadelphia, 2002.
%
%\bibitem{Pituk2002}
% M.~Pituk, The Hartman--Wintner theorem for functional differential equations, {\em J. Differential
%Equations}, 155, no. 1, 1--16, 1999. 
%
%\cite{Pituk2006}
%M.~Pituk, A Perron type theorem for functional differential equations, {\em J. Math. Anal. Appl.}, 316 (1), 24--41, 2006.

By focussing on the size of $\log |x(t)|$ in this way, we may characterise solution sizes very cleanly in terms of $y$, for both large and small perturbations, in the following Perron--type theorem. 
\begin{theorem} \label{thm.perronthm}
Suppose that $q\in (0,1)$ and \eqref{eq.unforcedasstable} holds. Let $f\in C([0,\infty);\mathbb{R})$ and $x$ be the unique continuous solution to \eqref{eq.forced}. Let $|\theta(t)|\nearrow  \infty$ as $t\to\infty$, and suppose $\theta(t)/\log t\to \eta\in [\kappa,\infty]$ as $t\to\infty$.  
\begin{enumerate}
	\item[(A)] If $\eta=\kappa$, then  
	\[
	\limsup_{t\to\infty} \frac{\log |y(t)|}{\theta(t)}\leq 1\, 
	\Longleftrightarrow\, \limsup_{t\to\infty} \frac{\log |x(t)|}{\theta(t)}\leq 1.
	\]
	\item[(B)] If $\eta\in (\kappa,\infty]$, then 
	\[
	\limsup_{t\to\infty} \frac{\log |y(t)|}{\theta(t)}=1\, 
	\Longleftrightarrow\, \limsup_{t\to\infty} \frac{\log |x(t)|}{\theta(t)}=1.
	\]
\end{enumerate}
\end{theorem}
Part (A) is equivalent to the more easily--understood statement 
\[
\limsup_{t\to\infty} \frac{\log |y(t)|}{\log t}\leq \kappa\, 
\Longleftrightarrow\, \limsup_{t\to\infty} \frac{\log |x(t)|}{\log t}\leq \kappa,
\] 
and part (B), for finite $\eta>\kappa$,  is equivalent to 
\[
\limsup_{t\to\infty} \frac{\log |y(t)|}{\log t}=\eta\, 
\Longleftrightarrow\, \limsup_{t\to\infty} \frac{\log |x(t)|}{\log t}=\eta,
\]
but we have stated Theorem~\ref{thm.perronthm} in the form above to show unity between the small and large perturbation cases, irrespective of the  solution or perturbation size. 

To prove part (B) of Theorem~\ref{thm.perronthm}, we break it into the cases when $\eta\in (\kappa,\infty)$ and $\eta=\infty$. We prove the former case, leaving the similar proof in the latter case to the interested reader. The reader will immediately note that the following proof requires mainly a re--interpretation or re--working of parts of earlier results, but besides this, no important new analysis is needed. 

If $\eta>\kappa$, the first statement in (B) means that for every $\epsilon>0$ there is a $C_\epsilon>0$ such that $|y(t)|\leq C_\epsilon t^{\eta+\epsilon}$ for all $t\geq 1$. Because $\eta>\kappa$, we can apply Theorem~\ref{thm.detlargepert} to give $|x(t)|\leq D_\epsilon t^{\eta+\epsilon}$ for all $t\geq 1$ and some $D_\epsilon>0$. As a result, we get that $\limsup_{t\to\infty} \log |x(t)|/\log t\leq \eta+\epsilon$. Letting $\epsilon\to 0$, we get that $\limsup_{t\to\infty} \log|x(t)|/\theta(t)\leq 1$. Suppose now, by way of contradiction, that the limsup is strictly less than unity; call it $l$. Then, for any $\epsilon>0$ so small that $(l+\epsilon)(\eta+\epsilon)<\eta$, there exists a $T(\epsilon)>0$ such that we have $\log |x(t)|< (l+\epsilon)\theta(t) < (l+\epsilon)(\eta+\epsilon) \log t$ for $t\geq T(\epsilon)$. Therefore, by Theorem~\ref{thm.detlargepert}, we have that there exists a $C_\epsilon>0$ such that $|y(t)|\leq C_\epsilon t^{(l+\epsilon)(\eta+\epsilon)}$ for all $t\geq 1$. Now, take logarithms, divide by $\log t$ and let $t\to\infty$; and noting that the inequality holds for all $\epsilon>0$ sufficiently small, we may further let $\epsilon\to 0^+$, giving 
\[
\limsup_{t\to\infty} \frac{\log |y(t)|}{\log t}\leq l\eta<\eta,
\]
while by hypothesis, this limsup is $\eta$, yielding the desired contradiction. Hence the forward implication holds. 

For the backward implication, the hypothesis implies that for every $\epsilon>0$ there is a $C_\epsilon>0$ such that $|x(t)|\leq C_\epsilon t^{\eta+\epsilon}$ for all $t\geq 1$. By \eqref{eq.yrepx} this implies that there exists $D_\epsilon>0$ such that $|y(t)|\leq D_\epsilon t^{\eta+\epsilon}$ for $t\geq 1$, and so, by taking logs, the limit as $t\to\infty$, and finally the limit as $\epsilon\to 0$ as before, we have $\limsup_{t\to\infty} \log |y(t)|/\theta(t)\leq  1$. As before, assume by way of contradiction that the limsup $l$ is strictly less than unity. Then for all $\epsilon$ sufficiently small we have 
$l+\epsilon<1$, and $(l+\epsilon)(\eta+\epsilon)<\eta$, and consequently for such an $\epsilon>0$ there is a $T(\epsilon)>0$ such that $\log|y(t)|<(l+\epsilon)\theta(t)< (l+\epsilon)(\eta+\epsilon)\log t$ for all $t\geq T(\epsilon)$.  

In the case that $l\eta<\kappa$, we may choose $\epsilon>0$ so that $(l+\epsilon)(\eta+\epsilon)\leq \kappa$; then by Theorem \ref{thm.smallstochpert}, we have that $x(t) = O(t^\kappa\log t)$, which means that $\limsup_{t\to\infty}\log x(t)/\log t\leq \kappa$. But $\limsup_{t\to\infty} \log |x(t)|/\log t=\eta$, so $\eta\leq \kappa$, which is false, and generates a contradiction in that case. 

In the other case, where  $l\eta\geq \kappa$, for every $\epsilon>0$ we have  $(l+\epsilon)(\eta+\epsilon)>\kappa$; in addition choose $\epsilon$ so small that $l+\epsilon<1$, and $(l+\epsilon)(\eta+\epsilon)<\eta$. Then by  Theorem~\ref{thm.detlargepert} we have that $|x(t)|\leq D_\epsilon t^{(l+\epsilon)(\eta+\epsilon)}$, which reduces to 
\[
\limsup_{t\to\infty} \frac{\log |x(t)|}{\log t}\leq (l+\epsilon)(\eta+\epsilon).
\]
By hypothesis, this limsup is $\eta$, so 
$\eta\leq  (l+\epsilon)(\eta+\epsilon)$. But $(l+\epsilon)(\eta+\epsilon)<\eta$, which generates the desired contradiction. We have thus shown (B) in the case where $l\eta\geq \kappa$.

The proof of part (A), where $\eta=\kappa$, is very similar to the argument given above, and again we leave it to the interested reader to complete.   

%In turns out that the analysis in part (D) is essentially unimprovable. Using the arguments used to prove Theorem~\ref{thm.smallstochpert}, we can show in the case that $\eta>\kappa$ and $y(t)=o(\gamma(t))$ as $t\to\infty$, that $x(t)=o(\gamma(t))$ as $t\to\infty$. Combining part (i)(D) of Theorem~\ref{thm.smallstochpert} with part (ii) and this newly claimed result, we have shown the following
%\begin{corollary}\label{cor.neccsuff}
%Suppose that $q\in (0,1)$ and \eqref{eq.unforcedasstable} holds. Let $f\in C([0,\infty);\mathbb{R})$ and $x$ be the unique continuous solution to \eqref{eq.forced}. Let $\gamma\in \text{RV}_\infty(\eta)$ with $\eta>\kappa$. Then 
%\begin{enumerate}
%	\item[(i)] $y(t)=o(\gamma(t))$ as $t\to\infty$ $\Longleftrightarrow$ $x(t)=o(\gamma(t))$ as $t\to\infty$
%	\item[(ii)] $y(t)=O(\gamma(t))$ as $t\to\infty$ $\Longleftrightarrow$ $x(t)=O(\gamma(t))$ as $t\to\infty$
%	\item[(iii)] $\limsup_{t\to\infty} |y(t)|/\gamma(t)\in (0,\infty)$  $\Longleftrightarrow$ 
%	$\limsup_{t\to\infty} |x(t)|/\gamma(t)\in (0,\infty)$.
%\end{enumerate} 	
%\end{corollary}

%
%\begin{theorem} 
%	Suppose that $q\in (0,1)$ and \eqref{eq.unforcedasstable} holds. Let $f\in C([0,\infty);\mathbb{R})$ and $x$ be the unique continuous solution to \eqref{eq.forced}. Let $\gamma\in \text{RV}_\infty(\eta)$, $\eta\leq \kappa$. 
%	\begin{enumerate}
%		\item[(A)] If $y(t)=\text{O}(\gamma(t))$, then 
%		\[
%		x(t)=\text{O}\left(t^\kappa\int_1^t \frac{\gamma(s)}{s^{\kappa+1}}\,ds\right), \quad t\to\infty.
%		\]
%		\item[(B)] If $x(t)=\text{O}(\gamma(t))$, then $y(t)=\text{O}(\gamma(t))$.
%	\end{enumerate}
%\end{theorem}
%Notice that the hypotheses on $y$ required to classify the asymptotic behaviour of $x$ are all equivalent to hypotheses on  $f$, as laid out in Theorem~\ref{thm.detlargepert}. Proofs of upper bounds in part (i) follow by considering $z=x-y$, noting that $y=O(\gamma)$ implies $\varphi=O(\gamma)$ and then applying Theorem \ref{thm.growthbounds2} to $z$. The proofs of lower bounds are by contradiction. Note first that (ii) can be shown by using \eqref{eq.yrepx}. Then, contradiction of the desired conclusion leads to supposing that $x(t)=o(\gamma(t))$. But then $y(t)=o(\gamma(t))$ by \eqref{eq.yrepx}, contradicting the hypothesis that $\limsup_{t\to\infty} |y(t)|/\gamma(t)>0$. 

The clean picture generated by the combination of the Hartman--Wintner type Theorem \ref{thm.smallstochpert}, the large perturbation characterisation Theorem~\ref{thm.detlargepert} and the Perron--type characterisation Theorem \ref{thm.perronthm} is spoiled by the inequality in part (A) of Theorem~\ref{thm.perronthm}, and the converse part (ii) in Theorem \ref{thm.detlargepert}. This converse does not give a result like 
\[
x(t)=O(t^\kappa), \, t\to\infty \Longrightarrow
y(t)=O(\gamma(t)), \,t\to\infty  \text{ where }  t\mapsto \gamma(t)/t^{\kappa+1}\in L^1(0,\infty),
\] 
which, together with Theorem~\ref{thm.smallstochpert}, would essentially characterise the maximal allowable size of perturbations; instead, the best we can do in the present work is to show that  
but rather $x(t)=O(t^\kappa)$ as $t\to\infty$ implies $y(t)=O(t^\kappa)$ as $t\to\infty$. 

Some scrutiny of part (ii) of  Theorem~\ref{thm.smallstochpert} may suggest a reason why a clean converse may be difficult to obtain. It suggests that if a solution of the perturbed equation decays faster than the unperturbed equation \eqref{eq.unforced}, this can only be achieved with a rapidly vanishing perturbation. It is not difficult to manufacture examples of this phenomenon. Let $\mu$ be a positive, continuous and integrable function.
%that tends to zero rapidly, in the sense that $\eta(t)/\int_t^\infty \eta(s)\,ds \to \infty$ as $t\to\infty$ (this hypothesis implies $\lim_{t\to\infty\infty} \log \int_t^\infty \eta(s)\,ds/t=-\infty$, which is consistent with faster than exponential decay to zero in $\eta$). 
Rescale $\mu$ so that $\int_0^\infty \mu(s)\,ds=1$. Let $x(0)=\xi\neq 0$ and suppose $f$ is given by 
\[
f(t)=\xi\left(-\mu(t)-b\int_t^\infty \mu(s)\,ds -a\int_{qt}^\infty \mu(s)\,ds\right), \quad  t\geq 0.
\] 
Then $x(t)=\xi \int_t^\infty \mu(s)\,ds$  for $t\geq 0$ is the unique continuous solution of \eqref{eq.forced}. If we suppose that $\mu$ tends to zero superexponentially, in the sense that $\lim_{t\to\infty} \log \mu(t)/t=-\infty$, then $x$ tends to zero superexponentially, and so does $f$; moreover, the more rapid the decay in $f$, the more rapid the decay in $x$. Thus, we may have the situation that the equation may be is perturbed by two ``small" perturbations, but that the smaller perturbation gives rise to a more slowly decaying solution. For instance, if $a>0$, $x(0)>0$ and $f(t)\equiv 0$, we have that there exist $0<c_1\leq c_2<+\infty$ such that $c_1t^\kappa \leq x(t)\leq c_2t^\kappa$ for all $t\geq 1$, while the solution in the above example, where there is a non--trivial perturbation, decays superexponentially fast when 
$\log \mu(t)/t\to-\infty$ as $t\to\infty$. 

The reader may also wonder whether the upper bound result in parts (i)(A) and (i)(B) of Theorem~\ref{thm.smallstochpert} can be improved in at least some situations to the stronger asymptotic estimate $x(t)\sim ct^\kappa$ for some $c\neq 0$. One situation where this seems especially plausible is that in which $f$ is positive, and $a>0$, since these conditions promote the positivity of the solution, and would tend to mitigate against the oscillation in the solution that is exhibited when $a<0$ and $f\equiv 0$. However, even in these situations where positivity of the solution is promoted, it is still not always possible to rule out the situation that $x(t)/t^\kappa$ is a bounded function which does not tend to a limit. Therefore, Theorem~\ref{thm.smallstochpert} is also sharp in terms of its conclusions; no better order estimate is available, granted the hypotheses imposed therein.   

To see that we cannot always expect solutions to be asymptotic to $t^\kappa$ as $t\to\infty$, consider the function $x$ defined by  
\[
x(t)=t^\kappa\left\{C+\sin\left(\frac{2\pi}{\log(1/q)}\log t\right)\right\}+Dt^\kappa\int_q^t \psi(s)\,ds, \quad t\geq q,
\]  
where $\psi\in RV_\infty(-1)$, $\psi\in L^1(0,\infty)$, and $C>1$, $D>0$. Let $a>0$, $b<0$ and $-b>a$. Let $\pi(t):=C+\sin(2\pi\log t/\log(1/q))$ for $t\geq q$. Define $f(t)=x'(t)-bx(t)-ax(qt)$ for all $t\geq 1$. Then $x$ solves \eqref{eq.forced} on $t\geq 1$ and $x(t)>0$ for all $t\geq q$. Moreover $x(t)/t^\kappa$ is bounded, but does not tend to a limit as $t\to\infty$. We now show that $f$ is asymptotically positive and satisfies the conditions of part (i)(A) of Theorem~\ref{thm.smallstochpert}. 

Note that $\pi(qt)=\pi(t)$ for all $t\geq 1$, so using $b+aq^\kappa=0$ and the uniform convergence theorem, we have that 
\[
\alpha(t):=-(bx(t)+ax(qt))=|b|Dt^\kappa\int_{qt}^t \psi(s)\,ds\sim |b|D\psi(t)t^{\kappa+1}\log(1/q),\quad t\to\infty.
\]
The function on the righthand side is positive, in $\text{RV}_\infty(\kappa)$, but is $o(t^\kappa)$ as
$t\to\infty$, because the integrability and asymptotic monotonicity of $\psi$ (the latter property follows from the regular variation of $\psi$) imply $t\psi(t)\to 0$ as $t\to\infty$. 

On the other hand, a direct calculation yields
\[
x'(t)=t^{\kappa-1} \left(\kappa\pi(t)+ \frac{2\pi}{\log(1/q)}\cos\left(\frac{2\pi}{\log(1/q)}\log t\right)
+D t\psi(t)+\kappa D \int_q^t \psi(s)\,ds\right).
\]
Note that the term in the brackets is uniformly bounded, since $\pi$ is bounded, $\psi$ is integrable, and $t\psi(t)\to 0$ as $t\to\infty$. Hence $x'(t)=O(t^{\kappa-1})$ as $t\to\infty$. Therefore, as $\alpha$ is asymptotic to a $RV_\infty(\kappa)$ function, $x'(t)=o(\alpha(t))$ as $t\to\infty$. Thus $f(t)=x'(t)+\alpha(t)\sim |b|D\log(1/q)\psi(t) t^{\kappa+1}$ as $t\to\infty$. From this it is easy to show that $y(t)\sim f(t)$ as $t\to\infty$. Therefore, we may take $\gamma(t):=|b|D\log(1/q)\psi(t) t^{\kappa+1}$ in Theorem~\ref{thm.smallstochpert} because it is in $\text{RV}_\infty(\kappa)$. Moreover, since $\psi$ is integrable, 
\[
\int_1^T \frac{\gamma(s)}{s^{\kappa+1}}\,ds
=|b|D\log(1/q)\int_1^T \psi(s)\,ds\to L, \quad T\to\infty,
\]
so although $y$ fulfills the properties of part (i)(B) of  Theorem~\ref{thm.smallstochpert}, with $a>0$ and $f$ positive for all $t$ sufficiently large, $t\to x(t)/t^\kappa$ does not tend to a constant as $t\to\infty$. This demonstrates, as claimed above, that the ``big O" result in Theorem~\ref{thm.smallstochpert} is the best possible in general, given the hypotheses imposed here. This example also hints at the possible asymptotic form of the bounded component of $x(t)/t^\kappa$, under appropriate smallness conditions on $y$, and we propose to investigate this further in later works.   

A characterisation, purely in terms of $f_\theta(t)$, of the situation wherein $\log |y(t)|$ has certain upper rates of growth or decay, has not been presented here. However, it seems that such a result should be available, and would be especially useful when considering Perron--type results for differential equations with general unbounded delay.  

%If we want to write an equivalence, this can be done by comparing upper polynomial Liapunov exponents, although some loss of sharpness is incurred. 
%
%\begin{theorem}
%	Suppose that $q\in (0,1)$ and \eqref{eq.unforcedasstable} holds. Let $f\in C([0,\infty);\mathbb{R})$ and $x$ be the unique continuous solution to \eqref{eq.forced}.
%	Then the following are equivalent
%	\begin{enumerate}
%		\item[(A)] There exists $\Lambda_y\leq \kappa$ such that 
%		$\limsup_{t\to\infty} \log|y(t)|/\log t=:\Lambda_y$; 
%		\item[(B)]  There exists $\Lambda_x\leq \kappa$ such that $\limsup_{t\to\infty} \log|x(t)|/\log t=:\Lambda_x$. 
%	\end{enumerate} 
%	Moreover, either implies $\Lambda_x\leq\max(\kappa,\Lambda_y)$ and $\Lambda_y\leq \Lambda_x$. 
%\end{theorem}
%\begin{proof}
%	Suppose (A). Then, for every $\epsilon>0$, $y(t)=\text{O}(t^{\Lambda_y+\epsilon})$. Set $\mu(t)=t^{\Lambda_y+\epsilon}\in \text{RV}_\infty(\Lambda_y+\epsilon)$. Hence for every $\epsilon>0$ we have 
%	$x(t)=\text{O}\left(t^\kappa\int_1^t s^{\Lambda_y+\epsilon-\kappa-1}\,ds\right)$. If $\Lambda_y<\kappa$, the integrand is integrable for $\epsilon$ sufficiently small. Since $\epsilon$ can be chosen arbitrarily, we get $x(t)=\text{O}(t^\kappa)$. Thus $\Lambda_x$ is well defined, and $\Lambda_x\leq \kappa$.
%	
%	If $\Lambda_y=\kappa$, the integrand behaves like $t^{-1+\epsilon}$, so the integral behaves like $t^\epsilon$, and hence $x(t)=\text{O}(t^{\kappa+\epsilon})$. Thus
%	$\limsup_{t\to\infty} \log|x(t)|/\log t\leq \kappa+\epsilon$. Since $\epsilon>0$ is arbitrary, we may let $\epsilon\to 0$, and so $\Lambda_x\leq \kappa$. Hence $\Lambda_x\leq \max(\kappa,\Lambda_y)$. Thus under each of the cases in (A), $\Lambda_x$ is well--defined, and  $\Lambda_x\leq \kappa$, which is precisely (B).
%	
%	Suppose now, statement (B) holds. Then for every $\epsilon>0$, $x(t)=\text{O}(t^{\Lambda_x+\epsilon})$. But then $y(t)=\text{O}(t^{\Lambda_x+\epsilon})$. Hence 
%	$\limsup_{t\to\infty} \log|y(t)|/\log t\leq \Lambda_x+\epsilon$. Letting $\epsilon\to 0$, we get $\Lambda_y\leq \Lambda_x\leq \kappa$, so claim (A) is satisfied, and the equivalence proven. 
%\end{proof}
Next, we compare the results of Theorems \ref{thm.detlargepert} and \ref{thm.smallstochpert}  with perturbation theorems of Cermak (see e.g., Theorem 2.3 and Example 4.1 in \cite{cermak}), which, in our view provide the sharpest and most comprehensive estimates for the asymptotic behaviour of perturbed functional differential equations with unbounded delay. He supposes that $f(t)=O(t^\beta)$ and $f'(t)=O(t^{\beta-1})$ as $t\to\infty$. Applying Cermak's results to the pantograph equation, we get  
\begin{align*}
x(t)&=O(t^\kappa), \quad t \to\infty, \quad \text{if $\beta<\kappa$},\\
x(t)&=O(t^\kappa\log t), \quad t \to\infty, \quad\text{if $\beta=\kappa$},\\
x(t)&=O(t^\beta), \quad t \to\infty, \quad\text{if $\beta>\kappa$}.
\end{align*}
First, we observe that our results need neither need pointwise estimates on $f$, nor on $f'$, but merely on $y$, where $y'=-y+f$, and the smoothness of $f$ is not needed in our results. Second, it can readily be the case that $f$ has very bad pointwise behaviour, yet $y$ has good behaviour. For example, suppose that 
\[
f(t)=t^{\beta}+e^{t}\sin(e^{2t}), \quad t\geq 0,
\] 
Clearly, $f(t)=O(e^t)$ and $f'(t)=O(e^{3t})$ as $t\to\infty$, yet $y(t)=O(t^\beta)$ as $t\to\infty$. Thus, applying Theorem~\ref{thm.smallstochpert}, we are able to conclude that 
\[
x(t)=O(t^{\max(\kappa,\beta)}), \quad t\to\infty, \text{for $\beta\neq \kappa$} 
\]
and $x(t)=O(t^\kappa\log t)$ as $t\to\infty$ if $\beta=\kappa$, while Cermak's theorem cannot give any conclusions.

Furthermore, the existing results require an explicit parametric comparison between the solution of the unperturbed equation and perturbation. For example, if $f(t)=O(t^\beta\log t)$ as $t\to\infty$, we cannot apply Theorem 2.3 in \cite{cermak} directly; instead, we would note that $f(t)=O(t^{\beta+\epsilon})$ for every $\epsilon>0$, and hope to obtain the estimate  $f'(t)=O(t^{\beta-1+\epsilon})$ on the derivative. Thus, if $\beta<\kappa$, we get $x(t)=O(t^\kappa)$ and if $\beta>\kappa$ we get $x(t)=O(t^\beta)$, but no prediction can be made in the case $\beta=\kappa$. Using Theorem~\ref{thm.smallstochpert} however,  we get 
\[
x(t)=O\left(t^\kappa\int_1^t s^{\beta-\kappa-1}\log s\,ds\right), \quad t\to\infty.
\]   
The estimates concur with \cite[Theorem 2.3]{cermak} in the cases when $\beta\neq \kappa$; but when $\beta=\kappa$, we get 
\[
x(t)=O(t^\kappa \log^2(t)), \quad t\to\infty,
\]
which is not covered by Cermak's theorem. 

These considerations seem to show that the estimation in the Fundamental Lemma~\ref{lemma.fundamental} may be slightly sharper in characterising the transition from small to large perturbations than the proof of \cite[Theorem 2.3]{cermak}. It is an interesting open question to revisit Cermak's delicate and beautiful proof of \cite[Lemma 3.3]{cermak} for equations with general unbounded delay. This result underlies \cite[Theorem 2.3]{cermak}, and is a close analogue of Lemma~\ref{lemma.fundamental}. This is especially worthwhile, because we have shown by means of the general examples in Theorem~\ref{thm.growthbdssharp} that the estimates in Theorem~\ref{thm.smallstochpert} are, in a certain sense, unimprovable.

Finally, we note that Theorem \ref{thm.detlargepert} above deals with large perturbations; suppose that $y(t)=O(e^{\beta t})$ as $t\to\infty$ for some $\beta>0$. Then this perturbation cannot be put into the framework of \cite[Theorem 2.3]{cermak}. However, by Theorem \ref{thm.detlargepert}, we have that $x(t)=O(e^{\beta t})$ as $t\to\infty$. Furthermore, if it is the case that $\limsup_{t\to\infty} |x(t)|/e^{\beta t}\in (0,\infty)$, this can happen if and only if  $\limsup_{t\to\infty} |y(t)|/e^{\beta t}\in (0,\infty)$. However, stronger conclusions about $f$ are not available; for instance consider 
\[
f(t)=e^{\beta t} + ce^{\theta t}\sin(e^{3\theta t})
\]    
where $\theta>\beta>0$. If $c=0$, then $f(t)=O(e^{\beta t})$ and $y(t)=O(e^{\beta t})$ as $t\to\infty$. If $c\neq 0$, then we still have 
$y(t)=O(e^{\beta t})$ as $t\to\infty$, but $f(t)=O(e^{\theta t})$. Thus, both $f$'s generate the same asymptotic size of $y$, and hence of $x$, but knowing the size of $x$ does not allow conclusions to be drawn about the order of $f$, merely that of $y$, or equivalently, by Theorem 6, the order of $f_\theta$.  

In the case of small perturbations i.e. when $y(t)=O(\gamma(t))$ as $t\to\infty$ and $t\mapsto\gamma(t)/t^{\kappa+1}$ in integrable, we have that $x(t)=O(t^\kappa)$ as $t\to\infty$, or to put it another way, $t\mapsto x(t)/t^\kappa$ is a bounded function. In the unperturbed case, the asymptotic  structure of this bounded component has been found; it is the sum of a log--periodic function and a function which tends to zero as $t\to\infty$ (see \cite{KatoMc}). It is the subject of further work to determine whether these small perturbation conditions are sufficient to establish that the bounded component of the perturbed equation retains a similar (log--periodic) structure. 

Some preliminary ideas as to how this might be achieved are presented in Section 5 and at the end of the paper in Section 7. Space constraints, and the necessity to develop a rate characterisation for related stochastic differential equations (see the next section) preclude more exact analysis and discussion in this work. 

Thus, we see that a correct determination of the exact asymptotic order is likely a first step in understanding the detailed asymptotic structure of solutions. Moreover, it opens up the prospect that the solution \textit{might retain the fine structure of its solutions despite bad pointwise behaviour in forcing functions, irregular stochastic terms, and relatively large forcing terms (as measured by $y$)}. In other words, the refined asymptotic bounds presented in this section can be seen as a critical first step in establishing that the solution structure is robust to relatively large and irregular perturbations.     

\vspace{-0.25in}

\section{Stochastically forced equation and auxiliary SDEs}
In this section, we consider the stochastically perturbed pantograph equations
\begin{equation} \label{eq.stochforced2}
dX(t)=(bX(t)+aX(qt)+f(t))\,dt + \sigma(t)\,dB(t), \quad t\geq 0; \quad X(0)=\xi,
\end{equation}
where once again $a$ and $b$ are real numbers and $q\in (0,1)$. For simplicity, we assume $\xi$ is non--random, $f$ is continuous as before, and   
\begin{equation} \label{eq.sigma}
\sigma\in C([0,\infty);\mathbb{R}).
\end{equation}
Here $B=\{B(t):t\geq 0\}$ is a standard Brownian motion. In order that a solution of \eqref{eq.stochforced2} is well--defined and has appropriate properties, we work on a complete probability space $(\Omega,\mathcal{F},\mathbb{P})$, where Brownian motion $B$ generates its own natural filtration $\mathcal{F}^B(t)=\sigma(\{B(s):0\leq s\leq t\})$ with  $\mathcal{F}^B(t)\subset \mathcal{F}$ for each $t\geq 0$. Under these assumptions, \eqref{eq.stochforced2} is considered a shorthand for the integral equation 
\[
X(t)=\xi+\int_{0}^t \{bX(s)+aX(qs) +f(s)\}\,ds + \int_0^t \sigma(s)\,dB(s), \quad t\geq 0
\]
which has a unique, continuous, and ($\mathcal{F}^B(t)$--) adapted solution on $[0,\infty)$ (see e.g., Mao~\cite{MaoBK}). 
%We can also consider a stochastically and deterministically forced equation; let $f$ obey \eqref{eq.fcns} and $\sigma$ obey \eqref{eq.sigma}, and consider with the same notation as above, the equation 
%\begin{equation} \label{eq.stochforced2}
%	d\tilde{X}(t)=(b\tilde{X}(t)+a\tilde{X}(qt)+f(t))\,dt + \sigma(t)\,dB(t), \quad t\geq 0; \quad \tilde{X}(0)=\xi.
%\end{equation}
%Once again, interpreting this as an integral equation 
%\[
%\tilde{X}(t)=\xi+\int_{0}^t \{b\tilde{X}(s)+a\tilde{X}(qs)+f(s)\}\,ds + \int_0^t %\sigma(s)\,dB(s), \quad t\geq 0,
%\]
%we have that \eqref{eq.stochforced2} has a unique, continuous, and adapted solution. 
%$X$ and $\tilde{X}$ are very closely related. Consider $D(t)=\tilde{X}(t)-X(t)$ for $t\geq 0$. Then $D$ is a continuous and adapted process with $D(0)=0$; indeed, $D$ is in $C^1((0,\infty);\mathbb{R})$ and obeys 
%\[
%D'(t)=aD(qt)+bD(t)+f(t), \quad t\geq 0; \quad D(0)=0,
%\] 
%which is exactly an equation of the form \eqref{eq.forced}, to which the theory of the last section can be applied.  

In just the way the auxiliary ordinary differential equation \eqref{eq.y} allows sharp results to be proven for the deterministically forced equation \eqref{eq.forced}, solutions of related auxiliary stochastic (ordinary) differential equations enable us to characterise the asymptotic behaviour of solution of \eqref{eq.stochforced2}. Let $Y_0$ and $Y$ be solutions of the SDEs
\begin{gather} \label{eq.Y0}
dY_0(t)=-Y_0(t)\,dt+\sigma(t)\,dB(t), \quad t\geq 0; \quad Y_0(0)=0,	\\
\label{eq.Y}
dY(t)=(-Y(t)+f(t))\,dt + \sigma(t)\,dB(t), \quad t\geq 0; \quad Y(0)=0.
\end{gather}
Both equations possess unique, continuous, adapted solutions. Recalling that $y$ is the solution of \eqref{eq.y}, the representations of the solutions of \eqref{eq.Y0} and \eqref{eq.Y} are 
\begin{eqnarray} \label{eq.Y0rep}
Y_0(t)&=&e^{-t}\int_0^t e^s \sigma(s)\,dB(s), \quad t\geq 0, \\
\label{eq.Yrep}
Y(t)&=&e^{-t}\int_0^t e^s f(s)\,ds+e^{-t}\int_0^t e^s \sigma(s)\,dB(s) = y(t)+Y_0(t), \quad t\geq 0.
\end{eqnarray} 
We show momentarily that the pathwise asymptotic behaviour of \eqref{eq.Y0} and \eqref{eq.Y} can be characterised very precisely. For now, we show how 
%$Y_0$ can be connected with $X$ and 
$Y$ can be connected with $X$. Let 
%$Z_0=X-Y_0$ and 
$Z=X-Y$. 
%Let $\phi_0(t)=aY_0(qt)+(1+b)Y_0(t)$, for $t\geq 0$. Then $Z_0(0)=\xi$ and $Z_0$ obeys an equation of the form \eqref{eq.forced}: 
%\begin{equation} \label{eq.Z0}
%Z_0'(t)=aX(qt)+bX(t)+Y_0(t)=aZ_0(qt)+bZ_0(t)+\phi_0(t), \quad t\geq 0. 
%\end{equation}
Set $\phi_1(t)=ay(qt)+(1+b)y(t)$ for $t\geq 0$ and %and $\phi(t)= aY(qt)+(1+b)Y(t)$ for $t\geq 0$, we have 
\begin{equation} \label{eq.phistoch}
\phi(t)= aY(qt)+(1+b)Y(t), \quad t\geq 0
\end{equation}
so $\phi(t)=\phi_0(t)+\phi_1(t)$ where  $\phi_0(t)=aY_0(qt)+(1+b)Y_0(t)$ and  
\begin{equation} \label{eq.Z}
Z'(t)=aX(qt)+bX(t)+Y(t) = aZ(qt)+bZ(t)+\phi(t), \quad t\geq 0.
\end{equation}
Thus, we can apply the machinery of the previous section to $Z$, % and $Z_0$, 
once the asymptotic behaviour of $\phi$ is known. This will be the case, since we will good information about the asymptotic behaviour of $y$ and $Y_0$, on which $\phi$ entirely depends. Furthermore, since %$X=Z+Y_0$ and 
$X=Z+Y$, the asymptotic behaviour of %$X$ and 
$X$ can be found. 

In particular, suppose $b<0$ and $|a|<|b|$ and that $Y(t)\to 0$ as $t\to\infty$ a.s. Then $\phi(t)\to 0$ as $t\to\infty$ a.s. Applying Theorem \ref{thm.1} part (a), we see from \eqref{eq.Z} that $Z(t)\to 0$ as $t\to\infty$ a.s. Thus, $X(t)=Z(t)+Y(t)\to 0$ as $t\to\infty$ a.s. In the case that $Y$ is a.s. bounded, $\phi$ is bounded a.s., and this leads to $Z$ being bounded a.s., by applying part (b) of Theorem \ref{thm.1} to \eqref{eq.Z}. Hence $X=Z+Y$ is bounded a.s. 

%Analogously, if $b<0$ and $|a|<|b|$ and $Y(t)\to 0$ as $t\to\infty$ a.s., then $\phi_0(t)\to 0$ as $t\to\infty$ a.s. Applying Theorem \ref{thm.1} part (a), we see from \eqref{eq.Z} that $Z(t)\to 0$ as $t\to\infty$ a.s. Thus, $X(t)=Z(t)+Y(t)\to 0$ as $t\to\infty$ a.s.  In the case that $Y$ is a.s. bounded, $\phi$ is bounded a.s., and this leads to $Z$ being bounded a.s., by applying part (b) of Theorem \ref{thm.1} to \eqref{eq.Z}. Hence $X=Z+Y$ is bounded a.s. 

Converses result by writing $Y$ in terms of $X$. 
Write $dX(t)=-X(t)+ f(t)+\{(1+b)X(t)+aX(qt)\}\,dt + \sigma(t)\,dB(t)$. 
By stochastic integration by parts, we get    
%establish
%\[
%X(t)=\xi e^{-t}+\int_0^t e^{-(t-s)}\{(1+b)X(s)+aX(qs)\}\,ds + Y(t),
%\]
%so rearranging gives 
%\begin{equation} \label{eq.Y0repX}
%Y(t)=X(t)-\xi e^{-t} - \int_0^t e^{-(t-s)}\{(1+b)X(s)+aX(qs)\}\,ds, \quad t\geq 0.
%\end{equation}
%Applying the same approach to $X$, we get 
%\begin{multline*}
%X(t)= \xi e^{-t} + \int_0^t  e^{-(t-s)}\{(1+b)X(s)+aX(qs)\}\,ds\\ + \int_0^t %e^{-(t-s)}f(s)\,ds) + \int_0^t e^{-(t-s)}\sigma(s)\,dB(s), 
%\end{multline*}
%and since the last two terms are $Y$, rearranging gives 
\begin{equation}  \label{eq.YreptildeX}
Y(t)=X(t)-\xi e^{-t} - \int_0^t e^{-(t-s)}\{(1+b)X(s)+aX(qs)\}\,ds, \quad t\geq 0.
\end{equation}
%From \eqref{eq.YrepX}, it is evident that if $X(t)\to 0$ as $t\to\infty$ a.s., then 
%$Y(t)\to 0$ as $t\to\infty$ a.s., and likewise, if $X$ is bounded a.s., then $Y$ is bounded a.s. 
From \eqref{eq.YreptildeX}, it is clear that if $X(t)\to 0$ as $t\to\infty$, then $Y(t)\to 0$ as $t\to\infty$ a.s., and likewise, if $X$ is bounded a.s., then $Y$ is bounded a.s. 

%The above discussion therefore yields the following theorems. 
%
%\begin{theorem}\label{thm.stochstabX1}
%	Suppose that $q\in (0,1)$ and \eqref{eq.unforcedasstable} holds. Let $\sigma\in C([0,\infty);\mathbb{R})$ and $X$ be the unique continuous adapted solution to \eqref{eq.stochforced1}. If $Y_0$ is the unique continuous adapted solution to \eqref{eq.Y0}, then the following are equivalent:
%	\begin{enumerate}
%		\item[(A)] $Y_0(t)\to 0$ as $t\to\infty$ a.s.;
%		\item[(B)] $X(t)\to 0$ as $t\to\infty$ a.s. 
%	\end{enumerate}
%\end{theorem}
%
%\begin{theorem}\label{thm.stochbddX1}
%	Suppose that $q\in (0,1)$ and \eqref{eq.unforcedasstable} holds.  Let $\sigma\in C([0,\infty);\mathbb{R})$ and $X$ be the unique continuous adapted solution to \eqref{eq.stochforced1}. If $Y_0$ is the unique continuous adapted solution to \eqref{eq.Y0}, then the following are equivalent:
%	\begin{enumerate}
%		\item[(A)] $Y_0$ is bounded, a.s.;
%		\item[(B)] $X$ is bounded, a.s. 
%	\end{enumerate}
%\end{theorem}
%
%\begin{theorem}\label{thm.stochstabtilX1}
%	Suppose that $q\in (0,1)$ and \eqref{eq.unforcedasstable} holds. Let $f,\sigma\in C([0,\infty);\mathbb{R})$ and $X$ and $Y$ be the unique continuous adapted solutions to \eqref{eq.stochforced2} and \eqref{eq.Y}. Then
%	\begin{enumerate}
%		\item[(i)] $Y(t)\to 0$ as $t\to\infty$ a.s. $\Longleftrightarrow$ $X(t)\to 0$ as $t\to\infty$ a.s.;
%		\item[(ii)] $Y$ is bounded a.s. $\Longleftrightarrow$  $X$ is bounded a.s. 
%		\item[(iii)] $\limsup_{t\to\infty} |Y(t)|\in (0,\infty)$ a.s. $\Longleftrightarrow$ 
%		$\limsup_{t\to\infty} |X(t)|\in (0,\infty)$ a.s.
%	\end{enumerate}
%\end{theorem}
%
%\begin{theorem}\label{thm.stochbddtilX1}
%	Suppose that $q\in (0,1)$ and \eqref{eq.unforcedasstable} holds. Let $f,\sigma\in C([0,\infty);\mathbb{R})$ and $X$ be the unique continuous adapted solution to \eqref{eq.stochforced2}. If $Y$ is the unique continuous adapted solution to \eqref{eq.Y}, then the following are equivalent:
%	\begin{enumerate}
%		\item[(A)] $Y$ is bounded a.s.;
%		\item[(B)] $X$ is bounded a.s. 
%	\end{enumerate}
%\end{theorem}

The above discussion shows the importance of possessing good asymptotic information concerning $Y_0$ and $Y$. Fortunately, some of this is already known for $Y_0$, and using known results, more can be established readily from scratch for $Y$. The following was established concerning $Y_0$ in \cite{ACR:2011(DCDS)}, generalising results from \cite{ChanWilliams:1989}.
\begin{theorem}\label{thm.Y0asy}
Suppose $\sigma$ obeys \eqref{eq.sigma}. Define $S(\epsilon)$ for $\epsilon>0$ by \eqref{eq.S} (if the sum  diverges, we take $S(\epsilon)=+\infty$)
and let $Y_0$ be the unique continuous adapted solution of \eqref{eq.Y0}. 
\begin{enumerate}
	\item[(a)] $S(\epsilon)<+\infty$ for all $\epsilon>0$ $\Longleftrightarrow$ $Y_0(t)\to 0$ as $t\to\infty$ a.s.
	\item[(b)] There exists $\epsilon'>0$ such that $S(\epsilon)<+\infty$ for all $\epsilon>\epsilon'$ 
	$\Longleftrightarrow$ $Y_0$ is bounded a.s. 
	\item[(c)] (There exists $\epsilon'>0$ such that $S(\epsilon)<+\infty$ for all $\epsilon>\epsilon'$ and $S(\epsilon)=+\infty$ for all $\epsilon<\epsilon'$)  $\Longleftrightarrow$ $0<\limsup_{t\to\infty} |Y_0(t)|<+\infty$ a.s. 
	\item[(d)] $S(\epsilon)=+\infty$ for all $\epsilon>0$  $\Longleftrightarrow$ $\limsup_{t\to\infty} |Y_0(t)|=+\infty$ a.s. 
\end{enumerate}
\end{theorem}
It can be shown that $S$ can only obey assumption (A) in part (a), (c) and (d) of the theorem, so solutions of \eqref{eq.Y0} either tend to zero with probability one, are bounded but non--convergent with probability one, or are unbounded with probability one.  

Notice the asymptotic behaviour of $\sigma$, measured through the function $\sigma^2_1(t)$ entirely governs the asymptotic behaviour of $Y_0$.
%; by virtue of Theorem~\ref{thm.stochstabX1} and \ref{thm.stochbddX1}, this means that the a.s. convergence to zero or a.s. boundedness of solutions $X$ of the stochastic equation \eqref{eq.stochforced1} are governed by the decay in $\int_t^{t+1}\sigma^2(s)\,ds$. 
This mirrors the manner in which the asymptotic decay or boundedness of $f_\theta(t)$ governs the convergence or boundedness of solutions of the deterministic equations \eqref{eq.y} and \eqref{eq.forced}.   

%Since it is hard to see how the finiteness of $S(\epsilon)$ is related to decay rates in $\sigma$, by imposing extra control on $\sigma^2$, the conditions (A) in Theorem
%\ref{thm.Y0asy} can be replaced by more easily checked conditions. The following result again appears in \cite{ACR:2011(DCDS)}.
%
%\begin{theorem} \label{thm.Y0asysuff}
%Let $\sigma$ obey \eqref{eq.sigma}, $Y_0$ be the unique continuous adapted solution to \eqref{eq.Y0}, and suppose that 
%$L:=\lim_{n\to\infty} \int_n^{n+1} \sigma^2(s)\,ds \log n \in [0,\infty]$ exists. 
%\begin{enumerate}
%	\item[(A)] If $L=0$, then $Y_0(t)\to 0$ as $t\to\infty$ a.s.
%	\item[(B)] If $L\in (0,\infty)$, then $\limsup_{t\to\infty} |Y_0(t)|\in (0,\infty)$ a.s.
%	\item[(C)] If $L=+\infty$, then $\limsup_{t\to\infty} |Y_0(t)|=+\infty$ a.s.
%\end{enumerate}
%\end{theorem} 
It is also shown in \cite{ACR:2011(DCDS)} that if $\sigma^2_1(n)$ is asymptotic to a monotone (non--increasing) sequence, $Y_0(t)\to 0$ as $t\to\infty$ a.s. if and only if $\sigma^2_1(n) \log n \to 0$ as $n\to\infty$. In rough terms, therefore,  $\sigma^2_1(t)$ must tend to zero faster than $1/\log t$ in order for $Y_0(t)$ to tend to zero as $t\to\infty$, a.s. 

We now give a characterisation of the a.s. convergence and boundedness of $Y$. In broad terms, we have convergence to zero (respectively, boundedness) for the ``doubly'' perturbed equation if and only if both the ``singly''
perturbed equations tend to zero (respectively, are bounded).
\begin{theorem} \label{thm.Yto0bddchar}
Let $f,\sigma$ obey \eqref{eq.fcns} and \eqref{eq.sigma}. Let $Y$ be the unique continuous adapted solution to \eqref{eq.Y}, and suppose $S$ is given by \eqref{eq.S}. 
\begin{itemize}
	\item[(a)] The following are equivalent:
	\begin{enumerate}
		\item[(A)] There exists $\Delta>0$ such that $f_\theta(t) \to 0$ as $t\to\infty$ for all $\theta\in [0,\Delta]$, $S(\epsilon)<+\infty$ for all $\epsilon>0$;
		\item[(B)] $Y(t)\to 0$ as $t\to\infty$ a.s. 
	\end{enumerate}
	\item[(b)] The following are equivalent:
	\begin{enumerate}
		\item[(A)] There exists $\Delta>0$, $\bar{F}>0$ and $\epsilon'>0$ such that 
		\[
		\sup_{0\leq \theta\leq \Delta}|f_\theta(t)|\leq \bar{F}, \quad t\geq 0, \quad \text{and}\quad 
		S(\epsilon)<+\infty \quad \epsilon>\epsilon'.
		\]
		\item[(B)] $Y$ is a.s. bounded. 
	\end{enumerate}
\end{itemize} 
\end{theorem}
\begin{proof}
Note $Y(t)=Y_0(t)+y(t)$ for $t\geq 0$, where $y$ is the solution of \eqref{eq.y}. 
For part (a), we prove first  (A) implies (B). The first part of (A) gives $y(t)\to 0$ as $t\to\infty$, by the implication in Theorem~\ref{thm.detasymptoticstab} that (B) implies (A). The second part of (A) gives $Y_0(t)\to 0$ as $t\to\infty$ a.s., by part (a) of Theorem~\ref{thm.Y0asy}. Hence $Y(t)\to 0$ as $t\to\infty$.

To show (B) implies (A), note $Y(t)$ is Gaussian distributed for each $t$, with mean $y(t)$. Consider any sequence $t_n\uparrow \infty$. Then $Y_n:=Y(t_n)$ converges to zero a.s., by (B). But since $Y_n$ is Gaussian, it means that it also converges in distribution, which forces $\mathbb{E}[Y_n]\to 0$ as $n\to\infty$, or $y(t_n)\to 0$ as $n\to\infty$. Suppose now, by way of contradiction, that $\limsup_{t\to\infty} |y(t)|>0$. 
Then there exists $c>0$ and a sequence $t_n\uparrow \infty$ such that $|y(t_n)|>c$ for all $n$ sufficiently large. But this contradicts the fact that (B) implies $y(t_n)\to 0$ as $n\to\infty$ for any divergent sequence $(t_n)$. Hence, we must have $\limsup_{t\to\infty} |y(t)|=0$, or $y(t)\to 0$ as $t\to\infty$. But by the implication in Theorem~\ref{thm.detasymptoticstab} that (A) implies (B), we must have that $\int_t^{t+\theta} f(s)\,ds\to 0$ as $t\to\infty$, which proves the first part of (A). To prove the second part, note we have shown $y(t)\to 0$ as $t\to\infty$, and assumed $Y(t)\to 0$ as $t\to\infty$ a.s. Hence $Y_0(t)=Y(t)-y(t)\to 0$ as $t\to\infty$ a.s. Now apply part (a) of Theorem \ref{thm.Y0asy} to see that  we must have $S(\epsilon)<+\infty$ for all $\epsilon>0$, which is the second part of (A). Thus we have shown that (B) implies (A), establishing the required equivalence to complete the proof of part (a).  

For part (b), we prove first (A) implies (B). Note the first part of (A) implies $y$ is bounded, applying Theorem~\ref{thm.detasymptoticbound}
with  $\gamma(t)\equiv 1$. The second part of (A) implies $Y_0$ is a.s. bounded, by part (b) of Theorem \ref{thm.Y0asy}. Thus $Y=Y_0+y$ is a.s. bounded, giving (B). 

To show that (B) implies (A), it is enough to show that $y$ is bounded; in that case the first part of (A) holds, by the converse half of Theorem~\ref{thm.detasymptoticbound}
with  $\gamma(t)\equiv 1$. On the other hand, if $y$ is bounded, then $Y_0=Y-y$ is a.s. bounded, by hypothesis. Thus, by part (b) of Theorem \ref{thm.Y0asy}, the second part of (A) holds, and hence (B) implies (A). Hence (A) and (B) are equivalent, establishing part (b) of the theorem. 

It remains to show that (B) implies  $y$ is bounded. Let $m$ be any positive, non--decreasing continuous function with $m(t)\to \infty$ as $t\to\infty$ 
%Then $y(t)/m(t)\to 0$ as $t\to\infty$. Hence, for any sequence $t_n\uparrow \infty$, $y(t_n)/m(t_n)\to 0$ as $n\to\infty$. 
and $(t_n)$ any sequence such that $t_n\uparrow \infty$ as $n\to\infty$. 
Now, $Y_n:=Y(t_n)/m(t_n)$ is a sequence of Gaussian random variables with mean $y(t_n)/m(t_n)$, and $Y_n\to 0$ as $n\to\infty$ a.s. since $Y$ is a.s. bounded. But then $Y_n$ converge in distribution, so $y(t_n)/m(t_n)\to 0$ as $n\to\infty$.   

Suppose next, by way of contradiction, that $\limsup_{t\to\infty} |y(t)|=+\infty$. Then $y^\ast(t):=\max_{0\leq s\leq t} |y(s)|\to\infty$ as $t\to\infty$. Pick $m(t)=\sqrt{y^\ast(t)}$, so that $m(t)\to\infty$ as $t\to\infty$, and let $t_n$ be a sequence along which $y^\ast(t) = |y(t)|$. Then 
$|y(t_n)|/m(t_n) = y^\ast(t_n)/\sqrt{y^\ast(t_n)} = \sqrt{y^\ast(t_n)}\to \infty$ as $n\to\infty$, 
which contradicts what was shown at the end of the last paragraph. Hence the supposition $\limsup_{t\to\infty} |y(t)|=+\infty$ is false, and $y$ is bounded, as required. 
\end{proof}
%
%\begin{theorem} \label{thm.Yto0bddchar}
%Let $f,\sigma$ obey \eqref{eq.fcns} and \eqref{eq.sigma}. Let $Y$ be the unique continuous adapted solution to \eqref{eq.Y}, and suppose $S$ is given by \eqref{eq.S}. 
%\begin{itemize}
%	\item[(a)] \hspace{0.1in} The following are equivalent:
%	\begin{enumerate}
%		\item[(A)] $\int_t^{t+\theta} f(s)\,ds \to 0$ as $t\to\infty$ for all $\theta\in [0,1]$, $S(\epsilon)<+\infty$ for all $\epsilon>0$;
%		\item[(B)] $Y(t)\to 0$ as $t\to\infty$ a.s. 
%	\end{enumerate}
%	\item[(b)] \hspace{0.1in} The following are equivalent:
%	\begin{enumerate}
%		\item[(A)] There exists $\bar{F}>0$ and $\epsilon'>0$ such that 
%		$\left|\int_t^{t+\theta} f(s)\,ds\right|\leq \bar{F}$ for all $\theta\in [0,1]$ and $t\geq 0$, and $S(\epsilon')<+\infty$ for some $\epsilon'>0$.
%		\item[(B)] $Y$ is a.s. bounded. 
%	\end{enumerate}
%	\item[(c)] \hspace{0.1in} The following are equivalent:
%	\begin{enumerate}
%		\item[(A)] The condition (A) in (b) holds, as well as at least one of the following:
%		\begin{gather*}
	%			\text{There exists $\theta'>0$ such that $\limsup_{t\to\infty} \left|\int_{t-\theta'}^t f(s)\,ds\right| >0$}; \\
	%			S(\epsilon'')=+\infty \text{ for some $\epsilon''>0$}. 
	%		\end{gather*}
%		\item[(B)] $\limsup_{t\to\infty} |Y(t)|\in (0,\infty)$ a.s.
%	\end{enumerate}
%\end{itemize} 
%\end{theorem}
%\begin{proof}
%Note $Y(t)=Y_0(t)+y(t)$ for $t\geq 0$, where $y$ is the solution of \eqref{eq.y}. 
%
%To prove part (a), consider first the implication (A) implies (B). If (A) holds, the first part of (A) gives $y(t)\to 0$ as $t\to\infty$, by the implication in Theorem~\ref{thm.detasymptoticstab} that (B) implies (A). The second part of (A) gives $Y_0(t)\to 0$ as $t\to\infty$ a.s., by part (a) of Theorem~\ref{thm.Y0asy}. Hence $Y(t)\to 0$ as $t\to\infty$.
%
%To show that (B) implies (A), note that $Y(t)$ is Gaussian distributed for each $t$, with mean $y(t)$. Consider any sequence $t_n\uparrow \infty$. Then $Y_n:=Y(t_n)$ converges to zero a.s., by (B). But since $Y_n$ is Gaussian, it means that it also converges in distribution, which forces $\mathbb{E}[Y_n]\to 0$ as $n\to\infty$, or $y(t_n)\to 0$ as $n\to\infty$. Suppose now, by way of contradiction, that $\limsup_{t\to\infty} |y(t)|>0$. 
%Then there exists $c>0$ and a sequence $t_n\uparrow \infty$ such that $|y(t_n)|>c$ for all $n$ sufficiently large. But this contradicts the fact that (B) implies $y(t_n)\to 0$ as $n\to\infty$ for any divergent sequence $(t_n)$. Hence, we must have $\limsup_{t\to\infty} |y(t)|=0$, or $y(t)\to 0$ as $t\to\infty$. But by the implication in Theorem~\ref{thm.detasymptoticstab} that (A) implies (B), we must have that $\int_t^{t+\theta} f(s)\,ds\to 0$ as $t\to\infty$, which proves the first part of (A). To prove the second part, we note we have just shown $y(t)\to 0$ as $t\to\infty$, and assumed that $Y(t)\to 0$ as $t\to\infty$ a.s. Hence $Y_0(t)=Y(t)-y(t)\to 0$ as $t\to\infty$ a.s. Now apply part (a) of Theorem \ref{thm.Y0asy} to see that  we must have $S(\epsilon)<+\infty$ for all $\epsilon>0$, which is the second part of (A). Thus we have shown that (B) implies (A), establishing the required equivalence to complete the proof of part (a).  
%
%To prove part (b), we start by proving that (A) implies (B). Notice that the first part of (A) yields that $y$ is bounded, by Theorem \ref{thm.detasymptoticbound}. Likewise, the second part of (A) yields that $Y_0$ is a.s. bounded, by part (b) of Theorem \ref{thm.Y0asy}. Thus $Y=Y_0+y$ must be a.s. bounded, proving (B). 
%
%To show that (B) implies (A), it will be enough to show that $y$ is bounded, for if $y$ is bounded, then the first part of (A) holds, by the converse half of Theorem \ref{thm.detasymptoticbound}. On the other hand, if $y$ is bounded, then $Y_0=Y-y$ is a.s. bounded too, by hypothesis. Therefore, by part (b) of Theorem \ref{thm.Y0asy}, it follows that the second part of (A) holds, and hence (B) implies (A). Thus (A) and (B) are equivalent, and hence we have established part (b) of the theorem. 
%
%It remains to show that (B) implies that $y$ is bounded. Let $m$ be any positive, non--decreasing continuous function with $m(t)\to \infty$ as $t\to\infty$. 
%%Then $y(t)/m(t)\to 0$ as $t\to\infty$. Hence, for any sequence $t_n\uparrow \infty$, $y(t_n)/m(t_n)\to 0$ as $n\to\infty$. 
%Let $(t_n)$ be any sequence such that $t_n\uparrow \infty$ as $n\to\infty$. 
%Now, notice that $Y_n:=Y(t_n)/m(t_n)$ is a sequence of Gaussian random variables with mean $y(t_n)/m(t_n)$ and that $Y_n\to 0$ as $n\to\infty$ a.s., by virtue of the a.s. boundedness of $Y$. Hence, by the convergence in distribution of the $Y_n$ we have $y(t_n)/m(t_n)\to 0$ as $n\to\infty$.   
%
%Suppose next, by way of contradiction, that $\limsup_{t\to\infty} |y(t)|=+\infty$. Then $y^\ast(t):=\max_{0\leq s\leq t} |y(s)|\to\infty$ as $t\to\infty$. Pick $m(t)=\sqrt{y^\ast(t)}$, so that $m(t)\to\infty$ as $t\to\infty$, and let $t_n$ be a sequence along which $y^\ast(t) = |y(t)|$. Then 
%\[
%\frac{|y(t_n)|}{m(t_n)} = \frac{y^\ast(t_n)}{\sqrt{y^\ast(t_n)}} = \sqrt{y^\ast(t_n)}\to \infty, \quad n\to\infty.
%\] 
%But this contradicts what was shown at the end of the last paragraph. Hence the supposition $\limsup_{t\to\infty} |y(t)|=+\infty$ must be false, and so $y$ is bounded, as required. 
%\end{proof}
%
%\section{Quantitative estimates for the singly perturbed SDE}
%Most of the work in this section concerns the SDE \eqref{eq.Y0}. However, once a characterisation of its asymptotic behaviour is complete, we may use the line of proof in 
%Theorem~\ref {thm.Yto0bddchar} to immediately characterise growth or decay bounds. 

%As for the SDE \eqref{eq.Y0} itself, #
This result, together with the discussion before Theorem~\ref{thm.Y0asy},  gives a characterisation of stability and  boundedness for \eqref{eq.stochforced2}.
\begin{theorem}\label{thm.stochstabtilX1}
Suppose that $q\in (0,1)$ and \eqref{eq.unforcedasstable} holds. Let $f,\sigma\in C([0,\infty);\mathbb{R})$ and $X$ and $Y$ be the unique continuous adapted solutions to \eqref{eq.stochforced2} and \eqref{eq.Y}. Then
\begin{enumerate}
	\item[(i)] Condition (a)(A) of Theorem~\ref{thm.Yto0bddchar}
	$\Longleftrightarrow$ $X(t)\to 0$ as $t\to\infty$ a.s.;
	\item[(ii)] Condition (b)(A) of Theorem~\ref{thm.Yto0bddchar} $\Longleftrightarrow$  $X$ is bounded a.s. 
\end{enumerate}
\end{theorem}
Some remarks are in order; first, Theorem~\ref{thm.Yto0bddchar} is to the best of our knowledge (and remarkably, given the simplicity of the equation) the first pathwise characterisation of asymptotic stability and boundedness for a stochastic differential equation which is subjected to both deterministic and stochastic perturbations. \textit{We consider this an auxiliary result of tremendous importance}. We see straightaway that it can be used to characterise the a.s. convergence and boundedness of the stochastically and deterministically forced pantograph equation, an equation far more complicated that the SDE, revealing that \textit{the stability conditions are independent of the structure of the underlying differential system}. Again, such a stability characterisation for deterministically and stochastically perturbed functional differential equations is, to the best of our knowledge, the first to appear in the literature.  

To see the potential for Theorem~\ref{thm.Yto0bddchar} to have very wide applicability to characterise many other complicated differential systems which are forced state--independently, we give a sketch of the possible direction of research. Suppose, concretely, we have the following stochastic functional differential equation
\[
dX(t)=(F(t,X_t)+f(t))\,dt + \sigma(t)\,dB(t), \quad t\geq 0,
\]    
where $F$ is a functional on $[0,\infty)\times C([0,\infty))$. We view $X$ as a perturbation of the underlying deterministic system $u'(t)=F(t,u_t)$. We use the common notational convention for functional differential equations that for a continuous function $x:\mathbb{R}^+\to \mathbb{R}$, $x_t$ refers to the history of $x$ i.e., to $\{x(s):0\leq s\leq t\}$. Now, consider $Y$ as above, and let $Z:=X-Y$. Then 
\[
Z'(t)=F(t,Z_t+Y_t)+Y(t), \quad t\geq 0.
\]  
Let $F$ satisfy a uniform functional Lipschitz condition in the state i.e., if   for all $\phi$, $\psi\in C(\mathbb{R}^+;\mathbb{R})$, there exists a constant $K>0$ such that 
\[
|F(t,\phi_t)-F(t,\psi_t)|\leq K\max_{0\leq s\leq t}|\phi(s)-\psi(s)|, \quad t\geq 0.
\]
Note that this condition is sufficient to guarantee a unique continuous adapted solution of the original SFDE on $[0,\infty)$ (see e.g., Mao~\cite{MaoBK}). 
Thus, writing $\tilde{f}(t):=F(t,Z_t+Y_t)-F(t,Z_t)+Y(t)$ for $t\geq 0$, we see that $\tilde{f}(t)\to 0$ as $t\to\infty$, $\tilde{f}$ is continuous on $[0,\infty)$, and $Z$ obeys 
\[
Z'(t)=F(t,Z_t)+\tilde{f}(t), \quad t\geq 0.
\] 
Now, if a good perturbation theory exists for $C_0(\mathbb{R}^+;\mathbb{R})$ perturbations of the equation $u'(t)=F(t,u_t)$, we may be able to show that $Z(t)\to 0$ as $t\to\infty$, and hence that $X(t)=Y(t)+Z(t)\to 0$ as $t\to\infty$. 

In the converse direction, stability of the unperturbed functional differential equation means that we should impose the condition 
\[
F(t,\phi_t)\to 0 \quad \text{as $t\to\infty$ for each $\phi\in C_0(\mathbb{R}^+;\mathbb{R})$}.
\]
By our usual arguments, $Y$ can be connected to $X$ via
\[
Y(t)=X(t)-X(0)e^{-t}-\int_{0}^t e^{-(t-s)}\{X(s)+F(s,X_s)\}\,ds, \quad t\geq 0.
\]
Therefore, if $X(t)\to 0$ as $t\to\infty$, the term in curly brackets in the integrand is a continuous function in $s$, and tends to $0$ as $s\to\infty$. Therefore, the integral tends to zero as $t\to\infty$, and hence $Y(t)\to 0$ as $t\to\infty$. From this brief discussion, it can be seen how an effective convergence theory for many ``doubly perturbed'' functional differential equations can be developed. In rough terms, provided the asymptotic stability of the unperturbed is robust to deterministic $C_0$ perturbations, then the convergence of the stochastic FDE system is equivalent to the convergence of the SDE, which has already been characterised in Theorem \ref{thm.Yto0bddchar}. 
%Indeed, since the conditions (A) in Theorem~\ref{thm.Yto0bddchar} are exhaustive, we see that solutions of \eqref{eq.stochforced2} either (a) tend to zero with probability one; (b) are bounded, but not convergent with probability one; or (c) are unbounded with probability one. 

%For growth or decay bounds, define for a continuous positive function $\gamma$ the sum  
%\begin{equation} \label{eq.Stilde}
%\tilde{S}_\gamma(\epsilon) = \sum_{n=1}^\infty \frac{\gamma(n-1)}{\left(\int_{n-1}^{n}\sigma^2(s)ds\right)^{1/2}}\exp\left(-\epsilon \frac{\gamma^2(n-1)}{\int_{n-1}^{n}\sigma^2(s)ds}\right),
%\end{equation}
%(we take the summand to be zero if $\int_{n-1}^{n}\sigma^2(s)ds$), and consider the condition 
%\begin{equation} \label{eq.gammauniformquotientcompact}
%\liminf_{t\to\infty}
%\inf_{0\leq \theta\leq 1}
%\frac{\gamma(t-\theta)}{\gamma(t)} >0, \quad 
%\limsup_{t\to\infty} \sup_{0\leq \theta\leq 1}
%\frac{\gamma(t-\theta)}{\gamma(t)} <+\infty,
%\end{equation}
%which is satisfied if $\gamma$ is subexponential, for instance. 
%, it turns out if we let $\Psi$ be the complementary normal distribution function and define 
%\begin{equation} \label{def.Sf}
%	S(\epsilon):= \sum_{n=1}^\infty  \Psi \left(\frac{\epsilon}{\theta(n)} \right).
%	%\quad \theta^2(n) := \frac{ \int_{(n-1)}^{n}{\sigma^2(s)ds} }{ \psi^2(n-1) }, \quad n\geq 1,
%\end{equation} 
%then 
%	\begin{align} \label{eq.Y0psito0}
%	\lim_{t \to \infty}\frac{Y_0(t)}{\psi(t)} = 0, \quad\text{a.s.}
%\end{align} 
%\[
%\lim_{t\to\infty} \frac{Y_0(t)}{\psi(t)}=0
%\]
%holds if and only if we have $S(\epsilon)<+\infty$ for all $\epsilon>0$. 
%The next lemma shows that \eqref{def.Sf} implies \eqref{eq.Y0psito0}, while a subsequent lemma gives the reverse implication. 
%
%We prove this result in a sequence of lemmata.
%\begin{lemma} \label{sigmaIntegral:zero}
%Let $\sigma$ be continuous. 
%	Let $h>0$ and define $S(\epsilon,h)$ as in \eqref{def.Sf}. If $S(\epsilon,h)< +\infty$ for all $\epsilon > 0$, then 
%	\begin{align*}
	%		\lim_{t \to \infty}\frac{ \int_{t}^{t+h}{\sigma(s)dB(s)} }{ \psi(t) }=0, \text{ a.s.}
	%	\end{align*}
%\end{lemma}
%\begin{proof}
%	Considering $\int_t^{t+h}{\sigma(s)dB(s)}$ we write, for $nh \leq t \leq (n+1)h$,
%	\begin{align*}
	%		\left|\int_t^{t+h}{\sigma(s)dB(s)}\right| \leq \left|-\int_{(n+1)h}^{t}{\sigma(s)dB(s)}\right| + \left|\int_{(n+1)h}^{t+h}{\sigma(s)dB(s)}\right|.
	%	\end{align*} 
%	It follows that 
%	\begin{align*}
	%		\sup_{nh \leq t \leq (n+1)h}\left|\int_t^{t+h}{\sigma(s)dB(s)}\right| \leq &\sup_{nh \leq t \leq (n+1)h} \left|-\int_{(n+1)h}^{t}{\sigma(s)dB(s)}\right| \\ + &\sup_{nh \leq t \leq (n+1)h} \left|\int_{(n+1)h}^{t+h}{\sigma(s)dB(s)}\right|.
	%	\end{align*} 
%	Similarly we obtain 
%	\begin{align} \label{eq:***}
	%		\nonumber \sup_{nh \leq t \leq (n+1)h}\left|\int_t^{t+h}{\sigma(s)dB(s)}\right| \leq &\left|-\int_{nh}^{(n+1)h}{\sigma(s)dB(s)}\right| \\ + &\sup_{(n+1)h \leq t \leq (n+2)h} \left|\int_{(n+1)h}^{t}{\sigma(s)dB(s)}\right|.
	%	\end{align}
%	Here we note that $S(\epsilon,h)< +\infty$ implies that 
%	\begin{align*}
	%		\lim_{t \to \infty}\frac{\int_{nh}^{(n+1)h}{\sigma(s)dB(s)}}{\psi(nh)} = 0, \quad \text{a.s.}
	%	\end{align*}
%	by means of the first Borel--Cantelli lemma. Next we proceed to estimate 
%	\begin{align*}
	%		\mathbb{P}\left[ Z((n+1)h) > \epsilon \right],\,\, \text{for some } \epsilon \in (0,1), 
	%	\end{align*}
%	where
%	\begin{align*} 
	%		Z((n+1)h) := \sup_{(n+1)h \leq t \leq (n+2)h} \frac{\left|\int_{(n+1)h}^{t}{\sigma(s)dB(s)}\right|}{\psi((n+1)h)}, \,\, n \geq 1.
	%	\end{align*}
%	We also define the function
%	\begin{align*}
	%		\tau(t) := \frac{ \int_{(n+1)h}^{t}\sigma^2(s)ds }{ \psi^2((n+1)h) }, \text{ for } t \in [(n+1)h, (n+2)h].
	%	\end{align*}
%	By the Martingale Time Change Theorem there exists a standard Brownian motion $B_n^*$ such that
%	\begin{align*}
	%		\lefteqn{
		%			\mathbb{P}[ Z((n+1)h) > \epsilon]} \\
	%		&=\mathbb{P}\left[ \sup_{t \in [(n+1)h, (n+2)h]}\left| B_{n+1}^*(\tau(t)) \right| > \epsilon \right] \\
	%		&=\mathbb{P}\left[\sup_{u \in [0, \tau((n+2)h)]} |B_{n+1}^*(u)| > \epsilon \right] \\
	%		&\leq \mathbb{P}\left[\sup_{u \in [0, \tau((n+2)h)]} B_{n+1}^*(u) > \epsilon \right] 
	%		+ \mathbb{P}\left[\sup_{u \in [0, \tau((n+2)h)]} -B_{n+1}^*(u) > \epsilon \right] \\ 
	%		&=\mathbb{P}\left[ |B_{n+1}^*(\tau((n+2)h))| > \epsilon \right] + \mathbb{P}\left[ |B_{n+1}^{**}(\tau((n+2)h))| > \epsilon \right],
	%	\end{align*}
%	where $-B_{n+1}^* = B_{n+1}^{**}$ is a standard Brownian motion. Thus, as $B_{n+1}^*(\tau((n+2)h))$ is normally distributed with zero mean we have 
%	\begin{align*}
	%		\lefteqn{
		%			\mathbb{P}[Z((n+1)h) > \epsilon]}\\ 
	%		&\leq 2\mathbb{P}\left[ |B_{n+1}^*(\tau((n+2)h))| > \epsilon \right] 
	%		= 4\mathbb{P}\left[ B_{n+1}^*(\tau((n+2)h)) > \epsilon \right] \\
	%		%&
	%		&= 4\Psi\left(\frac{\epsilon}{\sqrt{\tau((n+2)h)}} \right)  = 4\Psi\left(\frac{\epsilon}{\theta(n+1)} \right).
	%	\end{align*}
%	But since we assumed that $S(\epsilon) < +\infty$, we have
%	\begin{align*}
	%		\sum_{n=0}^\infty \mathbb{P}[Z((n+1)h) > \epsilon]  \leq 4\sum_{n=0}^\infty \Psi\left( \frac{\epsilon}{\sqrt{\tau((n+2)h)}} \right) 
	%		< +\infty.
	%	\end{align*}
%	We can then apply the first Borel--Cantelli Lemma to conclude that 
%	\begin{align*}
	%		\limsup_{n \to \infty}Z((n+1)h)< \epsilon \quad\text{a.s.}
	%	\end{align*}
%	Thus
%	\begin{align*}
	%		\lim_{n \to \infty} \sup_{(n+1)h \leq t \leq (n+2)h} 
	%		\frac{\left| \int_{(n+1)h}^{t} \sigma(s)\,dB(s) \right|}{ \psi((n+1)h) } = 0 \quad\text{a.s.}
	%	\end{align*}
%	Combining this with (\ref{eq:***}) we get 
%	\begin{align*}
	%		\lim_{n \to \infty} \sup_{nh \leq t \leq (n+1)h} \frac{\left| \int_{t}^{t+h} \sigma(s)\,dB(s) \right|}{ \psi(nh) } = 0,\quad \text{a.s.},
	%	\end{align*}
%	as required.
%\end{proof}
%We next need the asymptotic behaviour of an auxiliary process which solves an affine SDE. 
%\begin{lemma} \label{eq.YOUasy}
%Let $\sigma$ be continuous, $\psi$ be subexponential.
%	Let $h>0$ and define $S(\epsilon)$ as in \eqref{def.Sf}. Suppose that $S(\epsilon)< +\infty$ for all $\epsilon > 0$.
%	Let $Y_0$ be the unique continuous adapted process which solves \eqref{eq.Y0}.   
%%	\begin{align} \label{Y}
	%%		dY(t) = -Y(t)dt + \sigma(t)dB(t), \, t \geq 0, \quad Y(0)=0.
	%%	\end{align}
%	Then \eqref{eq.Y0psito0} holds. 
%\end{lemma}
%\begin{proof}
%	Define 
%	\[
%	V(n) = \int_{n-1}^{n} e^{s-n}\sigma(s)\,dB(s), \quad n\geq 1; \quad \tilde{V}(n)=\frac{V(n)}{\psi(n-1)}, \quad n\geq 1.
%	\]
%	Then $(\tilde{V}(n))_{n\geq 1}$ is a sequence of independent normal random variables with zero mean and variance
%	\[
%	\tilde{v}^2(n)=\frac{1}{\psi^2(n-1)}\int_{n-1}^{n} e^{2s-2n}\sigma^2(s)\,ds, \quad n\geq 1.
%	\]
%	We show first that $\tilde{V}_h(n)\to 0$ a.s. as $n\to\infty$. 
%	%By the fact that $f\circ F^{-1}$ is asymptotic to $\varphi\circ \Phi^{-1}$, 
%%	There is $N=N(h)$ such that for all $n\geq N(h)$ we have 
%Note that
%	\[
%	\tilde{v}^2(n)\leq \frac{1}{\psi(n-1)^2} \int_{(n-1)}^{n} \sigma^2(s)\,ds \leq  \theta^2(n).
%	\]
%	Hence $\tilde{v}(n)\leq \theta(n)$ for $n\geq 1$. Also
%	$\mathbb{P}[|\tilde{V}(n)|>\epsilon]=2\mathbb{P}[\tilde{V}(n)/\tilde{v}(n)>\epsilon/\tilde{v}(n)]$, so
%	\[
%	\mathbb{P}[|\tilde{V}(n)|>\epsilon]
%	=2\Psi(\epsilon/\tilde{v}(n)) 
%	%\leq 2\Psi(\epsilon/\tilde{v}_h(n))
%	\leq 2\Psi(\epsilon/\theta(n)), \quad n\geq 1,
%	\]
%	since $\Psi$ is decreasing, and $\epsilon/v(n)\geq \epsilon/\theta(n)$ for $n\geq 1$. Now, as $S(\epsilon)<+\infty$, it follows that 
%	\[
%	\sum_{n=1}^\infty \mathbb{P}[|\tilde{V}(n)|>\epsilon] <+\infty
%	\] 
%	for every $\epsilon>0$, and hence, by the first Borel--Cantelli lemma, it follows that $\mathbb{P}[\lim_{n\to\infty} \tilde{V}(n)=0]=1$. 
%	
%	Next, $Y_0$ obeys  
%	\begin{align*}
	%		Y(t) = e^{-t}\int_0^t e^s \sigma(s)\, dB(s),\quad t\geq 0.
	%	\end{align*}
%	Notice that  
%	$Y_0(n+1) = e^{-1}Y_0(n) + V(n+1)$ for $n\geq 0$. Now we define $\tilde{Y}(t) := Y_0(t)/\psi(t)$ for $t\geq 0$ and thus we have 
%%	\begin{align*}
	%%		\tilde{Y}((n+1)h) &= e^{-h}\frac{Y(nh)}{\psi((n+1)h)} + %\frac{V_h(n+1)}{\psi((n+1)h)} \\
	%%		&= e^{-h}\frac{\psi(nh)}{\psi((n+1)h)}\tilde{Y}(nh) + %\frac{\psi(nh)}{\psi((n+1)h)}\tilde{V}_h(n+1).
	%%	\end{align*}
%%	Hence
%	\begin{equation} \label{eq.tildeYn}
	%		\tilde{Y}(n+1) = a(n) \tilde{Y}(n) + a(n)e\tilde{V}(n+1), \quad n\geq 0,
	%	\end{equation}
%	where $a(n) := e^{-1}\psi(n) / \psi(n+1)$. We note that $a(n) > 0$ for all $n \in \mathbb{N}$ and that $\lim_{n \to \infty}a(n)=e^{-1}$ by the fact that $\psi$ is subexponential. Notice that $(1+1/2)e^{-1}<1$. Since $a(n)\to e^{-1}$ as $n\to\infty$, there exists $N_1\in \mathbb{N}$ such that $a(n)\leq (1+1/2)e^{-1}<1$ for all $n\geq N_1$. Next, we may write, for all $n > N_1$
%	\begin{align*}
	%		|\tilde{Y}(n+1)| &\leq a(n) |\tilde{Y}(n)| + a(n)e|\tilde{V}(n+1)| \\ 
	%		&\leq e^{-1}(1+1/2)|\tilde{Y}(n)| + (1+1/2)|\tilde{V}(n+1)|.
	%	\end{align*}
%	From this inequality we define
%	\begin{gather*}
	%		\bar{Y}((n+1)) = e^{-1}(1+1/2)\bar{Y}(n)+ (1+1/2)|\tilde{V}(n+1)|, \quad n \geq  N_1+1,\\ 
	%		\bar{Y}(n) = |\tilde{Y}(n)| +1, \quad n = N_1+1.
	%	\end{gather*}
%	Therefore we have that $|\tilde{Y}(n)|<\bar{Y}(n)$ for $n\geq N_1+1$. 
%	Since $\tilde{V}(n)\to 0$ as $n\to\infty$ a.s., it follows that $\bar{Y}(n)\to 0$ as $n\to\infty$ a.s. Hence $\tilde{Y}(n)\to 0$ as $n\to\infty$ a.s.
%	
%	Next, let $t \in [n, n+1]$. Then
%	\begin{align*}
	%		\frac{Y_0(t)}{\psi(t)} = \frac{1}{\psi(t)} Y_0(n)e^{-(t-n)} 
	%		+ \frac{e^{-t}}{\psi(t)}\int_{n}^{t} e^s\sigma(s)\,dB(s).
	%	\end{align*}
%	Now, by the uniform convergence property of subexponential functions, there is $N_2\in\mathbb{N}$ such that $\psi(t)>\psi(n)/2$ %, $\psi(t)>\psi((n+1)h)/2$ 
%	for all $t\in [n,n+1]$ and all $n\geq N_2$.
%	%Notice since $\varphi'(x)\to 0$ as $x\to 0$, and $(\Phi^{-1})'(t)=(\varphi\circ \Phi^{-1})(t)$ that 
%	%$t\mapsto e^{-t}/(\varphi\circ \Phi^{-1})(t)$ is decreasing on $(T_3,\infty)$ for some $T_3>0$. 
%	Define $T_2$ such that $N_2>T_2$. 
%	%Also $t\mapsto (\varphi\circ \Phi^{-1})(t)$ is decreasing on $[0,\infty)$. 
%	Then we have for $n\geq N_2$ that $t\geq n\geq N_2>T_2$, and so 
%	\begin{align*}
	%		\lefteqn{
		%			\sup_{t\in [n,n+1]}
		%			\frac{|Y_0(t)|}{\psi(t)}}
	%		\\ 
	%		&\leq  
	%		\sup_{t\in [n,n+1]}
	%		\frac{1}{\psi(t)} |Y_0(n)|  
	%		+ \sup_{t\in [n,n+1]}\frac{e^{-t}}{\psi(t)}\left|\int_{n}^{t} e^s\sigma(s)\,dB(s)\right|\\
	%		&\leq 2\frac{|Y_0(n)|}{\psi(n)}  
	%		+ 2e\frac{e^{-(n+1)}}{\psi(n)}\sup_{t\in [n,n+1]}\left|\int_{n}^{t} e^s\sigma(s)\,dB(s)\right|.
	%	\end{align*}
%	Since $Y_0(n)/\psi(n)\to 0$ as $n\to\infty$ a.s. we have that 
%	the first term on the right--hand side has zero limit as $n\to\infty$ a.s. Therefore it remains to prove that 
%	\[
%	U(n+1):=\frac{e^{-(n+1)}}{\psi(n)}\sup_{t\in [n,n+1]}\left|\int_{n}^{t} e^s\sigma(s)\,dB(s)\right|
%	\]
%	obeys $U(n)\to 0$ as $n\to\infty$ a.s., as this will demonstrate that 
%	\[
%	\lim_{n\to\infty} \sup_{t\in [n,n+1]} \frac{|Y_0(t)|}{\psi(t)}=0, \quad\text{a.s.}
%	\]
%	Next, we see that with $\rho(t)=\int_{n}^{t} e^{2s}\sigma^2(s)\,ds$ for $t\in [n,n+1]$, by the martingale time change theorem, 
%	there exists a standard Brownian motion $B^\ast_n$ such that
%	\begin{align*}
	%		\lefteqn{\mathbb{P}[U(n+1)>\epsilon]}\\
	%		&=\mathbb{P}\left[\sup_{t\in [n,n+1]}\left|\int_{n}^{t} e^s\sigma(s)\,dB(s)\right|>\epsilon \frac{e^{(n+1)}}{\psi(n)} \right]\\
	%		&=\mathbb{P}\left[\sup_{t\in [n,n+1]}|B^\ast_n(\rho(t))|>\epsilon \frac{e^{(n+1)}}{\psi(n)} \right]
	%		=\mathbb{P}\left[\sup_{t\in [0,\rho(n+1)]}|B^\ast_n(t)|>\epsilon \frac{e^{(n+1)}}{\psi(n)} \right].
	%	\end{align*}
%	By standard arguments, we get that 
%	\begin{multline*}
	%		\mathbb{P}[U(n+1)>\epsilon]
	%		\leq 2\mathbb{P}\left[ \sup_{t\in [0,\rho(n+1)]} B^\ast_n(t) >\epsilon \frac{e^{(n+1)}}{\psi(n)}  \right]\\
	%		= 2\mathbb{P}\left[|B^\ast_n(\rho(n+1))| >\epsilon \frac{e^{(n+1)}}{\psi(n)}\right]
	%		=4\Psi\left(\epsilon \frac{e^{(n+1)}}{\psi(n)} \frac{1}{\sqrt{\rho(n+1)}}  \right).
	%	\end{multline*}
%	Finally, we estimate the right--hand side of the above expression. Since 
%	\[
%	\sqrt{\rho(n+1)}=\left(\int_{n}^{n+1} e^{2s}\sigma^2(s)\,ds\right)^{1/2}\leq e^{(n+1)h}\left(\int_{n}^{n+1} \sigma^2(s)\,ds\right)^{1/2},
%	\]
%	we have
%	\[
%		\epsilon\frac{e^{n+1}}{\psi(n)}\cdot 
%		\frac{1}{\sqrt{\rho(n+1)}}
%		\geq 
%		\epsilon \frac{1}{\psi(n)} \left(\int_{h}^{n+1} \sigma^2(s)\,ds\right)^{-1/2} = \frac{\epsilon}{\theta(n)}.
%	\]	
%	%Next, the fact that $f\circ F^{-1}$ is asymptotic to $\varphi\circ\Phi^{-1}$ and that both are regularly varying functions 
%	%means there is an $N_4(h)\in\mathbb{N}$ such that
%	%\[
%	%\frac{(f\circ F^{-1})(nh)}{(\varphi \circ \Phi^{-1})((n+1)h)}\geq \frac{1}{2}, \quad %n\geq N_4(h). 
%	%\]
%	%Hence by the definition of $\theta(n)$, for $n\geq 1$ that N_4(h)$, we get
%%	\[
%%	\epsilon\frac{e^{(n+1)h}}{(\varphi \circ \Phi^{-1})((n+1)h)}\cdot 
%%	\frac{1}{\sqrt{\rho((n+1)h)}}
%%	\geq 
%%	\frac{\epsilon}{2} \frac{1}{\theta(n)},
%%	\]
%	so as $\Psi$ is decreasing, we have for $n\geq 1$ that
%	\begin{align*}
	%		\mathbb{P}[U(n+1)>\epsilon]
	%		\leq 4\Psi\left(\epsilon \frac{e^{(n+1)}}{\psi(n)} \frac{1}{\sqrt{\rho(n+1)}}  \right)
	%		\leq 4\Psi\left(\frac{\epsilon}{\theta(n)}\right).
	%	\end{align*}
%	Since $S(\epsilon)<+\infty$ for all $\epsilon>0$, it follows that 
%	$\sum_{n=1}^\infty  \mathbb{P}[U(n+1)>\epsilon] < +\infty$ 	for every $\epsilon>0$, and therefore, by the first Borel--Cantelli lemma, it follows that $\mathbb{P}[\lim_{n\to\infty} U(n)=0]=1$.
%	As noted above, this is the remaining fact that guarantees that $\lim_{t\to\infty} Y_0(t)/\psi(t)=0$ a.s., as required.
%\end{proof}  
%
%\begin{lemma} \label{eq.YOUasyconv}
%	Let $\sigma$ be continuous, $\psi$ be subexponential. 
%	Define $S(\epsilon)$ as in \eqref{def.Sf}. %Suppose that $S(\epsilon,h)< +\infty$ for all $\epsilon > 0$.
%	Let $Y_0$ be the unique continuous adapted process which solves \eqref{eq.Y0} and which obeys \eqref{eq.Y0psito0}. Then  $S(\epsilon)< +\infty$ for all $\epsilon > 0$.
%\end{lemma}
%\begin{proof}
%Revisit \eqref{eq.tildeYn} where, as before, $a(n) := e^{-1}\psi(n) / \psi(n+1)$, so $a(n)\to e^{-1}$ as $n\to\infty$. By hypothesis, $\tilde{Y}(n)\to 0$ as $n\to\infty$ a.s., so therefore we have $a(n)e\tilde{V}(n+1)\to 0$ as $n\to\infty$ a.s. Since $a(n)\to e^{-1}$ as $n\to\infty$, we have  $\tilde{V}(n+1)\to 0$ as $n\to\infty$ a.s. Since $\tilde{V}(n+1)$ are independent normal random variables with zero mean and variance $\tilde{v}(n)^2$, by the second Borel--Cantelli lemma, this implies $\sum_{n=1}^\infty \mathbb{P}\left[|\tilde{V}(n)|>\epsilon\right]<+\infty$ 
%for every $\epsilon>0$. As before
%%$\mathbb{P}[|\tilde{V}_h(n)|>\epsilon]
%%=2\mathbb{P}[\tilde{V}_h(n)/\tilde{v}_h(n)>\epsilon/\tilde{v}_h(n)]$%, so
%$
%\mathbb{P}[|\tilde{V}(n)|>\epsilon]
%=2\Psi(\epsilon/\tilde{v}(n))$. 
%%\leq 2\Psi(\epsilon/\tilde{v}_h(n))
%%\leq 2\Psi(\epsilon/\theta(n)), \quad n\geq 1,
%%\]
%Next, note that 
%	\[
%\tilde{v}^2(n)%=\frac{1}{\psi^2((n-1)h)}\int_{(n-1)h}^{nh} e^{2s-2nh}\sigma^2(s)\,ds
%\geq e^{-2}\frac{1}{\psi^2(n-1)}\int_{n-1}^{n}
%\sigma^2(s)\,ds = e^{-2}\theta^2(n),
%\]
%so $\tilde{v}(n)\geq e^{-1}\theta(n)$. Thus $\epsilon/\tilde{v}(n)\leq \epsilon/( e^{-1}\theta(n))$ and so 
%$\Psi(\epsilon/\tilde{v}(n))\geq \Psi(\epsilon/( e^{-1}\theta(n)))$. Therefore 
%\[
%\mathbb{P}[|\tilde{V}(n)|>\epsilon]
%=2\Psi(\epsilon/\tilde{v}(n))\geq 2\Psi(\epsilon/( e^{-1}\theta(n))). 
%\]
%Since the lefthand side is summable, so is the right; but this means that $S(\epsilon e)<+\infty$ for all $\epsilon>0$, as needed. 
%\end{proof}
%
%Combining these lemmas, we have the following result.
%\begin{theorem} \label{thm.Y0opsichar}
%	Let $\sigma$ be continuous, $\gamma$ be subexponential or obeys 
%	\begin{equation} \label{eq.gammauniformquotientcompact}
	%		\liminf_{t\to\infty}
	%		\inf_{0\leq \theta\leq 1}
	%		\frac{\gamma(t-\theta)}{\gamma(t)} >0, \quad 
	%		\limsup_{t\to\infty} \sup_{0\leq \theta\leq 1}
	%		\frac{\gamma(t-\theta)}{\gamma(t)} <+\infty.
	%	\end{equation}
%Define $\tilde{S}(\epsilon)$ as in \eqref{eq.Stilde}.  
%%Suppose that $S(\epsilon,h)< +\infty$ for all $\epsilon > 0$.
%Let $Y_0$ be the unique continuous adapted process which solves \eqref{eq.Y0}. 
%\begin{itemize}
%	\item[(i)] \hspace{0.1in} The following are equivalent:
%\begin{enumerate} 
%	\item[(A)] $\tilde{S}(\epsilon)<+\infty$ for all $\epsilon>0$;
%	\item[(B)] $Y_0(t)=\text{o}(\gamma(t))$, as $t\to\infty$ a.s.  
%\end{enumerate}
%\item[(ii)]  \hspace{0.1in}  The following are equivalent:	
%\begin{enumerate} 
%	\item[(A)]  $\tilde{S}(\epsilon)<+\infty$ for some $\epsilon>0$;
%	\item[(B)] $Y_0(t)=\text{O}(\gamma(t))$ as $t\to\infty$ a.s.
%\end{enumerate}
%\item[(iii)]  \hspace{0.1in} If $\tilde{S}(\epsilon)=+\infty$ for some $\epsilon>0$, then 
%	$\limsup_{t\to \infty} |Y_0(t)|/\gamma(t)>0$ a.s.  
%\end{itemize}
%\end{theorem}
%
%Moreover, scrutiny of the proofs reveals that boundedness can be characerised.

%\begin{theorem}  \label{thm.Y0bddpsichar}
%	Let $\sigma$ be continuous, $\psi$ be subexponential. 
%Define $S(\epsilon)$ as in \eqref{def.Sf}. 
%	%Suppose that $S(\epsilon,h)< +\infty$ for all $\epsilon > 0$.
%	Let $Y_0$ be the unique continuous adapted process which solves \eqref{eq.Y0}. Then the following are equivalent:
%	\begin{enumerate} 
	%		\item[(A)] 
	%		\[
	%		S(\epsilon) =\sum_{n=1}^\infty  \Psi \left(\frac{\epsilon \psi(n-1)}{\left(\int_{n-1}^{n}\sigma^2(s)ds\right)^{1/2}} \right)
	%		< +\infty \text{ for some $\epsilon > 0$};
	%		\]
	%		\item[(B)] 
	%		\[
	%		\limsup_{t\to \infty} \frac{|Y_0(t)|}{\psi(t)}<+\infty, \quad \text{a.s.} 
	%		\]
	%	\end{enumerate}
%	Moreover, if $S(\epsilon)=+\infty$ for some $\epsilon=\epsilon'>0$, we have 
%	\[
%	\limsup_{t\to \infty} \frac{|Y_0(t)|}{\psi(t)}>0, \quad \text{a.s.} 
%	\]
%\end{theorem}
%\begin{proof} 
%	As indicated, the proof of the equivalence follows the pattern of Theorem~\ref{thm.Y0opsichar}. We prove only the last claim, which rules out the possibility that $Y(t)=o(\psi(t))$ as $t\to\infty$ on an event of positive probability. 
%	
%	To start, note that 
%	\[
%	U(n) = \frac{1}{\psi(n-1)}\int_{n-1}^n \sigma(s)\,dB(s), \quad n\geq 1.
%	\]  
%	is a sequence of independent zero mean normal random variables with variance 
%	\[
%	w(n)^2:= \frac{1}{\psi^2(n-1)}\int_{n-1}^n \sigma^2(s)\,ds.	
%	\]
%	Therefore $\mathbb{P}[|U(n)|>\epsilon']/2=\Psi(\epsilon'/w(n))$. These are the summands in $S(\epsilon')$, so by hypothesis, we have 
%	$\sum_{n=1}^\infty \mathbb{P}[|U(n)|>\epsilon'] = +\infty$. By the second Borel--Cantelli lemma, this means that $
%	\mathbb{P}\left[|U(n)|>\epsilon' \text{ i.o.}\right]=1$,
%	or to put it another way $\limsup_{n\to\infty} |U(n)|\geq \epsilon'$ a.s. Next, integrate \eqref{eq.Y0} over $[n-1,n]$, divide by $\psi(n-1)$ and rearrange to get  
%	\[
%	U(n)=
%	\frac{\psi(n)}{\psi(n-1)}\frac{Y_0(n)}{\psi(n)} - \frac{Y_0(n-1)}{\psi(n-1)} + \frac{1}{\psi(n-1)}\int_{n-1}^n \psi(s) \frac{Y_0(s)}{\psi(s)}\,ds.
%	\]
%	Suppose there is an event of positive probability ($C$, say) on which $Y_0(t)/\psi(t)\to 0$ as $t\to\infty$. By the uniform convergence property of $\psi$, $\int_{n-1}^n \psi(s)\,ds/\psi(n)\to 1$ as $n\to\infty$, so on $C$, the last term on the right tends to zero, as do the first two terms on the right (invoking the fact that $\psi$ is subexponential once again). Hence for each $\omega\in C$, $U(n,\omega)\to 0$ as $n\to\infty$. But since $C$ is an event of positive probability, this contradicts the earlier conclusion that $\limsup_{n\to\infty} |U(n)|\geq \epsilon'>0$ a.s. Hence, the event $C$ must be a zero probability event, and its complement $\{\limsup_{t\to\infty} |Y_0(t)|/\psi(t)>0\}$ must be almost sure, as claimed.  
%	\end{proof}
%
%We close this section with some remarks. Note that the results deal with both decay and growth bounds, according the whether the subexponential function tends to zero or to infinity.  
%
%The complementary distribution function does not have a closed form expression in terms of elementary functions, making the conditions (A) in these Theorems hard to apply. However, note that 
%\[
%\Psi(x)\sim \frac{1}{\sqrt{2\pi}}\frac{1}{x}e^{-x^2/2}, \quad x\to\infty,
%\]
%so if follows in conditions (A) that we can instead consider the explicit sum 
%\begin{equation} \label{eq.Stilde}
%\tilde{S}(\epsilon) = \sum_{n=1}^\infty \frac{\psi(n-1)}{\left(\int_{n-1}^{n}\sigma^2(s)ds\right)^{1/2}}\exp\left(-\epsilon \frac{\psi^2(n-1)}{\int_{n-1}^{n}\sigma^2(s)ds}\right).
%\end{equation}
%To be definite: we have that $Y_0(t)/\psi(t)\to 0$ a.s. if and only if $\tilde{S}(\epsilon)<+\infty$ for all $\epsilon$, while $Y_0/\psi$ is bounded a.s. if and only if $\tilde{S}(\epsilon)<+\infty$ is finite for some $\epsilon>0$. 
%
%Careful scrutiny of the proofs reveals one further direct generalisation. 
%It also turns out if $\gamma$ is a function for which 
%\begin{equation} \label{eq.gammauniformquotientcompact}
%\liminf_{t\to\infty}
%\inf_{0\leq \theta\leq 1}
%\frac{\psi(t-\theta)}{\psi(t)} >0, \quad 
%\limsup_{t\to\infty} \sup_{0\leq \theta\leq 1}
%\frac{\psi(t-\theta)}{\psi(t)} <+\infty,
%\end{equation}
%all the statements in the above theorem still hold. 
%then the equivalences still hold. 
%
%A nice class of functions for which \eqref{eq.gammauniformquotientcompact} this is true are those of the form  $\gamma(t)=e^{\alpha t}\gamma_0(t)$ for $\alpha>0$ where $\gamma_0$ is subexponential, and for these we can characterise exponential--type growth bounds of order $\gamma$ for $Y_0$. 
%
%Note that we can always replace a subexponential function $\gamma_0$ with a differentiable function $\Gamma_0$ for which $\gamma_0(t)\sim \Gamma_0(t)$ and $\Gamma_0'(t)/\Gamma_0(t)\to 0$ as $t\to\infty$. Then $\Gamma(t)=e^{\alpha t}\Gamma_0(t)$ can be taken as the majorising function, and it has the same asymptotic behaviour as $\gamma$. However, it is also increasing for large $t$, because 
%\[
%\frac{\Gamma'(t)}{\Gamma(t)} = \alpha+ \frac{ \Gamma_0'(t)}{\Gamma_0(t)}\to \alpha>0, \quad t\to\infty.
%\]  
%This monotonicity is useful, since our deterministic results often rely on increasing majorising functions. Specifically, it will allow for necessary and sufficient conditions for exponential growth bounds for $Y$, where there are both stochastic and deterministic perturbations.  
%
%We are now in a position to give a very sharp result concerning the bounds on solutions of \eqref{eq.Y}, which is perturbed both deterministically and stochastically. 
%%In fact, the result is almost identical to that of Theorem~\ref{thm.Yto0bddchar}, once the results in the previous section are taken into account. For this reason, we merely state it. 
%\begin{theorem} \label{thm.Yto0bddpsichar}
%Let $f,\sigma$ obey \eqref{eq.fcns} and \eqref{eq.sigma}. Let $Y$ be the unique continuous adapted solution to \eqref{eq.Y}, and suppose $\tilde{S}_\gamma$ is given by \eqref{eq.Stilde}. 
%\begin{itemize}
%	\item[(a)] Let $\gamma$ be subexponential (or obey \eqref{eq.gammauniformquotientcompact}). The following are equivalent:
%	\begin{enumerate}
%		\item[(A)] $\int_{t-\theta}^t f(s)\,ds = \text{o}(\gamma(t)$ as $t\to\infty$ for all $\theta\in [0,1]$, $\tilde{S}_\gamma(\epsilon)<+\infty$ for all $\epsilon>0$;
%		\item[(B)] $Y(t)/\gamma(t)\to 0$ as $t\to\infty$ a.s. 
%	\end{enumerate}
%	\item[(b)] Let $\gamma$ be subexponential (or obey \eqref{eq.gammauniformquotientcompact}). The following are equivalent:
%	\begin{enumerate}
%		\item[(A)] There exists $\epsilon'>0$ such that $\tilde{S}_\gamma(\epsilon')<+\infty$ and \eqref{eq.fthetagammabdd} holds; 
%		\item[(B)] $Y(t)=O(\gamma(t))$ as $t\to\infty$ a.s. bounded. 
%	\end{enumerate}
%	\item[(iii)]  \hspace{0.1in} The following are equivalent:
%	\begin{enumerate}
%		\item[(A)] 	$\limsup_{t\to \infty} |Y(t)|/\gamma(t)\in (0,\infty)$, a.s.
%		\item[(B)] 	There exists $\bar{F}>0$ and $\epsilon'>0$ such that $\tilde{S}_\gamma(\epsilon')<+\infty$ and 
%		$\left|\int_{(t-\theta)^+}^t f(s)\,ds\right|\leq \bar{F}\gamma(t)$ for all $\theta\in [0,1]$ and $t\geq 0$, 
%		and at least one of the following holds:
%		\begin{gather} \label{eq.fgammabelow}
	%			\text{There exists $\theta'>0$ such that } \limsup_{t\to\infty} \frac{\left|\int_{t-\theta'}^t f(s)\,ds\right|}{\gamma(t)}>0, \\
	%			\label{eq.Sgammainfinite}
	%			\text{There exists $\epsilon''>0$ such that $\tilde{S}_\gamma(\epsilon'')=+\infty$.}
	%		\end{gather}
%	\end{enumerate}
%\end{itemize} 
%\end{theorem}
%Note that the results deal with both decay and growth bounds, according the whether the subexponential function tends to zero or to infinity.  
%
%A nice class of functions for which \eqref{eq.gammauniformquotientcompact} is true are those of the form  $\gamma(t)=e^{\alpha t}\gamma_0(t)$ for $\alpha>0$ where $\gamma_0$ is subexponential, and for these we can characterise exponential--type growth bounds of order $\gamma$ for $Y_0$; indeed, we can impose smoothness on $\gamma_0$ without cost, by virtue of the remarks after \eqref{eq.psisub}. 
%
%Note that we can always replace a subexponential function $\gamma_0$ with a differentiable function $\Gamma_0$ for which $\gamma_0(t)\sim \Gamma_0(t)$ and $\Gamma_0'(t)/\Gamma_0(t)\to 0$ as $t\to\infty$. Then $\Gamma(t)=e^{\alpha t}\Gamma_0(t)$ can be taken as the majorising function, and it has the same asymptotic behaviour as $\gamma$. However, it is also increasing for large $t$, because 
%$\Gamma'(t)/\Gamma(t) = \alpha+ \Gamma_0'(t)/\Gamma_0(t)\to \alpha>0$ as $t\to\infty$.
%This monotonicity is useful, since our deterministic results often rely on increasing majorising functions. Specifically, it will allow for necessary and sufficient conditions for exponential growth bounds for $Y$, where there are both stochastic and deterministic perturbations.  
%
%One final matter needs to be addressed; can we identify, as in the deterministic case, exactly the conditions under which $Y/\psi$ has a finite and positive limsup. This is more delicate, and requires  little more effort. 
%
%First, recall that Theorem~\ref{thm.Y0bddpsichar} shows that 
%\[
%\limsup_{t\to\infty} \frac{|Y_0(t)|}{\gamma(t)}\in (0,\infty),\, \text{a.s.} \Longleftrightarrow \exists \epsilon_1, \, \epsilon_2>0 \text{ such that } \tilde{S}(\epsilon_1)<+\infty, \tilde{S}(\epsilon_2)=+\infty
%\] 
%and we can likewise characterise when $\limsup_{t\to\infty} |y(t)|/\gamma(t)\in (0,\infty)$. What we need to show is that the following are equivalent:
%\begin{gather*}
%\limsup_{t\to\infty} \frac{|Y(t)|}{\gamma(t)}\in (0,\infty), \quad \text{a.s.}
%\\
%\limsup_{t\to\infty} \frac{|y(t)|}{\gamma(t)}<+\infty, \, 
%\limsup_{t\to\infty} \frac{|Y_0(t)|}{\gamma(t)}<+\infty\, \text{a.s.}\,
%\text{ with at least one limsup positive}. 
%\end{gather*}
%Call the $\limsup$s $\bar{Y}$, $\bar{y}$ and $\bar{Y}_0$. For the reverse implication, if $\bar{y}$ and $\bar{Y}_0$ are finite, then so is $\bar{Y}$, since $Y=y+Y_0$. To complete the proof, it is necessary to show that when both are positive, then so is $\bar{Y}$. 
%
%To this end, the following lemma is required. 
%\begin{lemma}  \label{lemma.Mn}
%	Suppose $g\in C([0,\infty);\mathbb{R})$ and
%	\[
%	\int_{0}^\infty g^2(s)\,ds = +\infty
%	\]
%	Let $t_0=0$ and $(t_n)_{n\geq 0}$ be a deterministic increasing sequence with $t_n\to\infty$ as $n\to\infty$.
%	Let $B=\{B(t):t\geq 0\}$ be a standard $\mathcal{F}^B$ Brownian motion and define $M=\{M(n):n\in\mathbb{N}\}$ by
%	\[
%	M(n)=\int_{0}^{t_n} g(s)\,dB(s), \quad n=0,1,2\ldots.
%	\]
%	Then
%	\begin{equation} \label{eq.MnsqrtMn}
%		\limsup_{n\to\infty} \frac{M(n)}{\sqrt{\langle M\rangle(n)}}=+\infty, \quad
%		\liminf_{n\to\infty}  \frac{M(n)}{\sqrt{\langle M\rangle(n)}}=-\infty, \quad \text{a.s.},
%	\end{equation}
%	and in particular
%	\begin{equation} \label{eq.Mnfluct}
%		\limsup_{n\to\infty} M(n)=+\infty, \quad \liminf_{n\to\infty}  M(n)=-\infty, \quad\text{a.s.}
%	\end{equation}
%\end{lemma}
%\begin{proof}
%	Since $-M(n)$ has the same distribution as $M(n)$, we need only prove the first statement in \eqref{eq.MnsqrtMn}.
%	Define $\mathcal{G}(n)=\mathcal{F}^{B}(t_n)$ for $n\in\mathbb{N}$. Then $M$ is a $\mathcal{G}_n$--martingale. We have
%	\[
%	\langle M\rangle (n)=\int_{0}^{t_n} g^2(s)\,ds\to\infty, \quad n\to\infty.
%	\]
%	Hence we see that \eqref{eq.Mnfluct} is a direct consequence of \eqref{eq.MnsqrtMn}.
%	
%	Define $v_{n}^2=\int_{t_{n-1}}^{t_n} g^2(s)\,ds$ for $n=1,2,3\ldots$ and
%	$\xi(n)=\int_{t_{n-1}}^{t_n} g(s)\,dB(s)$  for $n=1,2,3\ldots$. Observe that $(\xi(j))_{j\geq 1}$ is a sequence of
%	independent zero mean normal random variables, where $\text{Var}[\xi(n)]=v_n^2$,  that $M(n)=\sum_{j=1}^n \xi(j)$ and
%	that $\langle M\rangle(n)=\sum_{j=1}^n v_j^2$. Let $c>0$ and define
%	\[
%	A_c=\{\omega: \limsup_{n\to\infty} \frac{M(n,\omega)}{\sqrt{\langle M \rangle(n)}}>c\},
%	\quad
%	A=\{\omega: \limsup_{n\to\infty} \frac{M(n,\omega)}{\sqrt{\langle M \rangle(n)}}=+\infty\}.
%	\]
%	Then $A_c\downarrow A$ and $A_c$ is a tail event. If
%	\begin{equation} \label{eq.zeroonecondn}
%		\mathbb{P}[A_c]>0 \text{ for all $c>0$}
%	\end{equation}
%	then by the Zero--One Law, we have that $\mathbb{P}[A_c]=1$, and so $\mathbb{P}[A]=\lim_{c\to\infty} \mathbb{P}[A_c]=1$.
%	It remains therefore to prove \eqref{eq.zeroonecondn}.
%	
%	To do this, note that
%	\begin{itemize}
%		\item[(a)] For any process $\eta=\{\eta(n):n\in \mathbb{N}\}$ we have
%		\[
%		\{\omega: \limsup_{n\to\infty} \eta(n,\omega)>x\} \supseteq \{ \omega: \eta(n,\omega)>x \text{ i.o.}\} \quad\text{for all $x\in\mathbb{R}$};
%		\]
%		\item[(b)] For any sequence of events $(B_n)_{n\geq 1}$
%		\[
%		\mathbb{P}[B_n\text{ i.o.}]\geq \limsup_{n\to\infty}\mathbb{P}[B_n].
%		\]
%	\end{itemize}
%	Therefore applying observation (a) and then observation (b) we get
%	\begin{align*}
%		\mathbb{P}[A_c]&=\mathbb{P}\left[ \omega: \limsup_{n\to\infty}\frac{M(n,\omega)}{\sqrt{\langle M\rangle(n)}}>c \right]\\
%		&\geq \mathbb{P}\left[\omega: \frac{M(n,\omega)}{\sqrt{\langle M\rangle(n)}}>c  \quad\text{i.o.}\right]\\
%		&\geq \limsup_{n\to\infty} \mathbb{P}\left[\omega: \frac{M(n,\omega)}{\sqrt{\langle M\rangle(n)}}>c \right].
%	\end{align*}
%	Note that $M(n)/\sqrt{\langle M\rangle (n)}$ is normally distributed with zero mean and unit variance. Thus
%	\[
%	\mathbb{P}\left[\omega: \frac{M(n,\omega)}{\sqrt{\langle M\rangle(n)}}>c \right]=1-\Phi(c),
%	\]
%	so we have $\mathbb{P}[A_c]\geq 1-\Phi(c)>0$ for each $c>0$. This proves \eqref{eq.zeroonecondn} and with it, the result.
%\end{proof}
%
%\begin{lemma}
%	Suppose that $\gamma$ is either subexponential or increasing. Then 
%	$\bar{y}>0$, $\bar{Y}_0>0$ implies $\bar{Y}>0$. 
%\end{lemma}
%\begin{proof}
%	By hypothesis $\overline{y}\in (0,\infty)$. Since $\tilde{y}:=y/\gamma$ is deterministic, this means that there is an increasing and deterministic sequence $(t_n)_{n\geq 1}$ such that $|y(t_n)|/\gamma(t_n)\to\overline{y}$ as $n\to\infty$. Therefore there exists a subsequence $(t_{n_j})_{j\geq 1}$ along which either $\tilde{y}(t_{n_j})\to \overline{y}>0$ as $j\to\infty$ or $\tilde{y}(t_{n_j})\to -\overline{y}$ as $j\to\infty$. Define
%	\[
%	M(j)=\int_{0}^{t_{n_j}} e^s\sigma(s)\,dB(s), \quad j\geq 0.
%	\]
%	Notice that $Y_0(t_{n_j})=M(j)/e^{t_{n_j}}$. 
%	
%	Suppose that $\int_0^\infty e^{2s}\sigma^2(s)\,ds <+\infty$. Then by martingale convergence, it follows from the fact  
%	\[
%	Y_0(t) e^t = \xi + \int_0^t e^{s}\sigma(s)\,dB(s)
%	\]
%	that $Y_0(t)e^t$ tends to a finite (random) limit $C^\ast$ a.s. 
%	Then, if $\gamma$ is either subexponential or incerasing, we have that $\gamma(t)e^t\to\infty$ as $t\to\infty$, so 
%	\[
%	\lim_{t\to\infty} \frac{Y_0(t)}{\gamma(t)} = \lim_{t\to\infty} \frac{Y_0(t)}{e^{-t}}\frac{1}{e^t\gamma(t)} = 0, \quad \text{a.s.},
%	\]
%	which contradicts the fact that $\limsup_{t\to\infty} |Y_0(t)|/\gamma(t)>0$ a.s.
%	Therefore, we must have that $\int_0^\infty e^{2s}\sigma^2(s)\,ds =+\infty$.
%	
%	Then by Lemma~\ref{lemma.Mn} we have that there is an a.s. event $A_1$ such that
%	\[
%	A_1=\{\omega: \liminf_{j\to\infty}  M(j,\omega)=-\infty, \quad \limsup_{j\to\infty} M(j,\omega)=+\infty\}.
%	\]
%	Next, suppose that there is an event $A=\{\omega: \limsup_{t\to\infty} |Y(t,\omega)|/\gamma(t)=0\}$ for which $\mathbb{P}[A]>0$. If $A_2=A\cap A_1$, then $\mathbb{P}[A_2]=\mathbb{P}[A]>0$.
%	
%	Consider the case where $\tilde{y}(t_{n_j})\to \overline{y}>0$ as $j\to\infty$. If $\omega\in A_2$, then $Y_0(t_{n_j},\omega)/\gamma(t_{n_j})\to -\overline{y}$ as $j\to\infty$. 
%	Since 
%	\[
%	M(j)=e^{t_{n_j}}\gamma(t_{n_j})\frac{Y_0(t_{n_j})}{\gamma(t_{n_j})},
%	\] 
%	and the deterministic prefactor tends to infinity as $j\to\infty$  (by monotonicity or subexponetiality of $\gamma$), we have that $M(j,\omega)\to -\infty$ as $j\to\infty$,
%	which contradicts the fact $\limsup_{j\to\infty} M(j,\omega)=+\infty$ for $\omega\in A_2$. Therefore $\mathbb{P}[A]=0$.
%	
%	On the other hand, suppose that we have $\tilde{y}(t_{n_j})\to -\overline{y}>0$ as $j\to\infty$. If $\omega\in A_2$, then $Y(t_{n_j},\omega)/\gamma(t_{n_j})\to \overline{y}$ as $j\to\infty$. 
%	Since $M(j)=e^{t_{n_j}}Y(t_{n_j})$, we have that $M(j,\omega)\to +\infty$ as $j\to\infty$,
%	which contradicts the fact $\liminf_{j\to\infty} M(j,\omega)=-\infty$ for $\omega\in A_2$. Therefore $\mathbb{P}[A]=0$.
%	Therefore, we have that $\mathbb{P}[\overline{A}]=1$, which implies that
%	\[
%	\limsup_{t\to\infty} |Y(t)|/\gamma(t)>0, \quad\text{a.s.}
%	\]
%	as required.
%\end{proof}

%\begin{lemma}
%	Suppose $\gamma$ is subexponential or obey \eqref{eq.gammauniformquotientcompact}
%		\begin{equation} \label{eq.gammauniformquotientcompact}
%		\liminf_{t\to\infty}
%		\inf_{0\leq \theta\leq 1}
%		\frac{\gamma(t-\theta)}{\gamma(t)} >0, \quad 
%		\limsup_{t\to\infty} \sup_{0\leq \theta\leq 1}
%		\frac{\gamma(t-\theta)}{\gamma(t)} <+\infty.
%	\end{equation}
%The following are equivalent:
%\begin{enumerate}
%	\item[(A)] 
%	\[
%	\limsup_{t\to \infty} \frac{|Y(t)|}{\gamma(t)}\in (0,\infty), \quad\text{a.s.}
%	\] 
%		\item[(B)]  
%		\[
%		\limsup_{t\to\infty} \frac{|y(t)|}{\gamma(t)}<+\infty, \quad 
%		\limsup_{t\to \infty} \frac{|Y_0(t)|}{\gamma(t)}<+\infty, \quad\text{a.s.}
%		\] 
%		and at least one of these limsups is strictly positive.
%	\item[(B)] 
%Suppose additionally  $\gamma$ obeys 
%\begin{equation} \label{eq.gammauniformquotientcompact}
%	\liminf_{t\to\infty}
%	\inf_{0\leq \theta\leq 1}
%	\frac{\gamma(t-\theta)}{\gamma(t)} >0, \quad 
%	\limsup_{t\to\infty} \sup_{0\leq \theta\leq 1}
%	\frac{\gamma(t-\theta)}{\gamma(t)} <+\infty.
%\end{equation}
%	There exists $\bar{F}>0$ and $\epsilon'>0$ such that 
%	\begin{gather*}
%		\left|\int_{(t-\theta)^+}^t f(s)\,ds\right|\leq \bar{F}\gamma(t), \quad \theta\in [0,1], \, t\geq 0, \\
%		\tilde{S}(\epsilon) := \sum_{n=1}^\infty \frac{\gamma(n-1)}{\left(\int_{n-1}^{n}\sigma^2(s)ds\right)^{1/2}}\exp\left(-\epsilon \frac{\gamma^2(n-1)}{\int_{n-1}^{n}\sigma^2(s)ds}\right)<+\infty \text{ for some $\epsilon=\epsilon'$},
%	\end{gather*}
%	and at least one of the following prevails:
%	\begin{gather*}
%		\text{There exists $\theta'>0$ such that } \limsup_{t\to\infty} \frac{\left|\int_{t-\theta'}^t f(s)\,ds\right|}{\gamma(t)}>0, \\
%		\text{There exists $\epsilon''>0$ such that $\tilde{S}(\epsilon'')=+\infty$.}
%	\end{gather*}
%\end{enumerate}
%\end{lemma}
%\begin{proof}
%	We have already shown that (B) and (C) are equivalent. 	
%	To show (B) implies (A), it is obvious that both $y$ and $Y_0$ being bounded by $\gamma$ implies $Y$ is, so the limsup in (A) is finite. If exactly one of the limsups in (B) is positive, and the other is zero, then the limsup in (A) is positive. On the other hand, if both limsups in (B) are positive, then by the previous lemma, so is the the limsup in (A). Thus (B) implies (A). 
%	
%	Conversely, we have already shown that the statement in (A) implies the finiteness of the limsups in (B). But at least one of these must be positive, since if both are zero, the limsup in (A) would be zero, causing a contradiction.   
%\end{proof}
%
%\section{Asymptotic behaviour of the stochastic pantograph equation}
%
%With all the results in hand for the singly deterministically and singly stochastically perturbed equations, it is quite straightforward to give a good characterisation of the growth or decay rates of solutions of the doubly perturbed equation. 
%
%We start with results with  conditions that may be more delicate to check, and which have some side conditions on the problem data, but which characterise certain types of a.s. asymptotic behaviour. Later in this section, we give easier--to--check and sharp sufficient conditions which do not place side restrictions on the data, but do not give necessary and sufficient conditions for certain types of asymptotic behaviour. Therefore, it can be seen that both types of results are complementary.
%
%We start with a result with power--law type perturbations that are bigger than the solution of the unperturbed equation. These can be growing or decaying. 
%\begin{theorem}
%Suppose that $q\in (0,1)$ and \eqref{eq.unforcedasstable} holds. Let $f,\sigma\in C([0,\infty);\mathbb{R})$ and $X$ be the unique continuous adapted solution to \eqref{eq.stochforced2}. Suppose $\tilde{S}_\gamma$ is given by \eqref{eq.Stilde}.
%Suppose that $\gamma$ is  in $\text{RV}_\infty(\eta)$ with $\eta>\kappa$.  
%\begin{itemize}
%	\item[(I)] \hspace{.1in} The following are equivalent:
%	\begin{enumerate}
%		\item[(A)] There exists $\epsilon'>0$ such that $\tilde{S}_\gamma(\epsilon')<+\infty$ and
%		\eqref{eq.fthetagammabdd} holds;  
%		%\tilde{S}(\epsilon) := \sum_{n=1}^\infty %\frac{\gamma(n-1)}{\left(\int_{n-1}^{n}\sigma^2(s)ds\right)^{1/2}}\exp\left(-\epsilon %\frac{\gamma^2(n-1)}{\int_{n-1}^{n}\sigma^2(s)ds}\right)<+\infty \text{ for some $\epsilon=\epsilon'$};		
%		%\end{gather*} 
%		\item[(B)] $X(t)=\text{O}(\gamma(t))$ a.s.
%	\end{enumerate}
%	\item[(II)] \hspace{.1in} The following are equivalent:
%	\begin{enumerate}
%		\item[(A)]  For all  $\epsilon>0$, $\tilde{S}_\gamma(\epsilon)<+\infty$ and \eqref{eq.fthetagamma0} holds ; 
%		%	\tilde{S}(\epsilon) := \sum_{n=1}^\infty \frac{\gamma(n-1)}{\left(\int_{n-1}^{n}\sigma^2(s)ds\right)^{1/2}}\exp\left(-\epsilon \frac{\gamma^2(n-1)}{\int_{n-1}^{n}\sigma^2(s)ds}\right)<+\infty;		
%		%\end{gather*} 
%		\item[(B)] $X(t)=\text{o}(\gamma(t))$ a.s.
%	\end{enumerate}
%	\item[(III)] \hspace{0.1in} The following are equivalent
%	\begin{enumerate}
%		\item[(A)] There exists $\epsilon'>0$ such that $\tilde{S}_\gamma(\epsilon')<+\infty$ and \eqref{eq.fthetagammabdd} holds, and  
%		at least one of \eqref{eq.fgammabelow}
%		and \eqref{eq.Sgammainfinite} holds. 
%		\item[(B)] $\limsup_{t\to\infty} |X(t)|/\gamma(t)\in (0,\infty)$ a.s. 
%	\end{enumerate}
%\end{itemize}
%\end{theorem} 
%\begin{proof}
%The entire proof follows from fact that, when $\gamma\in \text{RV}_\infty(\eta)$ and $\eta>\kappa$, then  
%\[
%\limsup_{t\to\infty} \frac{|X(t)|}{\gamma(t)}\in (0,\infty), \quad\text{a.s.}
%\Longleftrightarrow 
%\limsup_{t\to\infty} \frac{|Y(t)|}{\gamma(t)}\in (0,\infty), \quad \text{a.s.} 
%\]
%The statement on $Y$ is equivalent to (A), proving the result. 
%However, the above equivalence can be proven using Theorem~\ref{thm.detlargepert} part (III), applying the argument path by path, noticing that $Y$ can be represented in terms of $X$ and that $Z=X-Y$ obeys a pantograph equation of the type considered in Theorem~\ref{thm.detlargepert}. 
%We now prove the equivalence. If $Y(y)=\text{O}(\gamma(t))$, then $\phi(t)=bY(t)+aY(qt)$ is also $\text{O}(\gamma(t))$, using the regular variation of $\gamma$. Since $\eta>\kappa$, this implies that $Z(t)=X(t)-Y(t)$ obeys $Z'(t)=bZ(t)+aZ(qt)+\phi(t)$, and therefore we have $Z(t)=\text{O}(\gamma(t))$. Hence $X(t)=Z(t)+Y(t)=\text{O}(\gamma(t))$. Next, given that) $\limsup_{t\to\infty}Y(t)/\gamma(t)>0$, we assume by way of contradiction that $X(t)=\text{o}(\gamma(t))$ on an event of positive probability. But this implies that then $Y(t)=\text{o}(\gamma(t))$ on that event, forcing a contradiction. Hence the reverse implication is true. 
%
%In the other direction, if $X(t)=\text{O}(\gamma(t))$, then $Y(t)=\text{O}(\gamma(t))$. Now, by hypothesis we have $\limsup_{t\to\infty}Y(t)/\gamma(t)>0$, but assume by way of contradiction that $X(t)=\text{o}(\gamma(t))$ on  an event of positive probability. But this implies that $Y(t)=\text{o}(\gamma(t))$ on that event, contradicting the hypothesis. Hence the forward implication is also true. 
%\end{proof}
%
%Notice the hypothesis $Y(t)=O(\gamma(t))$ as $t\to\infty$ for $\gamma\in \text{RV}_\infty(\eta)$ is equivalent to 
%\[
%S_\gamma(\epsilon')<+\infty \text{ for some $\epsilon'>0$}, \quad \left|\int_{t-\theta}^t f(s)\,ds \right|\leq K\gamma(t), \quad t\geq 0,\,\theta\in [0,1],
%\]
%which  can be checked from the problem data. Likewise, we have a characterisation of $Y(t)=o(\gamma(t))$ and $\limsup_{t\to\infty} |Y(t)|/\gamma(t)\in (0,\infty)$, all supplied by Theorem~\ref{thm.Yto0bddpsichar}.
%
%For perturbations decaying faster than the solution of \eqref{eq.unforced}, the following holds. 
%\begin{theorem} \label{thm.smallstochpert}
%Suppose $q\in (0,1)$ and \eqref{eq.unforcedasstable} holds. Let $f,\sigma\in C([0,\infty);\mathbb{R})$ and $X$, $Y$ be the unique continuous adapted solutions to \eqref{eq.stochforced2}, \eqref{eq.Y}. Let  $\gamma\in \text{RV}_\infty(\eta)$ for $\eta\leq \kappa$.
%%Suppose $\tilde{S}$ is given by \eqref{eq.Stilde}.
%%Let $\gamma$ be regularly varying with index $\eta>\kappa$. 
%\begin{itemize}
%	\item[(i)] Let $Y(t)=\text{O}(\gamma(t))$ as $t\to\infty$ a.s.
%	\begin{enumerate}
%		\item[(A)] 
%		If $\eta< \kappa$, then $X(t)=O(t^\kappa)$ as $t\to\infty$ a.s.
%		\item[(B)] If $\eta=\kappa$ and  $t\mapsto \gamma(t)/t^{\kappa+1}$ is in $L^1$, then 
%		$X(t)=O(t^\kappa)$ as $t\to\infty$ a.s. 
%		Moreover, if $\limsup_{t\to\infty} |Y(t)|/\gamma(t)\in (0,\infty)$, then 
%		\[
%		\limsup_{t\to\infty} \frac{|X(t)|}{\gamma(t)}>0, \quad\limsup_{t\to\infty} \frac{|X(t)|}{t^\kappa}<+\infty, \quad \text{a.s.}
%		\]
%		\item[(C)] If $\eta=\kappa$ and  $t\mapsto\gamma(t)/t^{\kappa+1}$ is not in $L^1$, then 
%		\[
%		X(t)=\text{O}\left(t^\kappa\int_1^t \frac{\gamma(s)}{s^{\kappa+1}}\,ds\right), \,t\to\infty \quad \text{a.s.}
%		\]
%		Moreover, if $\limsup_{t\to\infty} |Y(t)|/\gamma(t)\in (0,\infty)$, then 
%		\[
%		\limsup_{t\to\infty} \frac{|X(t)|}{\gamma(t)}>0, \quad \limsup_{t\to\infty} \frac{|X(t)|}{t^\kappa\int_1^t \frac{\gamma(s)}{s^{\kappa+1}}\,ds}<+\infty, \quad \text{a.s.}
%		\]
%	\end{enumerate}  
%	\item[(ii)] If $X(t)=\text{O}(\gamma(t))$ as $t\to\infty$, then $Y(t)=\text{O}(\gamma(t))$ as $t\to\infty$ a.s. Likewise, if $X(t)=\text{o}(\gamma(t))$ as $t\to\infty$, then $Y(t)=o(\gamma(t))$ as $t\to\infty$.   
%	%	\item[(C)] We have 
%	%	\[
%	%\limsup_{t\to\infty} \frac{|Y(t)|}{t^{\kappa+\epsilon}}=0, \,\text{a.s.}\,
%	%\Longleftrightarrow 
%	%	\limsup_{t\to\infty} \frac{\log |Y(t)|}{\log t}\leq \kappa, \,\text{a.s.}\,
%	%	\Longleftrightarrow \limsup_{t\to\infty} \frac{\log |X(t)|}{\log t}\leq \kappa,\, \text{a.s.}
%	%	\]
%\end{itemize}
%\end{theorem}
%\begin{proof}
%	For parts (A), (B), (C), we note that if $Y(t)=O(\gamma(t))$, then $\varphi(t)=O(\gamma(t))$. Hence $Z(t)=O(t^\kappa)$ 
%	To prove part (A), there are three cases to consider. 
%	If $Y(t)=\text{O}(\gamma(t))$ as $t\to\infty$ where $\gamma\in \text{RV}_\infty(\kappa)$ and $\gamma(t)/t^{\kappa+1}$ is integrable, then by Theorem~\ref{thm.growthbounds2}, $\gamma(t)=o(t^\kappa)$. Moreover, $Z(t)=\text{O}(t^\kappa)$ as $t\to\infty$ so $X(t)=\text{O}(t^\kappa)$ as $t\to\infty$, a.s. 
%	
%	If $Y(t)=\text{O}(\gamma(t))$ as $t\to\infty$ where $\gamma\in \text{RV}_\infty(\kappa)$ and $\gamma(t)/t^{\kappa+1}$ is not integrable, by Theorem~\ref{thm.growthbounds2}, we have $Z(t)=\text{O}(\gamma(t))$ as $t\to\infty$, so $X(t)=\text{O}(\gamma(t))$ as $t\to\infty$.
%	
%	If $Y(t)=\text{O}(\gamma(t))$ as $t\to\infty$ where $\gamma\in \text{RV}_\infty(\eta)$ and $\eta<\kappa$, then by Theorem~\ref{thm.growthbounds2}, $Z(t)=\text{O}(t^\kappa)$ as $t\to\infty$ and so $X(t)=\text{O}(t^\kappa)$ as $t\to\infty$.
%	
%	To prove part (B), note that  $X(t)=\text{O}(t^\kappa)$ implies $bX(t)+aX(qt) =\text{O}(t^\kappa)$, so using \eqref{eq.YreptildeX} we see that $Y(t)=\text{O}(t^\kappa)$ as $t\to\infty$. 
%	
%	The last equivalence in part (C) is a direct consequence of parts (A) and (B). 
%\end{proof}
%
%Now we consider perturbations which grow faster than a power law. However, the condition \eqref{eq.gammauniformquotientcompact}, on which Theorem~\ref{thm.Yto0bddpsichar} relies, limits a characterisation of the asymptotic behaviour to the case when the perturbations grow no faster than exponentially.
%
%\begin{theorem}
%Suppose that $q\in (0,1)$ and \eqref{eq.unforcedasstable} holds. Let $f,\sigma\in C([0,\infty);\mathbb{R})$ and $X$ and $Y$ be the unique continuous adapted solutions to \eqref{eq.stochforced2} and \eqref{eq.Y}. Let $\gamma$ be positive, increasing, in $C^1$ and $\gamma(t)\to\infty$ as $t\to\infty$. 
%\begin{enumerate}
%	\item[(i)] $Y(t)=O(\gamma(t))$ a.s. $\Longleftrightarrow$  $X(t)=O(\gamma(t))$ a.s.	
%	\item[(ii)]	$Y(t)=o(\gamma(t))$ a.s. $\Longleftrightarrow$ $X(t)=o(\gamma(t))$ a.s. 
%	\item[(iii)] $\limsup_{t\to\infty} |Y(t)|/\gamma(t) \in (0,\infty)$ a.s. $\Longleftrightarrow$ $\limsup_{t\to\infty} |X(t)|/\gamma(t) \in (0,\infty)$ a.s.
%\end{enumerate}
%\end{theorem}
%\begin{proof}
%We prove forward implications first. For (i) $Y(t)=O(\gamma(t))$, $t\to\infty$ gives  $\phi(t)=bY(t)+aY(qt) = \text{O}(\gamma(t))$ as $t\to\infty$. By part (A) of Theorem \ref{thm.monotonepert}, $Z(t)=X(t)-Y(t)$ is $\text{O}(\gamma(t))$ as $t\to\infty$. Thus $X(t)=Z(t)+Y(t)$ is $O(\gamma(t))$ as $t\to\infty$.
%
%For part (ii), if $Y(t)=o(\gamma(t))$, as $t\to\infty$ then $\phi(t)=bY(t)+aY(qt) = \text{o}(\gamma(t))$ as $t\to\infty$. By part (B) of Theorem \ref{thm.monotonepert}, $Z(t)=X(t)-Y(t)$ is $\text{o}(\gamma(t))$ as $t\to\infty$. Thus $X(t)=Z(t)+Y(t)$ is $o(\gamma(t))$ as $t\to\infty$.
%
%For part (iii), since the limsup for $Y$ is finite, the limsup for $X$ is also. Suppose by way of contradiction that, on an event of positive probability, the limsup for $X$ is zero. Then, by part (II) of Theorem \ref{thm.detlargepert}, we are led to $Y(t)=\text{o}(\gamma(t))$ as $t\to\infty$ on this event. This forces a contradiction, so the limsup for $X$ is positive.   
%
%Reverse implications follow by applying pathwise arguments in Theorem~\ref{thm.detlargepert}.
%\end{proof}
%
%We now give a result which relaxes exponential bounds on $\gamma$ implied by \eqref{eq.gammauniformquotientcompact} and gives conditions that are easier to check, at the expense of characterising the asymptotic behaviour. To do this, note that 
%$Y_0$ has the representation $Y_0(t)=e^{-t} M(t)$ where 
%$M(t) = \int_0^t e^s\sigma(s)\,dB(s)$ is a martingale with quadratic variation
%\[
%\langle M\rangle(t)=\int_0^t e^{2s}\sigma^2(s)\,ds, \quad t\geq 0.
%\]
%In the case that $\int_0^\infty e^{2s}\sigma^2(s)\,ds<+\infty$, the quadratic variation converges, so by the martingale convergence theorem, we have $Y_0(t)\sim Y^\ast e^{-t}$ as $t\to\infty$ a.s.  
%
%In the case that $\int_0^\infty e^{2s}\sigma^2(s)\,ds<+\infty$, the quadratic variation diverges, so by the Law of the Iterated Logarithm we have
%\[
%\limsup_{t\to\infty} \frac{e^t |Y_0(t)|}{\sqrt{2\int_0^t e^{2s}\sigma^2(s)\,ds\log\log \int_0^t e^{2s}\sigma^2(s)\,ds}}=
%\limsup_{t\to\infty} \frac{|M(t)|}{\sqrt{2\langle M\rangle(t) \log\log \langle %M\rangle(t)}}=1,
%\]
%so
%$\limsup_{t\to\infty}|Y_0(t)|/\gamma_2(t)=1$ a.s. where 
%\begin{equation} \label{eq.gamma2}
%\gamma_2(t) = \sqrt{2\int_0^t e^{-2(t-s)}\sigma^2(s)\,ds\log\log \int_0^t e^{2s}\sigma^2(s)\,ds}
%\end{equation}
%Thus we have $\limsup_{t\to\infty}|Y_0(t)|/\gamma_2(t)\in (0,\infty)1$ for $\gamma_2$ given by \eqref{eq.gamma2} or $\gamma_2(t) = e^{-t}$ according as to whether $e^{2t}\sigma^2(t)$ is non--integrable or integrable. Let us rule out the integrable case, since in this case the stochastic perturbation is very small, and the dynamics will be determined largely by the deterministic perturbation. 
%
%Next, let us assume the existence of a $\gamma_1$, and $K>0$ and a $\theta'>0$ such that 
%\begin{equation} \label{eq.gamma1}
%\left|\int_{t-\theta}^{t} f(s)\,ds\right| \leq K\gamma_1(t), \quad \limsup_{t\to\infty} \frac{|\int_{t-\theta'}^t f(s)\,ds|}{\gamma_1(t)}>0,
%\end{equation}
%which is equivalent to $\limsup_{t\to\infty} |y(t)|/\gamma_1(t)\in (0,\infty)$ if  $\gamma_1$ is increasing or subexponential. The following result is easy to apply, does not impose pointwise bounds on $f$ and $\sigma$, is applicable to many perturbing functions, and gives sharp results.
%\begin{theorem} 
%Suppose that $q\in (0,1)$ and \eqref{eq.unforcedasstable} holds. Let $f,\sigma\in C([0,\infty);\mathbb{R})$ and $X$ and $Y$ be the unique continuous adapted solutions to \eqref{eq.stochforced2} and \eqref{eq.Y}. Let $\gamma_1$ and $\gamma_2$ be given by \eqref{eq.gamma1} and \eqref{eq.gamma2} respectively, and suppose $\gamma=\max(\gamma_1,\gamma_2)$. 
%\begin{itemize}
%	\item[(I)] \hspace{0.1in} Suppose either (i) $\gamma$ is increasing; or (ii) $\gamma$ is  subexponential and obeys $\liminf_{t\to\infty}\log \gamma(t)/\log t>\kappa$. 
%	If $\gamma_2(t)=o(\gamma_1(t))$ as $t\to\infty$, or 
%	$\gamma_1(t)=o(\gamma_2(t))$ as $t\to\infty$, or $\gamma_1(t)=O(\gamma_2(t))$ and $\gamma_2(t)=O(\gamma_1(t))$ as $t\to\infty$, then 
%	$
%	\limsup_{t\to\infty} |X(t)|/\gamma(t)\in (0,\infty)$, a.s.
%	\item[(II)] \hspace{0.1in} If $\limsup_{t\to\infty} \log\gamma(t)/\log t <\kappa$, then $X(t)=O(t^\kappa)$ as $t\to\infty$ a.s.
%	\item[(III)] \hspace{0.1in} If $\lim_{t\to\infty} \log\gamma(t)/\log t =\kappa$, and $\gamma$ is regularly varying, then 
%	\[
%	X(t)=\text{O}\left(t^\kappa \int_1^t \frac{\gamma(s)}{s^{\kappa+1}}\,ds \right) \quad t\to\infty, \text{a.s.}
%	\]
%\end{itemize}
%\end{theorem} 
%An important ingredient in proving the positivity of the limsup in part (I)(c) is that, if $\gamma$ is increasing or subexponential, and $\limsup_{t\to\infty} |y(t)|/\gamma(t)\in (0,\infty)$, $\limsup_{t\to\infty} |Y_0(t)|/\gamma(t)\in (0,\infty)$, then $\limsup_{t\to\infty} |Y_0(t)+y(t)|/\gamma(t)\in (0,\infty)$ a.s. (see Lemma in Appleby and Lawless). 
\vspace{-0.25in}

\section{Developments and Further Work}
The principal purpose of this paper was to develop necessary and sufficient conditions for the convergence of perturbed solutions of a celebrated unbounded delay differential equation to the  equilibrium solution of the underlying unforced equation. Making these very particular choices, we have demonstrated that the forcing terms which are necessary and sufficient to preserve the asymptotic stability are precisely those needed to give the convergence to equilibrium of simple linear differential equations (whether deterministic or stochastic). 

One reason to focus on the pantograph equation, rather than to consider more general equations, is to exploit the fact that the asymptotic behaviour of the unperturbed pantograph equation is (essentially) completely understood, whereas for equations with general unbounded delay, and finite dimensional equations, for example, complete characterisations of the asymptotic behaviour are not always available. This makes it harder to determine whether gaps in the analysis for perturbed equations are purely the result of shortcomings in the perturbed (or stochastic) analysis, or may instead reflect the incomplete state of knowledge in the unperturbed systems. The abundance of necessary and sufficient conditions for various phenomena in the preceding questions supports this conservative approach.   

On the other hand, pantograph equations represent only a particular example of equations with unbounded delay, and it can be the case that the arguments used to analyse scalar equations can possess special features which make generalisation to higher dimensions difficult, or even impossible. Therefore, in this section, we give  arguments to show that a number of generalisations are possible, despite possibly incomplete knowledge about underlying unperturbed equations, or other important auxiliary equations (e.g., the SDE satisfied by $Y$). Leaving aside questions about growth or decay rates for general problems momentarily, several avenues for generalisation may be readily stated:
\begin{enumerate}
\item[(i)] Equations with general unbounded delay, rather than simply proportional delay;
\item[(ii)] Multi--dimensional equations;
\item[(iii)] Equations with multiplicative, rather than additive noise e.g., equations of the form 
\[
dX(t)=(bX(t)+aX(qt))\,dt + \sigma X(t)\,dB(t), \quad t\geq 0.
\]
\end{enumerate}   
In the following subsections, we sketch details of how the results in the earlier part of the paper can be generalised (rather successfully) to cover these cases. In cases (i) and (ii), the machinery developed in earlier sections can mostly be generalised in reasonably obvious ways (with some new problems emerging), while in case (iii), although the approach of the earlier part of the paper is not useful in the particulars, it  nevertheless suggests an analogous approach to attacking the equations in question which seems to generate some interesting preliminary results.    
\subsection{Equations with general unbounded delay}
For the first question, the framework of the present work is sufficient to provide a positive answer. Consider now the SFDE
\begin{equation}\label{eq.sfde}
dX(t)=(bX(t)+aX(t-\tau(t)))+f(t))\,dt + \sigma(t)\,dB(t), \quad t\geq 0,
\end{equation}
where $\tau$ is a positive continuous function with $\sup_{t\geq 0} t-\tau(t)=:-\bar{\tau}>-\infty$ and $t-\tau(t)\to \infty$ as $t\to\infty$. In this case, we need an initial function $\psi\in C([-\bar{\tau},\infty);\mathbb{R})$ so that if  $X=\psi$ on $[-\bar{\tau},0]$, then there will exist a unique continuous adapted solution of the equation on $[-\bar{\tau},\infty)$. Confining attention to the situation when $X(t)\to 0$ as $t\to\infty$ a.s., we may proceed as above to show that if $X(t)\to 0$ as $t\to\infty$ a.s., then $Y(t)\to 0$ as $t\to\infty$ a.s. We do this by writing $dX(t)=\{-X(t)+(1+b)X(t)+aX(t-\tau(t))+f(t)\}\,dt + \sigma(t)\,dB(t)$, and now, by treating everything apart from $-X(t)$ on the right hand side as a perturbation, using stochastic integration by parts, and rearranging for $Y$, we get  
\begin{equation} \label{eq.YXgeneral}
Y(t)= X(t)-e^{-t}X(0)-\int_{0}^t e^{-(t-s)}\{(1+b)X(s)+aX(s-\tau(s))\}\,ds, \quad t\geq 0.
\end{equation}
Since $t-\tau(t)\to \infty$ as $t\to\infty$, the term in curly brackets tends to zero as $s\to\infty$, so the integral tends to zero, and thus $Y(t)\to 0$ as $t\to\infty$ a.s., as claimed.

To show $Y(t)\to 0$ as $t\to\infty$ implies  $X(t)\to 0$ as $t\to\infty$ a.s., we may define 
$Z(t)=X(t)-Y(t)$ for $t\geq -\bar{\tau}$ (extending $Y$ to be zero on $[-\bar{\tau},0]$), in which case $Z=\psi$ on $[-\bar{\tau},0]$ and $Z$ obeys
\[
Z'(t)=bZ(t)+aZ(t-\tau(t))+\phi(t), \quad t\geq 0
\]  
where $\phi(t)=(1+b)Y(t)+aY(t-\tau(t))$ for $t\geq 0$. Clearly, if $Y(t)\to 0$ as $t\to\infty$ a.s., then $\phi(t)\to 0$ as $t\to\infty$ a.s. Evidently, if $Y(t)\to 0$ as $t\to\infty$ a.s., then $Z(t)\to 0$ as $t\to\infty$ a.s. is a necessary and sufficient condition to ensure that $X(t)\to 0$ as $t\to\infty$ a.s.

Thus, under the condition that the perturbations $f$ and $\sigma$ are ``small'' in such a way that the solutions of the SDE converge to zero, the question of whether deterministic and stochastic perturbations which preserve the convergence of solutions to the equilibrium in the complicated SFDE reduces to whether the deterministically perturbed FDE is convergent if there is a fading perturbation. By virtue of Cermak~\cite{cermak}, it is known that the conditions $b<0$ and $|a|<|b|$ are sufficient for the solution of the unperturbed equation 
\[
u'(t)=bu(t)+au(t-\tau(t)), \quad t\geq 0,
\]
to obey $u(t)\to 0$ as $t\to\infty$. Keeping this condition on $a$ and $b$, we can take the equation for $Z$ and show that $Z$ obeys
\[
D_+|Z(t)|\leq b|Z(t)|+|a||Z(t-\tau(t))|+|\phi(t)|,\quad t\geq 0,
\]
where $D_+$ is the Dini--derivative earlier introduced. Now, by comparison arguments, $|Z(t)| < z(t)$ for $t\geq -\tau$ where $z=|\psi|+1$ on $[-\tau,0]$ and 
\[
z'(t)=bz(t)+|a|z(t-\tau(t))+|\phi(t)|+e^{-t}, \quad t\geq 0.
\]
See Appleby and Buckwar~\cite{AppBuck10}, for example. Finally, it remains to show that $z(t)\to 0$ as $t\to\infty$, and this will give $Z(t)\to 0$ as $t\to\infty$, as needed. 

The proof that $z(t)\to 0$ as $t\to\infty$ can be achieved by means of comparison arguments, showing first that $z$ is bounded, and then by showing that a positive limsup of $z$ leads to a contradiction, and hence to the conclusion that $z(t)\to 0$ as $t\to\infty$. 

The boundedness of $z$ may also be proven by a comparison argument: simply choose a (sufficiently large) constant upper solution to the differential equation for $z$, exploiting the facts that the last terms on the right hand side are uniformly bounded, and that $-b>|a|\geq 0$. The key to the proof of convergence is to view all the terms in $z$, apart from the instantaneous term, as perturbations of the differential equation $v'=bv$,  writing down the variation of constants formula, majorising, and then taking the limsup as $t\to\infty$. From this, it is seen that a positive limsup for $z$ is incompatible with the inequality $-b>|a|\geq 0$.  

Parallel arguments can be developed for the characterisation of bounded solutions. 

Thus, the following theorem can be established.
\begin{theorem}
Let $b<0$, $|a|<|b|$, $\tau\in C([0,\infty);(0,\infty))$ be such that 
$t-\tau(t)\to \infty$ as $t\to\infty$ and $-\bar{\tau}:=\sup_{t\geq 0}(t-\tau(t)) >-\infty$. Let $f$ and $\sigma$ be continuous, and suppose that $\psi\in C([-\bar{\tau},0];\mathbb{R})$. Then there is a unique continuous adapted process $X$ which is the solution of \eqref{eq.sfde} with 
with $X=\psi$ on $[-\bar{\tau},0]$, and moreover:
\begin{enumerate}
	\item[(a)] The following are equivalent:
	\begin{enumerate}
		\item[(A)] $f$ and $\sigma$ obey condition (a)(A) of Theorem \ref{thm.Yto0bddchar};
		\item[(B)] $X(t)\to 0$ as $t\to\infty$ a.s.
	\end{enumerate} 
	\item[(b)] The following are equivalent:
	\begin{enumerate}
		\item[(A)] $f$ and $\sigma$ obey condition (b)(A) of Theorem \ref{thm.Yto0bddchar};
		\item[(B)] $X$ is a.s. bounded. 
	\end{enumerate} 
\end{enumerate}
\end{theorem}
We note that this theorem does not require any smoothness or monotonicity properties on the delay $\tau$, nor on the delayed argument $\text{id}-\tau$. This improves upon the regularity required on $\tau$ in \cite{cermak}, for instance.
\subsection{Decay estimates for general unbounded delay equations}
The question about rates of decay to zero if the solution of \eqref{eq.sfde} demands that a good asymptotic estimate be found to the rate at which $Y(t)\to 0$ as $t\to\infty$. This can then be used to get the rate of decay of $\phi(t)\to 0$ as $t\to\infty$. At this point, it may be hard to use \cite[Theorem 2.3]{cermak} on perturbed equations, since the asymptotic behaviour of $\phi'$ is also needed, and $\phi'$ will not be well--defined for stochastic perturbations, because the non--differentiability of $Y$ will make $\phi$ non--differentiable. There are possible routes to circumvent this problem. One is to invoke a kind of comparison argument in the manner of \cite{AppBuck10} or this paper, but this may not lead to sharp results. Another possibility is to apply a ``double mollification'' to the equation, as follows. Consider $X$, $Y$ and $Z$ as above. Now, define 
$Y_2'(t)=bY_2(t)+\phi(t)$ for $t\geq 0$,
and set $Y_2(t)=0$ for $t\leq 0$. Since $b<0$ and $\phi(t)\to 0$ as $t\to\infty$, we have that $Y_2(t)\to 0$ as $t\to\infty$. Moreover, an upper estimate on the rate of decay of $Y_2$ can be obtained from the rate of decay of $\phi$ (which is known from $Y$). Next, define $Z_2=Z-Y_2$. Then 
\[
Z_2'(t)= bZ_2(t)+aZ_2(t-\tau(t)))+\phi_2(t)
\]
where $\phi_2(t)=aY_2(t-\tau(t))$. Now, provided $\tau$ is continuously differentiable, as it is in Cermak's theory, $\phi_2$ is continuously differentiable, and asymptotic estimates can be obtained for $\phi_2$ and $\phi_2'$ in terms of $Y_2$ and $Y_2'=bY_2+\phi$, all of which can be estimated using pathwise estimates of $Y$. At this point, Cermak's results could be applied off the shelf to give an estimate on $Z_2$. But since $Z=Y_2+Z_2$ and $X=Z+Y$, an estimate on the decay rate of $X$ could be found, again appealing to the decay estimate on $Y$. As a result, we see that a theory that gives good upper bounds for the rate of decay of $X$ can probably be deduced from existing results, provided we can establish good asymptotic estimates for the solution of the SDE. Such a characterisation of the asymptotic behaviour of the SDE will form part of a future work.   

The sharpness of the sufficient conditions generated in this approach could be examined by returning to \eqref{eq.YXgeneral}. If
$X$ has a known a.s. rate of decay, by estimating the terms on the righthand side, an estimate on the necessary rate of decay of $Y$ could be found. 

\subsection{Multidimensional equations with additive noise}
In this subsection, we discuss the multidimensional pantograph equation (rather than equations with general unbounded delay) which is forced both deterministically and stochastically (although, as we mention at the end, most of our argument remains valid for equations with general unbounded delay). We consider $d$--dimensional equations. The analogue of the scalar deterministic perturbation is an $f\in C([0,\infty);\mathbb{R}^d)$ (and we denote its $i$-th component $f_i$). The multidimensional analogue for deterministically perturbed scalar equations is therefore 
\[
x'(t)=Bx(t)+Ax(qt)+f(t), \quad t\geq 0,
\]
where $A$ and $B$ are $d\times d$ real matrices, based on the underlying unforced equation being 
\begin{equation} \label{eq.multidpantou}
u'(t)=Bu(t)+Au(qt), \quad t\geq 0.
\end{equation}
We assume hereinafter that $A$ is not the zero matrix, since in that case all delay differential equations reduce to being  ordinary ones. 

As for the stochastic equation, since there are many dimensions, there may reasonably be many (independent) sources of randomness affecting the dynamics. Focussing on the dynamics of the $i$--th component of the solution vector $X(t)$, if we are to have a state--independent additive noise, we arrive at:
\[
dX_i(t)=\left\{[BX(t)+AX(qt)]_i+f_i(t)\right\}\,dt+\sum_{j=1}^r \sigma_{ij}(t)\,dW_j(t), 
\]
where $W=(W_1,\ldots,W_r)$ is an $r$--dimensional standard Brownian motion (or Wiener process). The $d\times r$ scalar continuous functions $\sigma_{ij}$ can be formed into a $d\times r$ continuous $d\times r$ matrix--values function (its value at time $t$ is denoted $\Sigma(t)$), and then $X$ can be expressed in the differential shorthand 
\begin{equation} \label{eq.stochmultidpanto}
dX(t)=\left\{BX(t)+AX(qt)+f(t)\right\}\,dt+\Sigma(t)\,dW(t), \quad t\geq 0. 
\end{equation}
With $X(0)=\xi\in \mathbb{R}^d$ given, there is a unique continuous adapted process satisfying this equation. 
The question now arises: can we characterise the situation wherein $x(t)\to 0$ as $t\to\infty$ and $X(t)\to 0$ as $t\to\infty$, granted that $u(t)\to 0$ as $t\to\infty$?

As before, we start with a converse result. Let $Y$ be the $d$--dimensional process given by 
\[
dY(t)=\left\{-Y(t)+f(t)\right\}\,dt+\Sigma(t)\,dW(t), \quad t\geq 0,
\] 
with $Y(0)=0\in\mathbb{R}^d$. Then generalising existing arguments to the multidimensional case, we get 
\[
Y(t)= X(t)-e^{-t}X(0)-\int_{0}^t e^{-(t-s)}\{(I_d+B)X(s)+AX(qs)\}\,ds, \quad t\geq 0, 
\] 
where $I_d$ is the $d\times d$ identity matrix. From this, we see that whenever $X(t)\to 0$ as $t\to\infty$ a.s., we have $Y(t)\to 0$ as $t\to\infty$ a.s. This last condition holds provided $Y_i(t)\to 0$ as $t\to\infty$ for each $i$ a.s., and noticing that $Y_i$ obeys the scalar SDE 
\[
dY_i(t)=\{-Y_i(t)+f_i(t)\}\,dt + \sum_{j=1}^r \sigma_{ij}(t)\,dW_j(t),
\]
we see that the equation satisfied by $Y_i$ is of essentially the same form as that satisfied by the solution of \eqref{eq.Y}. In fact, if we define 
\begin{equation} \label{eq.sigmai}
\sigma_i(t):=\left(\sum_{j=1}^r \sigma_{ij}^2(t)\right)^{1/2}, \quad t\geq 0,
\end{equation}
we may use the martingale representation theorem (see e.g., \cite{KS}) to rewrite the continuous martingale 
\[
M_i(t):=\sum_{j=1}^r \int_{0}^t \sigma_{ij}(s)\,dW_j(s), \quad t\geq 0,
\]
as 
\[
M_i(t)=\int_0^t \sigma_i(s)\,d\tilde{B}_i(s), \quad t\geq 0,
\]
where each $\tilde{B}_i$ is a standard one--dimensional Brownian motion (see e.g., Karatzas and Shreve~\cite{KS}). Therefore, $Y_i$ obeys the SDE
\[
dY_i(t)=\{-Y_i(t)+f_i(t)\}\,dt + \sigma_i(t)\,d\tilde{B}_i(t),
\]
which is in precisely the form of \eqref{eq.Y}. Now, define 
$S_i(\epsilon)$ as $S(\epsilon)$ in \eqref{eq.S}, but with $\sigma_i$ in the role of $\sigma$. Since we require $Y_i(t)\to 0$ as $t\to\infty$ for each $i$, a.s., it follows from the converse result (i.e. (B) implies (A)) of part (a) of  Theorem~\ref{thm.Yto0bddchar}, applied to each $Y_i$, that $S_i(\epsilon)<+\infty$ for all $\epsilon>0$ and that for each $f_i$ and each $\theta>0$ the condition 
\[
\int_{t-\theta}^t f_i(s)\,ds \to 0 \quad t\to\infty
\]
holds. This latter condition is nothing but 
\begin{equation} \label{eq.multidimf}
\int_{t-\theta}^t f(s)\,ds\to 0, \quad t\to\infty, \quad \text{for all $\theta>0$}.
\end{equation}
Therefore this condition and $S_i(\epsilon)<+\infty$ for all $\epsilon>0$ and $i\in \{1,\ldots,d\}$ is \textit{necessary} in order that $X(t)\to 0$ as $t\to\infty$ a.s.

Conversely, if each $S_i(\epsilon)<+\infty$ for all $\epsilon>0$ and $i\in \{1,\ldots,d\}$ and $\int_{t-\theta}^t f(s)\,ds\to 0$ as $t\to\infty$, then $Y_i(t)\to 0$ as $t\to\infty$ a.s. for each $i$, and therefore $Y(t)\to 0$ as $t\to\infty$ a.s. As before, set $Z(t):=X(t)-Y(t)$ for $t\geq 0$, and let $\phi(t)=(I_d+B)Y(t)+AY(qt)$ for $t\geq 0$. Then $Z(0)=X(0)$ and $Z$ obeys the perturbed equation 
\begin{equation} \label{eq.pantomultidim}
Z'(t)=BZ(t)+AZ(qt)+\phi(t), \quad t\geq 0.
\end{equation}
Moreover, since $Y(t)\to 0$ as $t\to\infty$ a.s., we have that $\phi(t)\to 0$ as $t\to\infty$ a.s., and since the paths of $Y$ are continuous, $\phi$ is continuous on $[0,\infty)$. We notice that if $Z(t)\to 0$ as $t\to\infty$ a.s., then $X(t)\to 0$ as $t\to\infty$ a.s. 

Therefore, by considering \eqref{eq.pantomultidim} path--by--path (i.e., for each $\omega$ in the a.s. event where $X$ and $Y$ are well--defined, and $Y$ tends to zero), it can be seen that if the \textit{deterministic} property 
\begin{equation} \label{eq.perturbedpropertymultid}
\text{For every $\phi\in C_0([0,\infty);\mathbb{R}^d)$, the solution $Z$ of \eqref{eq.pantomultidim} obeys $Z(t)\to 0$ as $t\to\infty$}
\end{equation}
holds,  
then we have that $Y(t)\to 0$ as $t\to\infty$ a.s. implies $X(t)\to 0$ as $t\to\infty$ a.s. 

We have therefore proven the following result, \textit{which eliminates the need for further analysis of stochastically perturbed equations}, and reduces the study of \eqref{eq.stochmultidpanto} to that of  \eqref{eq.pantomultidim} with deterministic forcing term $\phi\in C_0([0,\infty);\mathbb{R}^d)
$.

\begin{theorem} \label{thm.multi1}
Let $f\in C([0,\infty);\mathbb{R}^d)$ and $\Sigma\in C([0,\infty);\mathbb{R}^{d\times r})$. With $\sigma_{ij}=\Sigma_{ij}$,  define $\sigma_i$ for each $i\in \{1,\ldots,d\}$ according to \eqref{eq.sigmai}. 
Suppose that the property \eqref{eq.perturbedpropertymultid} holds for solutions of \eqref{eq.pantomultidim}. 
Let $X$ be the unique continuous adapted process which obeys \eqref{eq.stochmultidpanto}.
The following are then equivalent:
\begin{enumerate}
	\item[(A)] $f$ obeys \eqref{eq.multidimf} and for each $i$ and $\epsilon>0$ 	\begin{equation} \label{eq.Si}
		S_i(\epsilon):= \sum_{n=1}^\infty \sqrt{\int_{n-1}^n\sigma_i^2(s)\,ds} \exp\left(-\epsilon\int_{n-1}^{n}\sigma_i^2(s)\,ds\right)<+\infty.
	\end{equation}
	\item[(B)] $X(t)\to 0$ as $t\to\infty$ a.s.
\end{enumerate}
\end{theorem} 
Of course, the theorem would be vacuous if property \eqref{eq.perturbedpropertymultid} does not hold. We now explore the property \eqref{eq.perturbedpropertymultid} a little further, and show that it holds under conditions which are, in a certain sense, reasonable multidimensional analogues of \eqref{eq.unforcedasstable}. Since \eqref{eq.perturbedpropertymultid} must hold in the case $\phi(t)=0$ for all $t\geq 0$, this means the all solution $u$ of the unperturbed pantograph equation 
\eqref{eq.multidpantou} must obey $u(t)\to 0$ as $t\to\infty$. Iserles has shown in \cite{Iser} that this is equivalent to 
\begin{equation} \label{eq.iser1}
\text{All eigenvalues of $B$ have negative real parts}
\end{equation}
(a condition which implies $B$ is invertible), together with the spectral radius condition 
\begin{equation} \label{eq.iser2}
\rho(B^{-1}A)<1.
\end{equation}
Therefore, the conditions \eqref{eq.iser1} and \eqref{eq.iser2} should be imposed on $A$ and $B$. However, to the best of our knowledge, no result exists in the literature which states that \eqref{eq.iser1} and \eqref{eq.iser2} are sufficient to ensure that property \eqref{eq.perturbedpropertymultid} holds. It would clearly be of value to prove such a result, since, if it holds, a complete characterisation of convergence of solutions of  perturbed stochastic and deterministic pantograph equations would be achieved. A cursory reading of the work of Iserles in \cite{Iser}, however, does not lead to optimism that his approach could fulfill this easily. The proofs in that paper rely on analyticity of the solutions of the unperturbed equation, and so--called Dirichlet series solutions that can be formed as a consequence of the analyticity. For stochastic equations at least, the forcing function $\phi$ is continuous, but nowhere differentiable, so $Z$ is a $C^1$ function, but not even $C^2$, for instance.     

We show next, however, that for a subclass of $A$ and $B$ obeying properties \eqref{eq.iser1} and \eqref{eq.iser2}, that property \eqref{eq.perturbedpropertymultid} does hold. We retain property \eqref{eq.iser1}, and note that \eqref{eq.iser2} in a certain sense demands that $A$ be small relative to $B$. Our sufficient condition imposes such a smallness condition on $A$ in lieu of \eqref{eq.iser2}. 

It is worth noting that although our argument is not sharp in the pantograph case, it works without modification for general delay; this is appealing, because, to the best of our knowledge, necessary and sufficient conditions such as \eqref{eq.iser1} and \eqref{eq.iser2} are not known which characterise the stability of unperturbed multidimensional equations with unbounded delay. 

To prove property \eqref{eq.perturbedpropertymultid}, we continue to assume \eqref{eq.iser1}; therefore  it follows that there is a unique symmetric, positive definite $d\times d$ matrix $Q$ (with real entries) such that 
\begin{equation} \label{eq.liapB}
B^TQ+QB=-I_d,
\end{equation}
where, here and in the sequel, the superscripted $T$ denotes the transpose of a matrix. Now consider $z(t):=Z(t)^TQZ(t)$ for $t\geq 0$. We see that $z(t)\in \mathbb{R}$ for each $t\geq 0$. Furthermore, since $Q$ is positive definite, $z(t)\geq 0$ for all $t\geq 0$. Our strategy is to derive a perturbed scalar pantograph differential inequality for $z$. Using \eqref{eq.liapB} we get 
\begin{align}
z'(t)%&=Z'(t)^T QZ(t)+Z(t)^TQZ'(t)\\
%&=(BZ(t)+AZ(qt)+\phi(t))^TQZ(t)
%+Z(t)^TQ %(BZ(t)+AZ(qt)+\phi(t))\\
&=(Z(t)^TB^T+Z(qt)^TA^T+\phi(t)^T)QZ(t)
\nonumber\\
&\qquad+Z(t)^T  (QBZ(t)+QAZ(qt)+Q\phi(t))\nonumber\\
\label{eq.zDDEmaster}
&=-\|Z(t)\|^2+2Z(t)^TQAZ(qt) +2\phi(t)^TQZ(t), \quad t\geq 0.
\end{align}
Notice next that the positive definiteness of $Q$ means that 
$c_1\|Z(t)\|^2\leq z(t)\leq c_2\|Z(t)\|^2$ for some $0<c_1\leq c_2$, where $\|\cdot\|$ is the Euclidean norm on $\mathbb{R}^d$ (in fact, $c_1$ and $c_2$ are the minimum and maximum eigenvalues of $Q$, respectively).

Now we estimate the terms on the righthand side of 
\eqref{eq.zDDEmaster}. 
We start with the last term. By the Cauchy--Schwarz inequality, we get
\[
|2\phi(t)^TQZ(t)|
\leq 2\|\phi(t)\|\|QZ(t)\|
\leq \|\phi(t)\|\|Q\|\|Z(t)\|=:2\epsilon_1(t)\|Z(t)\|,
\]
where $\epsilon_1(t):=\|\phi(t)\|\|Q\|$ is continuous, nonnegative and tends to zero as $t\to\infty$, and for any $d\times d$ matrix $M$, $\|M\|$ denotes the matrix norm of $M$ induced from the Euclidean norm.    
Now, recall that $2xy\leq x^2+y^2$ for all $x,y\geq 0$, so, putting $x=\sqrt{\epsilon_1(t)}$ and $y=\sqrt{\epsilon_1(t)}\|Z(t)\|$, we get 
\begin{equation} \label{eq.estt3}
|2\phi(t)^TQZ(t)|\leq
2\epsilon_1(t)\|Z(t)\|
\leq \epsilon_1(t)+\epsilon_1(t)\|Z(t)\|^2.
\end{equation}
For the second term on the righthand side of \eqref{eq.zDDEmaster}, we have 
$
2Z(t)^TQAZ(qt)=2Z(t)^TQ^TAZ(qt)=2 (A^TQZ(t))^T Z(qt)$,
so, setting $\beta:=\|A^TQ\|>0$ and using the  Cauchy--Schwarz inequality once again, we have 
\begin{align*}
|2Z(t)^TQAZ(qt)|&\leq 2\|A^TQZ(t)\|\|Z(qt)\|\leq 2\beta\|Z(t)\|\|Z(qt)\|\\
&\leq \beta\left(\alpha\|Z(t)\|^2+\frac{1}{\alpha}\|Z(qt)\|^2\right),
\end{align*}
% Employing these estimates, we have 
%\[
%z'(t)\leq -\|Z(t)\|^2 + 2\|A^TQ\|\|Z(t)\|\|Z(qt)\| + \epsilon_1(t)\|Z(t)\|^2+\epsilon_1(t), \quad t \geq 0.
%\]
where, for arbitrary $\alpha>0$ we have used the inequality $2xy\leq \alpha x^2+y^2/\alpha$ for $x, y\geq 0$ at the last step. 

Hence, using this estimate and \eqref{eq.estt3}, we may return to \eqref{eq.zDDEmaster} to get  
\[
z'(t)\leq -\|Z(t)\|^2+\beta\alpha\|Z(t)\|^2+\frac{\beta}{\alpha}\|Z(qt)\|^2+\epsilon_1(t)\|Z(t)\|^2 + \epsilon_1(t).
\]
Next, we note that $\beta>0$; to see this, assume to the contrary that $\beta=0$. Then $A^TQ=0$. Since $Q$ is positive definite, $Q$ is invertible. This forces $A=0$, which we have ruled out by hypothesis, and we have the desired contradiction. Thus, we have that $\beta>0$ and we may    pick $\alpha:=1/(2\beta)>0$. Then 
\[
z'(t)\leq -\frac{1}{2}\|Z(t)\|^2
+2\beta^2\|Z(qt)\|^2+\epsilon_1(t)\|Z(t)\|^2 + \epsilon_1(t),\quad t\geq 0.
\]
Since $\epsilon_1(t)\to 0$ as $t\to\infty$, for every $\epsilon\in (0,1/2)$ there is $T(\epsilon)>0$ such that for $t\geq T(\epsilon)$, $\epsilon_1(t)<\epsilon$. Hence for $t\geq T(\epsilon)$, we have 
\[
z'(t)\leq -\left(\frac{1}{2}-\epsilon\right)\|Z(t)\|^2
+2\beta^2\|Z(qt)\|^2+ \epsilon_1(t).
\]
Finally, $\|Z(qt)\|^2\leq c_1^{-1}z(qt)$ and $\|Z(t)\|^2\geq c_2^{-1}z(t)$, so for $\epsilon\in (0,1/2)$ we have 
\begin{equation} \label{eq.zpantoineq}
z'(t)\leq -\left(\frac{1}{2}-\epsilon\right)\frac{1}{c_2}z(t)
+\frac{2\beta^2}{c_1}z(qt)+ \epsilon_1(t), \quad t\geq T(\epsilon).
\end{equation}
Since for any $d\times d$ matrix $M$ the identity $\|M\|=\|M^T\|$ holds for the 2--norm, and $Q$ is symmetric, we have $\beta=\|Q^TA\|=\|QA\|$. Now, suppose that 
\begin{equation} \label{eq.stabcondn2}
4\|QA\|^2< \frac{c_1}{c_2},
\end{equation}
or equivalently, $2\beta^2c_2/c_1<1/2$. Choose 
$\epsilon\in (0,1/2-2\beta^2 c_2/c_1)$. Then the coefficient of $z(t)$ in \eqref{eq.zpantoineq} is negative, and in modulus exceeds the coefficient of $z(qt)$. Since $\epsilon_1(t)\to 0$ as $t\to\infty$, we may apply our scalar theory to the inequality \eqref{eq.zpantoineq} to deduce that  $z(t)\to 0$ as $t\to\infty$. But since $\|Z(t)\|^2\leq c_1^{-1}z(t)$, this implies $Z(t)\to 0$ as $t\to\infty$. 

Thus we see the sufficient conditions required to ensure property \eqref{eq.perturbedpropertymultid} are \eqref{eq.iser1} together with \eqref{eq.stabcondn2} in place of \eqref{eq.iser2}. We discuss now the condition \eqref{eq.stabcondn2}, and show that it retains much of the spirit of the condition \eqref{eq.iser2}.

The righthand side of \eqref{eq.stabcondn2} is the minimal eigenvalue of $Q$, divided by the maximal eigenvalue of $Q$, so is a quantity which depends on $B$, but not on $A$. Therefore, as earlier claimed, \eqref{eq.stabcondn2} can be interpreted as a smallness condition on $A$, similar to \eqref{eq.iser2}, because, using the submultiplicative property of the matrix $2$--norm,
and the fact that  $\|Q\|=c_2$ for this norm,  we see that 
\[
4\|A\|^2< \frac{c_1}{c_2^3}
\]
implies \eqref{eq.stabcondn2}. Note also that the condition \eqref{eq.stabcondn2} is ``delay independent'' in the sense that \eqref{eq.stabcondn2} does not depend on $q$; it also shares this property with the condition \eqref{eq.iser2}. 

Using the fact that $Q=\int_0^\infty e^{Bs}e^{B^Ts}\,ds$ and $\int_0^\infty e^{2Bs}\,ds = -\frac{1}{2}B^{-1}$, we see that we can rewrite \eqref{eq.iser2} and \eqref{eq.stabcondn2} as 
\[
\rho\left(\int_0^\infty e^{2Bs}\,ds A \right)<\frac{1}{2},
\quad
\left\|\int_0^\infty e^{Bs} e^{B^Ts}\,ds A\right\|<\frac{1}{2}\sqrt{\frac{\text{minimal eigenvalue of $Q$}}{\text{maximal eigenvalue of $Q$}}},
\]
respectively. Noting that $\rho(M)\leq \|M\|$ for any matrix $M$, we see the close connection between the conditions.

To see further similarities between \eqref{eq.stabcondn2} and \eqref{eq.iser2}, note in the one--dimensional case that we $b:=B<0$ and $2BQ=-1$. Since $c_1/c_2=1$, the condition \eqref{eq.stabcondn2} is equivalent to  $1=\sqrt{c_1/c_2}>2\|QA\|=2QA=-A/B=-a/b$, and hence in the case $d=1$, \eqref{eq.iser1} and \eqref{eq.stabcondn2} are exactly \eqref{eq.unforcedasstable}, which is exactly \eqref{eq.iser1} and \eqref{eq.iser2} in the case $d=1$. 

One other connection between \eqref{eq.iser2} and \eqref{eq.stabcondn2} is worth pointing out. If $B$ is real diagonal, then $Q=-\frac{1}{2}B^{-1}$. Suppose $A$ is symmetric. Then $QA=B^{-1}A$ is symmetric, and so $\|QA\|=\rho(QA)=\rho(B^{-1}A)$ and \eqref{eq.stabcondn2} becomes 
\[
\rho(B^{-1}A) < \frac{\text{minimal eigenvalue of $|B|$}}{\text{maximal eigenvalue of $|B|$}},
\]
a condition very near in form to,  but more restrictive than, condition \eqref{eq.iser2}.

The above calculations therefore lead to the following theorem, which is a corollary of Theorem~\ref{thm.multi1}.
\begin{theorem}  \label{thm.multi2}
Let $f\in C([0,\infty);\mathbb{R}^d)$ and $\Sigma\in C([0,\infty);\mathbb{R}^{d\times r})$. With $\sigma_{ij}=\Sigma_{ij}$,  define $\sigma_i$ for each $i\in \{1,\ldots,d\}$ according to \eqref{eq.sigmai}. 
Suppose that \eqref{eq.iser1} holds, and with $Q$ defined by \eqref{eq.liapB}, the condition \eqref{eq.stabcondn2} holds. 
Let $X$ be the unique continuous adapted process which obeys \eqref{eq.stochmultidpanto}.
Then the following are equivalent:
\begin{enumerate}
	\item[(A)] $f$ obeys \eqref{eq.multidimf} and for each $i$ and $\epsilon>0$ 	\begin{equation} \label{eq.Si}
		S_i(\epsilon):= \sum_{n=1}^\infty \sqrt{\int_{n-1}^n\sigma_i^2(s)\,ds} \exp\left(-\epsilon\int_{n-1}^{n}\sigma_i^2(s)\,ds\right)<+\infty.
	\end{equation}
	\item[(B)] $X(t)\to 0$ as $t\to\infty$ a.s.
\end{enumerate}
\end{theorem} 
As stated above, we believe an optimal result would replace the sufficient condition \eqref{eq.stabcondn2} with \eqref{eq.iser2}. 

We finish with the remarks that the sketched results in this section do not really depend on proportional delay, and therefore should hold for any multidimensional equation with unbounded delay, and also that a result characterising boundedness can be formulated and proved with equal ease. 

\subsection{Equations with state--dependent noise}

There are many studies already present concerning the asymptotic behaviour of stochastic pantograph equations with multiplicative noise, or more generally, state--dependent noise. To date, the results on stability, growth bounds etc tend to give upper estimates, and sufficient conditions for asymptotic stability. The question arises, therefore, whether the techniques of this paper might be helpful in obtaining more precise asymptotic results. 

A direct application of the technique is unlikely to be successful in dealing with equations of the form 
\begin{equation} \label{eq.multipanto}
dX(t)=(bX(t)+aX(qt))\,dt+\sigma X(t)\,dB(t)
\end{equation}
since there is no obvious small perturbation to mollify. However, part of the sharpness achieved in this paper is due to 
a successful decomposition of non--differentiable dynamics of the delay--differential equation into an ordinary non--smooth stochastic differential equation, and a pathwise differentiable stochastic delay differential equation, so an adoption of the that aspect of the philosophy of this paper may be useful. We show briefly here how, in the special case that $a>0$, solutions of \eqref{eq.multipanto} have the following asymptotic behaviour. 
\begin{theorem}
Define $\lambda=b-\sigma^2/2$. Let $a>0$ and $X(0)>0$. Let $X$ be the unique strong (positive) solution of \eqref{eq.multipanto}.
\begin{itemize}
	\item[(a)] If $\lambda>0$, then 
	\[
	\lim_{t\to\infty} \frac{1}{t}\log X(t)=\lambda,\quad \text{a.s.}
	\]
	\item[(b)] If $\lambda\leq 0$, then 
	\[
	\lim_{t\to\infty} \frac{1}{t}\log X(t)=0,\quad \text{a.s.}
	\]
\end{itemize}
\end{theorem}
Therefore, the situation is similar to the case for the deterministic pantograph equation (wherein $\sigma=0$), with subexponential asymptotic behaviour when the underlying instantaneous equation 
\[
dv(t)=bv(t)\,dt+\sigma v(t)\,dB(t)
\]
is asymptotically stable (which occurs when $\lambda<0$). When the underlying instantaneous equation is unstable however (when $\lambda>0$), exponential growth occurs, and does so at a rate $\lambda=b-\sigma^2/2>0$.

This result can be established by means of the following multiplicative decomposition. Let 
$\rho$ be given by $\rho(0)=1$ and 
\[
d\rho(t)=b\rho(t)\,dt + \sigma \rho(t)\,dB(t), \quad t\geq 0.
\]
Then $\rho$, like $Y$, is the solution of a linear SDE, and has an explicit solution 
\[
\rho(t) = e^{(b-\sigma^2/2)t+\sigma B(t)}, \quad t\geq 0.
\]
Since $\rho$ is strictly positive, $1/\rho$ is a well--defined process and obeys the SDE
\[
d\left(\frac{1}{\rho(t)}\right) = \frac{\sigma^2-b}{\rho(t)}\,dt -\frac{\sigma}{\rho(t)}\,dB(t), \quad t\geq 0.
\]
The process $Z=X/\rho$ is likewise 
well--defined, and by stochastic integration by parts obeys 
\[
dZ(t)=X(t)d\left(\frac{1}{\rho(t)}\right)+\frac{1}{\rho(t)}dX(t)+\sigma X(t)\cdot -\frac{\sigma}{\rho(t)}\,dt,
\]
which simplifies to 
\[
dZ(t)=\frac{1}{\rho(t)}aX(qt)\,dt,\quad t\geq 0.
\]
Therefore, we have 
\[
Z(t)=X(0)+\int_0^t a Z(qs) \frac{\rho(qs)}{\rho(s)}\,ds, \quad t\geq 0.
\]
Since the integrand is continuous, $Z$ is in fact differentiable, and we have 
\begin{equation} \label{eq.Zstochapuredelay}
Z'(t)= a \frac{\rho(qt)}{\rho(t)}Z(qt)=:a(t)Z(qt), \quad t\geq 0,
\end{equation}
where $\lambda=b-\sigma^2/2$ and $a(t)=a  e^{\lambda (q - 1) t+\sigma(B(qt)-B(t))}$. 
The strategy now is to study the dynamics of $Z$ pathwise, exploiting the fact that $a(t)$ (or equally $\rho$) is known explicitly in terms of the Brownian motion $B$. Finally, we can recover information about $X$ from $X(t)=\rho(t)Z(t)$, again exploiting knowledge of $\rho$. 

When $\lambda>0$, this technique will work well, since in that case, due to the Strong Law of Large Numbers for standard Brownian motion, 
\[
\lim_{t\to\infty} \frac{1}{t}\log |a(t)| = -\lambda(1-q)<0, \quad \text{a.s.},
\]
so 
$a\in L^1(0,\infty)$ a.s. Then, arguing pathwise, $Z$ tends to a finite limit a.s. (call it $Z^\ast$). Therefore 
\[
X(t)\sim Z^\ast \rho(t) = Z^\ast e^{\lambda t+\sigma B(t)}, \quad  t\to\infty, \quad \text{a.s.}, 
\]
so $X$ grows exponentially (under the assumption that the random limit $Z^\ast$ is non--trivial). We have $Z^\ast>0$ in the case that $X(0)>0$ and $a>0$, where $Z(t)\uparrow Z^\ast\in (0,\infty)$ as $t\to\infty$. 

In the case when $Z^\ast=0$, $Z$ tends to  limit $Z^\ast$ according to $\limsup_{t\to\infty} \log|Z(t)|/t\leq -\lambda$. This can be proved by obtaining an a priori exponential rate of decay on $Z$ from \eqref{eq.Zstochapuredelay}, and then recursively estimating the rate of decay by feeding the exponential estimate back into \eqref{eq.Zstochapuredelay}. From the final estimate on $Z$, we end up with $\limsup_{t\to\infty} \log|X(t)|/t\leq 0$. 

Therefore, in the general case (without sign  restrictions on $a$ and $X(0)$), we have a rather complete picture of the long--run dynamics when $\lambda>0$. Either $Z^\ast\neq 0$ and $X$ inherits the asymptotic behaviour of the $\rho$ (and does so rapidly, because $Z$ converges exponentially fast to $Z^\ast$) or $Z^\ast =0$, and $X$ grows more slowly than any exponential function.  

The case $\lambda<0$ is more analogous to the situation covered in this paper, since when $\sigma=0$, $\lambda=b<0$. In that situation, $a(t)$ still has exponential asymptotic behaviour, and for every $\epsilon\in (0,1-q)$ there is a random variable $T(\epsilon)>0$ such that 
\[
D_+|Z(t)|\leq \left(\frac{-\lambda(1-q-\epsilon)}{1-q}+\epsilon\right)e^{-\lambda(1-q-\epsilon)t}|Z(qt)|, \quad t\geq T(\epsilon).
\]
Now, using comparison techniques, one can show by choosing $\bar{Z}$ sufficiently large (take $\bar{Z}>\max_{s\in [qT,T]} |Z(s)|$) that $|Z(t)|\leq \bar{Z} e^{-\lambda(1-q-\epsilon)/(1-q)t}=:Z_+(t)$ for all $t\geq T$. This is done by showing the $Z_+$ obeys the upper differential inequality satisfied by $|Z|$. From this we rapidly deduce that 
\[
\limsup_{t\to\infty} \frac{1}{t}\log |Z(t)|\leq -\lambda, \quad \text{a.s.}
\]
and since $X(t)=\rho(t)Z(t)$ and 
$\lim_{t\to\infty} \log \rho(t)/t=\lambda$ a.s.
we have that 
\[
\limsup_{t\to\infty} \frac{1}{t}\log |X(t)|\leq 0, \quad \text{a.s.}
\]
so that growth of solutions cannot be exponentially fast. 

In the case when $a>0$ and $X(0)>0$, we can strengthen the result to 
\[
\lim_{t\to\infty} \frac{1}{t}\log|X(t)|=0, \quad\text{a.s.},
\]
so the subexponential dynamics $u$ are inherited by $X$. This can be done by building a lower bound for $Z(t)$ by analogous comparison techniques to those outlined above (this works well, because $Z$ is positive, and indeed, increasing), and prove that 
\[
\liminf_{t\to\infty} \frac{1}{t}\log Z(t)\geq -\lambda, \quad \text{a.s.},
\]
which leads to 
\[
\liminf_{t\to\infty} \frac{1}{t}\log X(t)\geq 0, \quad \text{a.s.}
\]
Combining this with the limsup, we get the claimed zero Liapunov exponent. 

The argument in the case when $\lambda=0$ is similar to that when $\lambda<0$. 


%The methods of this paper may also work in the case where the noise term is small and the underlying unforced equation is stable i.e., for an equation of the form 
%\[
%dX(t)=(bX(t)+aX(qt))\,dt+\sigma(X(t))\,dB(t)
%\]
%where $\sigma(x)=o(x)$ as $x\to 0$ and $b<0$, $|a|<|b|$. What is needed now is an a priori bound on $X$. Using the Tanaka--Meyer formula, and by establishing that the local time in zero of $X$ is zero, one can show that $x(t):=\mathbb{E}[|X(t)|]$ obeys 
%\[
%D_+x(t)\leq bx(t)+|a|x(qt), \quad t\geq 0,
%\] 
%from which we see that $x(t)=O(t^\kappa)$. From this, it is possible to use standard techniques to show that $X(t)=O(t^(\kappa+\epsilon))$. Now, introduce the process 
%\[
%dY(t)=-Y(t)\,dt+\sigma(X(t))\,dB(t), \quad t\geq 0.
%\]
%direct approach 


\section{Proofs}
\begin{proof}[Proof of Lemma~\ref{lemma.fundamental}]
The proof up to \eqref{eq.Knexplicitestimate} follows closely the line of \cite[Theorem 3(i)]{KatoMc}.
By hypothesis, we rule out the case $a=0$, since this reduces the equation to an ODE. With $a\neq 0$, there exists solutions $k\in \mathbb{C}$ to the equation $aq^k+b=0$. Let $k_0$ be any solution of this equation. Then $\kappa=\Re(k_0)$ is uniquely defined and given by $\kappa=-\log(|b/a|)/\log(1/q)$. Let $c=\log q<0$. Then $e^c=q$. Define $w:(-\infty,\infty)\to\mathbb{C}$ by 
$w(s):=e^{-sk_0} z(e^s)$ for $s\in \mathbb{R}$. Then 
\begin{equation} \label{eq.ws+c}
	w(s+c)=e^{-sk_0} z(qe^s)/q^{k_0},
\end{equation}
and since $z$ is differentiable on $(0,\infty)$, for all $s\in\mathbb{R}$, $w'(s)$ exists and is given by 
\begin{align*}
	w'(s)%&=e^{-sk_0}z'(e^s)e^s-k_0e^{-sk_0}z(e^s) \\
	%	&= e^s\{bw(s)+ae^{-sk_0}z(qe^s)+e^{-sk_0}\varphi(e^s)\}-k_0w(s)\\
	&= e^s b(w(s)-w(s+c)) + e^{s-sk_0}\varphi(e^s)-k_0w(s),
\end{align*}
using \eqref{eq.ws+c} and $aq^{k_0}+b=0$. 
%at the last step. 
Gathering together the $w(s)$ terms on the righthand side and treating the equation as a perturbed linear ODE, we are led to 
\begin{equation} \label{eq.wvop}
	\frac{d}{ds}\left(w(s)e^{k_0s - be^s}\right)
	=-be^se^{k_0s-be^s}w(s+c)+e^{-be^s}e^s\varphi(e^s).
\end{equation}
Next, for $n\in \mathbb{N}$, define $I_n=[s_0-nc,s_0-(n+1)c]$ (note $c<0$), $\tau_n:=s_0-nc$ and $K_n:=\sup_{s\in I_n} |w(s)|$. 
%Note that $|w(s)|e^{s\kappa} = |e^{-sk_0} z(e^s)|e^{s\kappa} =  e^{-s\Re(k_0)} |z(e^s)| e^{s\kappa} = |z(e^s)|$, so the definition of $K_n$ matches that in \eqref{eq.Kn}. 
Now, let $s\in I_{n+1}$ and integrate \eqref{eq.wvop} over $[s_0-(n+1)c,s]$:
\[%\]begin{multline*}
w(t)e^{k_0 t - be^t}\big|_{t=\tau_{n+1}}^s
=-b\int_{\tau_{n+1}}^s e^t e^{k_0t-be^t}w(t+c)\,dt + \int_{\tau_{n+1}}^s e^{-be^t}e^t\varphi(e^t)\,dt.	
\]%end{multline*}
Now, since $b<0$, $-b=|b|>0$ and $s\in I_{n+1}$, taking the complex modulus on both sides, and using the triangle inequality, we get 
\begin{multline*}
	|w(s)||e^{k_0 s}| e^{|b|e^s} 
	\leq |w(s_0-(n+1)c)||e^{k_0(s_0-(n+1)c)}| e^{|b|e^{s_0-(n+1)c}}
	\\+|b|\int_{s_0-(n+1)c}^s e^t |e^{k_0t}|e^{|b|e^t}|w(t+c)|\,dt
	+\int_{s_0-(n+1)c}^s e^{|b|e^t}e^t|\varphi(e^t)|\,dt.	
\end{multline*}
Since $\text{Re}(k_0)=\kappa$, we get 
\begin{multline*}
	|w(s)| e^{\kappa s+|b|e^s} 
	\leq |w(-(n+1)c)|e^{\kappa(s_0-(n+1)c)} e^{|b|e^{s_0-(n+1)c}}
	\\+|b|\int_{s_0-(n+1)c}^s e^t e^{\kappa t+|b|e^t}|w(t+c)|\,dt
	+\int_{s_0-(n+1)c}^s e^{|b|e^t}e^t|\varphi(e^t)|\,dt.	
\end{multline*}
Since the terms in $w$ are bounded by $K_n$, and the term in $\varphi$ by $\epsilon_{n+1}$, we get 
\begin{multline*}
	|w(s)| e^{\kappa s+|b|e^s} 
	\leq K_n \left(e^{\kappa(s_0-(n+1)c)} e^{|b|e^{s_0-(n+1)c}}
	+\int_{s_0-(n+1)c}^s |b|e^t e^{\kappa t+|b|e^t}\,dt\right)
	\\+|b|\int_{s_0-(n+1)c}^s e^{|b|e^t}e^t\,dt \cdot \frac{\epsilon_{n+1}}{|b|}.	
\end{multline*}
Now evaluate the last integral, and bound it by the antiderivative at the upper limit: 
\begin{multline} \label{eq.est1}
	|w(s)| e^{\kappa s+|b|e^s} 
	\leq K_n \left(e^{\kappa(s_0-(n+1)c)} e^{|b|e^{s_0-(n+1)c}}
	-\int_{s_0-(n+1)c}^s be^t e^{\kappa t+|b|e^t}\,dt\right)
	\\+\frac{\epsilon_{n+1}}{|b|} e^{|b|s}.	
\end{multline}
Next, notice that we may take $s_0$ so large that $\kappa-k+|b|/2 e^{s_0-kc}>0$ for all integers $k\geq 0$, which makes $t\mapsto \kappa t+|b|e^t$ monotone on each $I_n$, for any $n$. We focus on the integral remaining on the righthand side in \eqref{eq.est1}. Integrating by parts gives
\begin{multline*}
	\int_{s_0-(n+1)c}^s b e^t e^{\kappa t+|b|e^t}\,dt
	=\frac{be^t}{\kappa-be^t} e^{\kappa t+|b|e^t}\Big|_{s_0-(n+1)c}^s
	\\+\text{O}\left(e^{\kappa s+|b|e^a}\int_{s_0-(n+1)c}^s \left|\frac{d}{ds} \frac{be^t}{\kappa-be^t}\right|\,dt\right),
\end{multline*}
where the big--O term is uniform in $s$, $\kappa$, $s_0$ and $n$, provided $s\in I_{n+1}$ and the stipulations on $s_0$ hold. The integral can therefore be further estimated according to 
\begin{multline*}
	-\left(1-\frac{\kappa}{be^t}\right)^{-1} e^{\kappa t+|b|e^t}\Big|_{s_0-(n+1)c}^s
	+O\left(\kappa e^{-s_0+(n+1)c} e^{\kappa s+|b|e^s}\right)
	\\= -e^{\kappa t+|b|e^t}\Big|_{s_0-(n+1)c}^s
	+O\left(\kappa e^{-s_0+(n+1)c} e^{\kappa s+|b|e^s}\right).
\end{multline*}
Now, we return to \eqref{eq.est1} with this estimate to get
\begin{multline*} %\label{eq.est1}
	|w(s)| e^{\kappa s+|b|e^s} 
	\leq \frac{\epsilon_{n+1}}{|b|} e^{|b|s} \\
	+ K_n \left(e^{\kappa(s_0-(n+1)c)} e^{|b|e^{s_0-(n+1)c}}
	+e^{\kappa t+|b|e^t}\Big|_{s_0-(n+1)c}^s
	+O\left(\kappa e^{-s_0+(n+1)c} e^{\kappa s+|b|e^s}\right)\right).
	%	\\+\frac{\epsilon_{n+1}}{|b|} e^{|b|s}.	
\end{multline*}
Thus for $s\in I_{n+1}$ we have 
$|w(s)| \leq K_n \left\{1+O\left(\kappa e^{-s_0+(n+1)c} \right)\right\}
+\epsilon_{n+1} e^{-\kappa s}/|b|$.	Taking the sup on both sides over $I_{n+1}$, and noting that $\kappa<0$ yields
\[
K_{n+1}\leq K_n(1+A e^{nc})+\frac{\epsilon_{n+1}}{|b|} e^{-\kappa (s_0-(n+2)c)}
=K_n(1+A e^{(n+1)c})+B'\epsilon_{n+1}e^{\kappa c n},
\]
where $A>0$ %can be made as small as desired by choosing $s_0$ sufficiently large, 
and $B'>0$. % is fixed. 
Since $e^{-\kappa c}=\alpha\in (0,1)$, with $\lambda_{n+1}=\epsilon_{n+1}/\alpha^{n+1}$ we have 
%$K_{n+1}\leq K_n(1+A e^{nc})+B\epsilon_{n+1} \alpha^{-(n+1)}$ for $n\geq 0$. Next, let $\lambda_{n+1}=\epsilon_{n+1}/\alpha^{n+1}$. 
%Then 
$K_{n+1}\leq K_n(1+A e^{nc})+B\lambda_{n+1}$ for $n\geq 0$. Thus for $n\geq 2$ 
\[
K_n\leq \prod_{j=0}^{n-1} (1+Ae^{jc}) \cdot K_0 + B\sum_{j=1}^{n-1} \prod_{l=j}^{n-1} (1+Ae^{lc}) \cdot \lambda_j + \lambda_n.
\]
Let $p_n=\prod_{j=0}^{n-1} (1+Ae^{jc})$ for $n\geq 1$. Then $p_n\to p_\ast\in (1,\infty)$ as $n\to\infty$ and so  
\begin{equation} \label{eq.Knexplicitestimate}
	K_n\leq p_n K_0 + B p_n\sum_{j=1}^n \frac{\lambda_j}{p_j}, \quad n\geq 2. 
\end{equation}
Thus, if $\sum_{n=1}^\infty \lambda_n<+\infty$, then $K_n\leq K^\ast$ for all $n\geq 0$. 
On the other hand, if $\sum_{j=1}^n \lambda_j\to \infty$ as $n\to\infty$, then 
$K_n\leq K^\ast \sum_{j=1}^n \lambda_j$ for all $n\geq 1$, by applying Toeplitz lemma to the sum on the righthand side, proving both parts of the lemma. 

If $\lambda_n$ is summable, $K_n\leq K^\ast$. Since $\kappa<0$,  
%\[
%e^{-\kappa(s_0-nc)}\sup_{s\in [s_0-nc,s_0-(n+1)c]} |z(e^s)|
%\leq \sup_{s\in [s_0-nc,s_0-(n+1)c]]} e^{-\kappa s} |z(e^s)|\leq K^\ast,
%\]
%so 
$\sup_{t\in [e^{s_0-nc},e^{s_0-nc}e^{-c}]} |z(t)|\leq K^\ast [e^{(s_0-nc)}]^\kappa$. Next, for all $t$ sufficiently large there is a (unique) $n\in \mathbb{N}$ such that 
$e^{s_0-nc}\leq t<e^{s_0-nc}e^{-c}$. Since $\kappa<0$, $t^\kappa\geq e^{\kappa(s_0-nc)}e^{-\kappa c} = \alpha e^{\kappa(s_0-nc)}$. Thus
\[
|z(t)|\leq \sup_{t\in [e^{s_0-nc},e^{s_0-nc}e^{-c}]} |z(t)|\leq 
K^\ast [e^{(s_0-nc)}]^\kappa \leq K^\ast \frac{1}{\alpha} t^\kappa,
\]
as required. In the case $\lambda_n$ is not summable, $K_n\leq K^\ast \sum_{j=1}^n \lambda_j$. Since $\kappa<0$, we have 
\[
e^{-\kappa(s_0-nc)}\sup_{s\in [s_0-nc,s_0-(n+1)c]} |z(e^s)|
\leq \sup_{s\in [s_0-nc,s_0-(n+1)c]]} e^{-\kappa s} |z(e^s)|= K_n,
\]
so 
\[
\sup_{t\in [e^{s_0-nc},e^{s_0-nc}e^{-c}]} |z(t)|\leq K_n \alpha^n 
\left(\frac{e^{-\kappa c}}{\alpha}\right)^n e^{\kappa s_0} 
= K_n \alpha^n  e^{\kappa s_0}
\leq K^{\ast\ast} \alpha^n \sum_{j=1}^n \frac{\epsilon_j}{\alpha^j}.
\]
as required. 
\end{proof}
\begin{proof}[Proof of Theorem \ref{thm.growthbounds2}] 
Write $\varphi(t)=\text{O}(\psi(t)t^{\kappa+1})$. We prove the following subcases:
\begin{enumerate}
	\item[(A)] If $\psi\in \text{RV}_{\infty}(-1)$, $\psi\in L^1(0,\infty)$, then 
	$z(t)=\text{O}(t^\kappa)$ as $t\to\infty$.
	\item[(B)] If $\psi\in \text{RV}_{\infty}(-1)$, $\psi\not\in L^1(0,\infty)$, then 
	$	z(t)=\text{O}\left(t^\kappa\int_0^t \psi(s)\,ds\right)$, $t\to\infty$.
	\item[(C)] If  $\psi\in \text{RV}_{\infty}(\beta)$ for $\beta>-1$, then 
	the estimate in (B) holds. Equivalently, if $\varphi=\text{O}(\gamma)$ where $\gamma\in \text{RV}_\infty(\eta)$ and $\eta>\kappa$, then $z(t)=\text{O}(\gamma(t))$ as $t\to\infty$.
	\item[(D)] If  $\psi\in \text{RV}_{\infty}(\beta)$ for $\beta<-1$, then 
	the estimate in (A) holds. Equivalently, if $\varphi=\text{O}(\gamma)$ where $\gamma\in \text{RV}_\infty(\eta)$ and $\eta<\kappa$, then $z(t)=\text{O}(t^\kappa)$ as $t\to\infty$.  
\end{enumerate}	

Since $t\mapsto t^{\kappa+1}\psi(t)$ is regularly varying, and an estimate of the form $|\varphi(t)|\leq Kt^{\kappa+1}\psi(t)$ holds, by the uniform convergence theorem for regularly varying functions, $\epsilon_n = \sup_{s\in [e^{s_0-nc},e^{s_0-nc}e^{-c}]} |\varphi(s)|
\leq \bar{K} e^{(s_0-nc)(\kappa+1)} \psi(e^{s_0-nc})$. 
Since $\psi$ is regularly varying, we get  
$\epsilon_n\leq K^\ast  e^{-nc(\kappa+1)} \psi(e^{-nc})$, so $\lambda_n = \epsilon_n/\alpha^n \leq K^\ast \psi(e^{-nc}) e^{-nc}$. 
By the  uniform convergence theorem for regularly varying functions
\[
\frac{\int_x^{xe^{-c}} \psi(s)\,ds}{\psi(x)x} = \int_1^{e^{-c}} \frac{\psi(\eta x)}{\psi(x)}\,d\eta \to \int_1^{e^{-c}} \eta^\beta\,d\eta =: k>0, \quad x\to\infty,
\]
and so we have that 
\begin{equation} \label{eq.asyconvertsumint}
	e^{-nc} \psi(e^{-nc}) \sim \frac{1}{k} \int_{e^{-nc}}^{e^{-(n+1)c}} \psi(s)\,ds, \quad n\to\infty,
\end{equation}

For part (A), $\psi$ is integrable, so \eqref{eq.asyconvertsumint} shows that $(e^{-nc} \psi(e^{-nc}))$ is summable, and hence $(\lambda_n)$ is summable. Hence by the fundamental lemma, $z(t)=\text{O}(t^\kappa)$, $t\to\infty$. 

For part (D), since $\beta<-1$, $\psi$ is integrable, and by the same argument $z(t)=\text{O}(t^\kappa)$, $t\to\infty$. The second half of the claim follows by making the identification $\psi=\gamma/t^{\kappa+1}$, and since $\eta<\kappa$, $\psi\in \text{RV}_\infty(\beta)$ with $\beta=\eta-\kappa-1<-1$. 

For (B) and (C), $\psi$ is not integrable by hypothesis. Using \eqref{eq.asyconvertsumint}, we have  
\[
\sum_{j=1}^n \lambda_j \leq K^\ast \sum_{j=1}^n \psi(e^{-jc}) e^{-jc}
\leq K^{\ast \ast} \sum_{j=1}^n \frac{1}{k} \int_{e^{-jc}}^{e^{-(j+1)c}} \psi(s)\,ds
\leq \bar{K}  \int_{1}^{e^{-(n+1)c}} \psi(s)\,ds.
\]
Now, since $p_n>1$ and $p_n\to p^\ast$ as $n\to\infty$, by \eqref{eq.Knexplicitestimate}, we get 
\begin{multline*}
	\sup_{t\in [e^{s_0-nc},e^{s_0-nc}e^{-c}]}  |z(t)| \leq \alpha^n K_n
	\leq \alpha^n\left(p_n K_0 + p_n \sum_{j=1}^n \lambda_j/p_j\right)\\
	\leq C_1 \alpha^n + C_2 \alpha^n \sum_{j=1}^n \lambda_j
	\leq C_1 \alpha^n + C_2 \alpha^n \bar{K}  \int_{1}^{e^{-(n+1)c}} \psi(s)\,ds.
\end{multline*}
Since the integral diverges, we have the consolidated estimate
\[
\sup_{t\in [e^{s_0-nc},e^{s_0-nc}e^{-c}]}  |z(t)| \leq C_3 \alpha^n  \int_{1}^{e^{-(n+1)c}} \psi(s)\,ds.
\]
Finally, take $t\in [e^{s_0-nc},e^{s_0-(n+1)c}]$. Then $e^{-nc} e^{-c} \leq e^{-s_0}e^{-c}t$. Since $\psi$ is regularly varying, and $\psi$ is not integrable, $\int \psi$ is also regularly varying. Hence
\[
\int_1^{e^{-(n+1)c}} \psi(s)\,ds \leq \int_1^{e^{-c} e^{-s_0}t} \psi(s)\,ds
\leq C_4  \int_1^{t} \psi(s)\,ds.
\]
On the other hand, $e^{-nc}\leq  e^{-s_0} t$, so $\alpha^n=e^{-nc\kappa}\geq  e^{-\kappa s_0} t^\kappa$. Thus
\[
|z(t)|\leq 	\sup_{t\in [e^{s_0-nc},e^{s_0-nc}e^{-c}]}  |z(t)| \leq C_3 \cdot e^{-\kappa s_0} t^\kappa \cdot C_4  \int_1^{t} \psi(s)\,ds,
\]
as required. In part (B), and for the first part of the claim in (C), this suffices. 

For the second claim in (C), we have $\varphi=\text{O}(\gamma)$, with $\gamma\in \text{RV}_\infty(\eta)$ and $\eta>\kappa$. This allows us to take  $\psi(t)=\gamma(t)/t^{\kappa+1}\in \text{RV}_\infty(\beta)$ with $\beta=\eta-\kappa-1>-1$, so the hypotheses in the first part of the claim apply. Further, by the regular variation of $\psi$,
\[
\int_0^t \psi(s)\,ds \sim \frac{1}{\beta+1} t \psi(t) = \frac{1}{\eta-\kappa} \frac{\gamma(t)}{t^\kappa}, \quad t\to\infty,
\]  
and so applying the first part of the claim gives $z(t)=\text{O}(\gamma(t))$, $t\to\infty$. 

The conclusions of the different parts of the theorem can be combined. Parts (A) and (B) can be combined: let $\varphi=\text{O}(\gamma)$ where $\gamma \in \text{RV}_\infty(\kappa)$. Then $\psi(t):=\gamma(t)/t^{\kappa+1}$ for $t\geq 1$ is in $\text{RV}_\infty(-1)$, and $\varphi(t)=\text{O}(t^{\kappa+1}\psi(t))$. Then, in (A) and (B), we have 
\[
z(t)=\text{O}\left(t^\kappa\int_1^t \psi(s)\,ds\right) = \text{O}\left(t^\kappa\int_1^t \frac{\gamma(s)}{s^{\kappa+1}}\,ds\right), \quad t\to\infty.
\] 
For part (C), $\eta>\kappa$, so $\psi$  %$\gamma_1(t):=\gamma(t)/t^{\kappa+1}$ 
is in $\text{RV}_\infty(\eta-\kappa-1)$. Since $\eta-\kappa-1>0$, the function in the big O--estimate is asymptotic to $\gamma(t)/(\eta-\kappa)$ as $t\to\infty$,  
%\[
%t^\kappa\int_1^t \frac{\gamma(s)}{s^{\kappa+1}}\,ds \sim t^\kappa\cdot \frac{1}{\eta-\kappa}t\frac{\gamma(t)}{t^{\kappa+1}} \sim \frac{1}{\eta-\kappa}\gamma(t),
%\]
so by part (C), 
$z(t)=\text{O}(\gamma(t)) = \text{O}\left(t^\kappa\int_1^t \frac{\gamma(s)}{s^{\kappa+1}}\,ds\right)$ as $t\to\infty$,
which agrees with (A) and (B). Finally, for (D), where $\eta<\kappa$, $\psi \in \text{RV}_\infty(\eta-\kappa-1)$. Since $\eta-\kappa-1<-1$, $\psi \in L^1(1,\infty)$, and so by (D)
\[
z(t)=\text{O}(t^\kappa) = \text{O}\left(t^\kappa\int_1^t \frac{\gamma(s)}{s^{\kappa+1}}\,ds\right), \quad t\to\infty.
\] 
This allows the desired consolidation of the statement of the theorem.
\end{proof}

\section{Conclusions}

The paper considers perturbations of the asymptotically stable deterministic pantograph equation 
\[
u'(t)=bu(t)+au(qt)
\]
where $a\neq 0$. As is well--known, all solutions of this equation tend to zero if and only if $b<0$ and $|a|<|b|$. We thus presume these conditions hold, and ask what is the asymptotic behaviour of forced equations 
\begin{align*}
x'(t)&=bx(t)+ax(qt)+f(t), \\
dX(t)&=bX(t)+aX(qt)+f(t)\,dt +\sigma(t)\,dB(t).
\end{align*}  
We can characterise, in terms of $f$ and $\sigma$, necessary and sufficient conditions such that solutions are bounded or convergent to zero, assuming only that the underlying equation has stable solutions. These conditions depend on the time averages of $f$ and $\sigma$ over moving compact intervals, but do not depend on the parameters ($a$, $b$ and $q$) of the underlying unperturbed equation.
The conditions allow $f$ and $\sigma$ to behave very badly on a pointwise basis, while being satisfied by often--assumed stability conditions such as $f(t)\to 0$ or $f\in L^1(0,\infty)$. The conditions on the stochastic perturbation are more delicate; the a.s. stability or a.s. boundedness is captured by the convergence or divergence of an infinite sum which depends on a parameter. This condition is known to characterise the a.s. stability or a.s. boundedness of solutions of the stochastic ordinary differential equation
\[
dY_0(t)=-Y_0(t)\,dt + \sigma(t)\,dB(t)
\]
a result proven in an earlier work. The necessary and sufficient conditions are a little opaque, but a very good sufficient condition is that 
\begin{equation} \label{eq.sig2nlognto0}
\sigma_1^2(n) \log n\to 0, \quad n\to\infty;
\end{equation}
where $\sigma_1^2(n)=\int_{n-1}^{n} \sigma^2(s)\,ds$; 
indeed, if $\sigma^2_n$ is decreasing, \eqref{eq.sig2nlognto0} is necessary, so once again we see that it is the average behaviour of $\sigma^2$ over moving compact intervals that determines the asymptotic behaviour. 

A key observation in the entire analysis is that the stability of the deterministic functional differential equation is inherited from the ordinary differential equation 
\[
y'=-y+f
\] 
and the SFDE from $Y=Y_0+y$, which solves a ``doubly perturbed'' SDE. It turns out $Y$ tends to zero if and only if both $Y_0$ and $y$ are, while $Y$ is bounded if and only if $Y_0$ and $y$ are.   

As well as characterising boundedness and convergence, the approach enables us to determine very sharp conditions describing the rate of growth or decay of deterministic equations. If the solution is to grow, or has a growing running maximum, the rate of growth of its running maximum is that of $y$, and this can be characterised entirely through the time averages $\int_{t-\theta}^t f(s)\,ds$, which has the same asymptotic order of $y$. In terms of estimating the growth order, and characterising when certain growth orders of solutions occur, these results are unimprovable. Crucially, they show that relying on pointwise estimates on $f$ may give very poor and misleading asymptotic information. 

When solutions tend to zero, they do so slowly in most cases, since the unperturbed equation has solutions which tend to zero at the rate $t^\kappa$, where $\kappa<0$ is known in terms of $a$, $b$ and $q$. Existing analysis for general equations with unbounded delay, applied to the pantograph equation, gives estimates of $x(t)=O(f(t))$ when $f(t)=O(t^\beta)$ and $f'(t)=O(t^{\beta-1})$, and $\beta>\kappa$, and $x(t)=O(t^\kappa)$ when $f(t)=O(t^\beta)$ for $\beta<\kappa$. We show, in contrast, that when $y(t)=O(\gamma(t))$ and $\gamma$ is regularly varying, then 
\[
x(t)=O\left(t^\kappa \int_1^t \frac{\gamma(s)}{s^{\kappa+1}}\,ds\right),\quad t\geq 0
\]     
which shows that neither pointwise nor derivative bounds are needed on $f$. Moreover, the estimate we establish is sharper in some cases than existing estimates, is never weaker than existing estimates, and is unimprovable in at least some situations (for equations in which $a>0$). 

The rate of growth of $y$ plays a determining role the asymtotic behaviour when $y$ dominates the solution of the unperturbed equation. For instance, if $\gamma$ is regularly varying with index greater than $\kappa$, or $\gamma$ is increasing, then $x(t)=o(\gamma(t))$ if and only if $y(t)=o(\gamma(t))$, $x(t)=O(\gamma(t))$ if and only if $y(t)=O(\gamma(t))$ and $\limsup_{t\to\infty} |x(t)|/\gamma(t)\in (0,\infty)$ if and only if $\limsup_{t\to\infty}|y(t)|/\gamma(t)\in (0,\infty)$. 

We have sufficient conditions which describe when the rate of decay of the underlying pantograph equation is preserved. If $y(t)=O(\gamma(t))$, $\gamma$ is regularly varying, and $t\mapsto\gamma(t)/t^{\kappa+1}$ is integrable, then $x(t)=O(t^\kappa)$ as $t\to\infty$. These conditions seem, in certain cases at least, unimprovable, since we have shown by means of general examples, that if the integrability condition is relaxed, the unperturbed rate of decay will not be preserved. Good necessary conditions prove for the moment, more elusive. We know that if $x(t)=O(t^\kappa)$, then $y(t)=O(t^\kappa)$; but taking the sufficient condition above into account, it seems that the maximal allowable rate of growth of $y$ would be such that $t\mapsto |y(t)|/t^{\kappa+1}\in L^1(0,\infty)$, a condition satisfied in the case $y(t)=O(t^\kappa\{\log t\log\log t(\log\log \log t)^2\}^{-1})$, for instance. Thus, our best converse result in the case of a small perturbation is still suboptimal. 

Although it might seem strange that we can characterise asymptotic bounds for ``large'' perturbations, but not ``small'' ones, we saw by means of an example that it is possible to arrange for the solution of the perturbed equation to decay at an arbitrarily fast rate, for a very carefully chosen and very rapidly decaying forcing term. Therefore, there is potentially great sensitivity in the rate of decay of solutions when there are small forcing terms, and this probably explains in part why it is hard to prove a very refined converse. As a first step though, our converse is not unacceptable: the above discussion suggests that our estimate on the critical $y$ is only off by a factor which grows at around the (very slow) rate of $\log t$ as $t\to\infty$.   

In this paper, we have not attempted to give quantitative asymptotic estimates for the solutions of stochastic equations. In part, this is because a characterisation of the rates of decay and growth of solutions of \eqref{eq.Y} is needed, and this task forms part of work in progress. However, once the task of characterising the asymptotic rates of growth and decay of $Y$ is completed, the analysis in this paper (see Theorems~\ref{thm.smallstochpert} and \ref{thm.detlargepert}) should be reusable with minor modifications, with Theorems~\ref{thm.growthbounds2} and \ref{thm.monotonepert} applicable to $Z=X-Y$, with the estimate on $\varphi$ being read off from $Y$.

There are two other directions where the results of this paper could readily be developed. One is in the direction of ``unstable'' equations, which is to say equations where the condition \eqref{eq.unforcedasstable} do not hold, and the unperturbed equation has unbounded solutions. The case when $b>0$ should be more straightforward, since it that case, solutions of the unperturbed equation grow like $e^{bt}$, so one would hope that the solution of the perturbed equation would grow at the faster rate between $e^{bt}$ and $\gamma$, where $\gamma$ describes the rate of growth of $y$.

The case when $b<0$, and $|a|>|b|$ is more interesting, since in that case, the unperturbed equation exhibits power law growth at rate $t^\kappa$; now $\kappa>0$. In this situation, we expect analogues of Theorems~\ref{thm.detlargepert} and \ref{thm.smallstochpert} to hold. Revisiting the proof of the fundamental Lemma~\ref{lemma.fundamental} reveals that the estimates hold up to \eqref{eq.Knexplicitestimate}, with perhaps some cosmetic changes needed earlier. Therefore it is probable that we will be able to estimate as in the proof of Theorem~\ref{thm.growthbounds2}. As a result, we expect a theorem analogous to Theorems~\ref{thm.detlargepert} to hold if $y(t)=O(\gamma(t))$ and  $\gamma$ is regularly varying with index greater than $\kappa$ (or $\gamma$ is increasing faster than a regularly varying function), while a theorem
analogous to Theorem~\ref{thm.smallstochpert} should hold if $\gamma$ has index less than or equal to $\kappa$. 

In Section 5, we commented that generalisations to equations with unbounded delay are possible, and we were able to sketch convergence and boundedness characterisation results directly. However, if growth bounds on solutions are required, we noted but that further work would be required, both on stochastic ordinary differential equations, as well as the analysis of Cermak in the deterministic case in order to sharpen estimates. Extensions of the convergence and boundedness results to the finite dimensional case are feasible too, but at the cost of stronger sufficient conditions being imposed on the structure of the unperturbed equation. We conjecture that a perturbation theorem which states that property \eqref{eq.perturbedpropertymultid} holds under the conditions \eqref{eq.iser1} and \eqref{eq.iser2} is true, but do not have immediate ideas how to proceed with it. The situation wherein there is multiplicative noise has been studied more extensively in the literature, but very refined results (i.e., results which classify the dynamics) appear to be lacking. However, by using the spirit of the approach taken here (decomposition of the solution into the solution of a stochastic ordinary differential equation, and the $C^1$ solution of a random functional differential equation), we have produced a result in the spirit of the very first results of Kato and McLeod: exponential asymptotic behaviour when the underlying non--delay equation is unstable, and subexponential asymptotic behaviour when the underlying non--delay equation is asymptotically stable. It is appealing that such progress is made for an equation with multiplicative noise by a multiplicative decomposition, since the bulk of this paper deals with additive noise, and is (we believe) successfully tackled by an additive decomposition of the solution. Further exploration of the multiplicative decomposition proposed here warrants a deeper investigation.     

We finish with a remark that when perturbations are ``small'', the solution $x$ of the perturbed pantograph equation is $O(t^\kappa)$. But it is well--known that solutions of the unperturbed equation exhibit log--periodic oscillation in the case that $a<0$. With further careful analysis, it may be possible to show, for sufficiently small forcing terms, that the perturbed equation inherits this oscillatory asymptotic behaviour, even in the stochastic case. Such results would require developing analogues of Kato and McLeod's arguments. Since some of their analysis involves repeated differentiation of the equation, the double mollification ideas (or even higher order mollifications) sketched in Section 5.2 will be necessary if the results are not to require estimates of $f$ and its derivatives in the deterministic case; indeed, such mollification would be necessary if one wished to prove more refined asymptotic results for stochastic equations, as suggested in Section 5.2.   

%and borrowing a factor of $e^{\kappa c}$ in the first integral on the right, this gives 
%\begin{multline*}
%	|w(s)|e^{\kappa s}  
%	\leq \{|w(-(n+1)c)|e^{-\kappa(n+1)c}\} e^{-|b|(e^s-e^{-(n+1)c})}
%	\\+|b|e^{-|b|e^s} e^{-\kappa c}\int_{-(n+1)c}^s e^t e^{|b|e^t}
%	\{e^{\kappa (t+c)}|w(t+c)|\}\,dt
%	\\+e^{-|b|e^s}\int_{-(n+1)c}^s e^{|b|e^t}e^t|\varphi(e^t)|\,dt.	
%\end{multline*}
%The factors in curly braces in the first term on the righthand side, as well as in the first integral, are both $\leq K_n$. Hence
%\begin{multline*}
%	|w(s)|e^{\kappa s}  
%	\leq K_n e^{-|b|(e^s-e^{-(n+1)c})}
%	+e^{-\kappa c} |b|e^{-|b|e^s} \int_{-(n+1)c}^s e^t e^{|b|e^t}\,dt \cdot K_n
%	\\+e^{-|b|e^s}\int_{-(n+1)c}^s e^{|b|e^t}e^t|\varphi(e^t)|\,dt.	
%\end{multline*}
%Thus, for $s\in I_{n+1}$, since $e^{-\kappa c}=|a/b|=:\alpha$, and for $t\in [-(n+1)c,s]\subseteq I_{n+1}$ we have $|\varphi(e^t)|\leq \epsilon_{n+1}$, we get 
%\begin{multline*}
%	|w(s)|e^{\kappa s}  
%	\leq K_n e^{-|b|(e^s-e^{-(n+1)c})}
%	+\alpha |b|e^{-|b|e^s} \int_{-(n+1)c}^s e^t e^{|b|e^t}\,dt \cdot K_n
%	\\+\frac{1}{|b|}\cdot |b|e^{-|b|e^s}\int_{-(n+1)c}^s e^{|b|e^t}e^t\,dt \cdot \epsilon_{n+1}.
%\end{multline*}
%Next, write 
%\[
%J_n(s):=|b|e^{-|b|e^s} \int_{-(n+1)c}^s e^t e^{|b|e^t}\,dt, \quad s\in I_{n+1},
%\]
%so that $J_n(s)=1-e^{-|b|(e^s-e^{-(n+1)c})}$, and for $s\in I_{n+}$, we have 
%\[
%|w(s)|e^{\kappa s}  
%\leq K_n e^{-|b|(e^s-e^{-(n+1)c})}
%+\alpha J_n(s) K_n+\frac{1}{|b|} J_n(s) \cdot \epsilon_{n+1}.
%\]
%Using the formula for $J_n$ in the first term, and the estimate $J_n(s)\leq 1$ in the second, we get 
%\[
%|w(s)|e^{\kappa s}  
%\leq K_n e^{-|b|(e^s-e^{-(n+1)c})}
%+\alpha (1-e^{-|b|(e^s-e^{-(n+1)c})}) K_n+\frac{1}{|b|} \epsilon_{n+1}
%=\alpha K_n + \frac{1}{|b|} \epsilon_{n+1} + (1-\alpha)
%\]

%	% Always give a unique label
%	% and use \ref{<label>} for cross-references
%	% and \cite{<label>} for bibliographic references
%	% use \sectionmark{}
%	% to alter or adjust the section heading in the running head
%	Instead of simply listing headings of different levels we recommend to let every heading be followed by at least a short passage of text.  Further on please use the \LaTeX\ automatism for all your cross-references and citations.
%	
%	Please note that the first line of text that follows a heading is not indented, whereas the first lines of all subsequent paragraphs are.
%	
%	Use the standard \verb|equation| environment to typeset your equations, e.g.
%	%
%	\begin{equation}
%		a \times b = c\;,
%	\end{equation}
%	%
%	however, for multiline equations we recommend to use the \verb|eqnarray| environment\footnote{In physics texts please activate the class option \texttt{vecphys} to depict your vectors in \textbf{\itshape boldface-italic} type - as is customary for a wide range of physical subjects}.
%	\begin{eqnarray}
%		\left|\nabla U_{\alpha}^{\mu}(y)\right| &\le&\frac1{d-\alpha}\int
%		\left|\nabla\frac1{|\xi-y|^{d-\alpha}}\right|\,d\mu(\xi) =
%		\int \frac1{|\xi-y|^{d-\alpha+1}} \,d\mu(\xi)  \\
%		&=&(d-\alpha+1) \int\limits_{d(y)}^\infty
%		\frac{\mu(B(y,r))}{r^{d-\alpha+2}}\,dr \le (d-\alpha+1)
%		\int\limits_{d(y)}^\infty \frac{r^{d-\alpha}}{r^{d-\alpha+2}}\,dr
%		\label{eq:01}
%	\end{eqnarray}
%	
%	\subsection{Subsection Heading}
%	\label{subsec:2}
%	Instead of simply listing headings of different levels we recommend to let every heading be followed by at least a short passage of text.  Further on please use the \LaTeX\ automatism for all your cross-references\index{cross-references} and citations\index{citations} as has already been described in Sect.~\ref{sec:2}.
%	
%	\begin{quotation}
%		Please do not use quotation marks when quoting texts! Simply use the \verb|quotation| environment -- it will automatically be rendered in line with the preferred layout.
%	\end{quotation}
%	
%	
%	\subsubsection{Subsubsection Heading}
%	Instead of simply listing headings of different levels we recommend to let every heading be followed by at least a short passage of text.  Further on please use the \LaTeX\ automatism for all your cross-references and citations as has already been described in Sect.~\ref{subsec:2}, see also Fig.~\ref{fig:1}\footnote{If you copy text passages, figures, or tables from other works, you must obtain \textit{permission} from the copyright holder (usually the original publisher). Please enclose the signed permission with the manuscript. The sources\index{permission to print} must be acknowledged either in the captions, as footnotes or in a separate section of the book.}
%	
%	Please note that the first line of text that follows a heading is not indented, whereas the first lines of all subsequent paragraphs are.
%	
%	% For figures use
%	%
%	\begin{figure}[b]
%		\sidecaption
%		% Use the relevant command for your figure-insertion program
%		% to insert the figure file.
%		% For example, with the graphicx style use
%		\includegraphics[scale=.65]{figure}
%		%
%		% If no graphics program available, insert a blank space i.e. use
%		%\picplace{5cm}{2cm} % Give the correct figure height and width in cm
%		%
%		\caption{If the width of the figure is less than 7.8 cm use the \texttt{sidecapion} command to flush the caption on the left side of the page. If the figure is positioned at the top of the page, align the sidecaption with the top of the figure -- to achieve this you simply need to use the optional argument \texttt{[t]} with the \texttt{sidecaption} command}
%		\label{fig:1}       % Give a unique label
%	\end{figure}
%	
%	
%	\paragraph{Paragraph Heading} %
%	Instead of simply listing headings of different levels we recommend to let every heading be followed by at least a short passage of text.  Further on please use the \LaTeX\ automatism for all your cross-references and citations as has already been described in Sect.~\ref{sec:2}.
%	
%	Please note that the first line of text that follows a heading is not indented, whereas the first lines of all subsequent paragraphs are.
%	
%	For typesetting numbered lists we recommend to use the \verb|enumerate| environment -- it will automatically rendered in line with the preferred layout.
%	
%	\begin{enumerate}
%		\item{Livelihood and survival mobility are oftentimes coutcomes of uneven socioeconomic development.}
%		\begin{enumerate}
	%			\item{Livelihood and survival mobility are oftentimes coutcomes of uneven socioeconomic development.}
	%			\item{Livelihood and survival mobility are oftentimes coutcomes of uneven socioeconomic development.}
	%		\end{enumerate}
%		\item{Livelihood and survival mobility are oftentimes coutcomes of uneven socioeconomic development.}
%	\end{enumerate}
%	
%	
%	\subparagraph{Subparagraph Heading} In order to avoid simply listing headings of different levels we recommend to let every heading be followed by at least a short passage of text. Use the \LaTeX\ automatism for all your cross-references and citations as has already been described in Sect.~\ref{sec:2}, see also Fig.~\ref{fig:2}.
%	
%	For unnumbered list we recommend to use the \verb|itemize| environment -- it will automatically be rendered in line with the preferred layout.
%	
%	\begin{itemize}
%		\item{Livelihood and survival mobility are oftentimes coutcomes of uneven socioeconomic development, cf. Table~\ref{tab:1}.}
%		\begin{itemize}
	%			\item{Livelihood and survival mobility are oftentimes coutcomes of uneven socioeconomic development.}
	%			\item{Livelihood and survival mobility are oftentimes coutcomes of uneven socioeconomic development.}
	%		\end{itemize}
%		\item{Livelihood and survival mobility are oftentimes coutcomes of uneven socioeconomic development.}
%	\end{itemize}
%	
%	\begin{figure}[t]
%		\sidecaption[t]
%		% Use the relevant command for your figure-insertion program
%		% to insert the figure file.
%		% For example, with the option graphics use
%		\includegraphics[scale=.65]{figure}
%		%
%		% If no graphics program available, insert a blank space i.e. use
%		%\picplace{5cm}{2cm} % Give the correct figure height and width in cm
%		%
%		%\caption{Please write your figure caption here}
%		\caption{If the width of the figure is less than 7.8 cm use the \texttt{sidecapion} command to flush the caption on the left side of the page. If the figure is positioned at the top of the page, align the sidecaption with the top of the figure -- to achieve this you simply need to use the optional argument \texttt{[t]} with the \texttt{sidecaption} command}
%		\label{fig:2}       % Give a unique label
%	\end{figure}
%	
%	\runinhead{Run-in Heading Boldface Version} Use the \LaTeX\ automatism for all your cross-references and citations as has already been described in Sect.~\ref{sec:2}.
%	
%	\subruninhead{Run-in Heading Boldface and Italic Version} Use the \LaTeX\ automatism for all your cross-refer\-ences and citations as has already been described in Sect.~\ref{sec:2}\index{paragraph}.
%	
%	\subsubruninhead{Run-in Heading Displayed Version} Use the \LaTeX\ automatism for all your cross-refer\-ences and citations as has already been described in Sect.~\ref{sec:2}\index{paragraph}.
%	% Use the \index{} command to code your index words
%	%
%	% For tables use
%	%
%	\begin{table}[!t]
%		\caption{Please write your table caption here}
%		\label{tab:1}       % Give a unique label
%		%
%		% Follow this input for your own table layout
%		%
%		\begin{tabular}{p{2cm}p{2.4cm}p{2cm}p{4.9cm}}
	%			\hline\noalign{\smallskip}
	%			Classes & Subclass & Length & Action Mechanism  \\
	%			\noalign{\smallskip}\svhline\noalign{\smallskip}
	%			Translation & mRNA$^a$  & 22 (19--25) & Translation repression, mRNA cleavage\\
	%			Translation & mRNA cleavage & 21 & mRNA cleavage\\
	%			Translation & mRNA  & 21--22 & mRNA cleavage\\
	%			Translation & mRNA  & 24--26 & Histone and DNA Modification\\
	%			\noalign{\smallskip}\hline\noalign{\smallskip}
	%		\end{tabular}
%		$^a$ Table foot note (with superscript)
%	\end{table}
%	%
%	\section{Section Heading}
%	\label{sec:3}
%	% Always give a unique label
%	% and use \ref{<label>} for cross-references
%	% and \cite{<label>} for bibliographic references
%	% use \sectionmark{}
%	% to alter or adjust the section heading in the running head
%	Instead of simply listing headings of different levels we recommend to let every heading be followed by at least a short passage of text.  Further on please use the \LaTeX\ automatism for all your cross-references and citations as has already been described in Sect.~\ref{sec:2}.
%	
%	Please note that the first line of text that follows a heading is not indented, whereas the first lines of all subsequent paragraphs are.
%	
%	If you want to list definitions or the like we recommend to use the enhanced \verb|description| environment -- it will automatically rendered in line with the preferred layout.
%	
%	\begin{description}[Type 1]
%		\item[Type 1]{That addresses central themes pertainng to migration, health, and disease. In Sect.~\ref{sec:1}, Wilson discusses the role of human migration in infectious disease distributions and patterns.}
%		\item[Type 2]{That addresses central themes pertainng to migration, health, and disease. In Sect.~\ref{subsec:2}, Wilson discusses the role of human migration in infectious disease distributions and patterns.}
%	\end{description}
%	
%	\subsection{Subsection Heading} %
%	In order to avoid simply listing headings of different levels we recommend to let every heading be followed by at least a short passage of text. Use the \LaTeX\ automatism for all your cross-references and citations citations as has already been described in Sect.~\ref{sec:2}.
%	
%	Please note that the first line of text that follows a heading is not indented, whereas the first lines of all subsequent paragraphs are.
%	
%	\begin{svgraybox}
%		If you want to emphasize complete paragraphs of texts we recommend to use the newly defined class option \verb|graybox| and the newly defined environment \verb|svgraybox|. This will produce a 15 percent screened box 'behind' your text.
%		
%		If you want to emphasize complete paragraphs of texts we recommend to use the newly defined class option and environment \verb|svgraybox|. This will produce a 15 percent screened box 'behind' your text.
%	\end{svgraybox}
%	
%	
%	\subsubsection{Subsubsection Heading}
%	Instead of simply listing headings of different levels we recommend to let every heading be followed by at least a short passage of text.  Further on please use the \LaTeX\ automatism for all your cross-references and citations as has already been described in Sect.~\ref{sec:2}.
%	
%	Please note that the first line of text that follows a heading is not indented, whereas the first lines of all subsequent paragraphs are.
%	
%	\begin{theorem}
%		Theorem text goes here.
%	\end{theorem}
%	%
%	% or
%	%
%	\begin{definition}
%		Definition text goes here.
%	\end{definition}
%	
%	\begin{proof}
%		%\smartqed
%		Proof text goes here.
%		%\qed
%	\end{proof}
%	
%	\paragraph{Paragraph Heading} %
%	Instead of simply listing headings of different levels we recommend to let every heading be followed by at least a short passage of text.  Further on please use the \LaTeX\ automatism for all your cross-references and citations as has already been described in Sect.~\ref{sec:2}.
%	
%	Note that the first line of text that follows a heading is not indented, whereas the first lines of all subsequent paragraphs are.
%	%
%	% For built-in environments use
%	%
%	\begin{theorem}
%		Theorem text goes here.
%	\end{theorem}
%	%
%	\begin{definition}
%		Definition text goes here.
%	\end{definition}
%	%
%	\begin{proof}
%		%\smartqed
%		Proof text goes here.
%		%\qed
%	\end{proof}
%	%
%	\begin{trailer}{Trailer Head}
%		If you want to emphasize complete paragraphs of texts in an \verb|Trailer Head| we recommend to
%		use  \begin{verbatim}\begin{trailer}{Trailer Head}
		%				...
		%		\end{trailer}\end{verbatim}
%	\end{trailer}
%	%
%	\begin{questype}{Questions}
%		If you want to emphasize complete paragraphs of texts in an \verb|Questions| we recommend to
%		use  \begin{verbatim}\begin{questype}{Questions}
		%				...
		%		\end{questype}\end{verbatim}
%	\end{questype}
%	\eject%
%	\begin{important}{Important}
%		If you want to emphasize complete paragraphs of texts in an \verb|Important| we recommend to
%		use  \begin{verbatim}\begin{important}{Important}
		%				...
		%		\end{important}\end{verbatim}
%	\end{important}
%	%
%	\begin{warning}{Attention}
%		If you want to emphasize complete paragraphs of texts in an \verb|Attention| we recommend to
%		use  \begin{verbatim}\begin{warning}{Attention}
		%				...
		%		\end{warning}\end{verbatim}
%	\end{warning}
%	
%	\begin{programcode}{Program Code}
%		If you want to emphasize complete paragraphs of texts in an \verb|Program Code| we recommend to
%		use
%		
%		\verb|\begin{programcode}{Program Code}|
	%			
	%			\verb|\begin{verbatim}...\end{verbatim}|
	%			
	%			\verb|\end{programcode}|
%		
%	\end{programcode}
%	%
%	\begin{tips}{Tips}
%		If you want to emphasize complete paragraphs of texts in an \verb|Tips| we recommend to
%		use  \begin{verbatim}\begin{tips}{Tips}
		%				...
		%		\end{tips}\end{verbatim}
%	\end{tips}
%	\eject
%	%
%	\begin{overview}{Overview}
%		If you want to emphasize complete paragraphs of texts in an \verb|Overview| we recommend to
%		use  \begin{verbatim}\begin{overview}{Overview}
		%				...
		%		\end{overview}\end{verbatim}
%	\end{overview}
%	\begin{backgroundinformation}{Background Information}
%		If you want to emphasize complete paragraphs of texts in an \verb|Background|
%		\verb|Information| we recommend to
%		use
%		
%		\verb|\begin{backgroundinformation}{Background Information}|
	%			
	%			\verb|...|
	%			
	%			\verb|\end{backgroundinformation}|
%	\end{backgroundinformation}
%	\begin{legaltext}{Legal Text}
%		If you want to emphasize complete paragraphs of texts in an \verb|Legal Text| we recommend to
%		use  \begin{verbatim}\begin{legaltext}{Legal Text}
		%				...
		%		\end{legaltext}\end{verbatim}
%	\end{legaltext}
%	%
%\begin{acknowledgement}
%EL was supported by Science Foundation Ireland (16/IA/4443) while a postgraduate student at Dublin City University. JA is supported by the RSE Saltire Facilitation Network on Stochastic Differential Equations: Theory, Numerics and Applications (RSE1832). 
%
%The authors wish to thank the anonymous referees for their very careful reading of the manuscript, and for prompting many interesting questions concerning generalisations of the original submission. This has resulted in a much more comprehensive paper, prompting among other things the treatment of general unbounded delay, multidimensional equations, and multiplicative noise. 
%
%We also wish to thank Prof. Elena Braverman for the opportunity to present this work at the session on ``Delay and stochastic differential equations with in Modelling in Finance, Life Sciences, and engineering'' at ICIAM2023, and for the chance to submit it to this special volume.    
%
%Finally, the first author also wishes to acknowledge the love and support of his late wife, Huizhong Appleby--Wu, with whom he discussed early results in this paper; She also was in Tokyo in August 2023, and gave advice on the presentation of our talk at the ICIAM meeting. HA--W died suddenly and unexpectedly, at the age of 41, in January 2024, and before this paper was submitted for publication. May she rest in peace.  
%\end{acknowledgement}
%%
%\ethics{Competing Interests}{The authors have no conflicts of interest to declare that are relevant to the content of this chapter.}
%
%\eject
%	
%	\ethics{Ethics Approval}{If your chapter includes primary studies with humans please declare the adherence of ethical standards. Example text: This study was performed in line with the principles of the Declaration of Helsinki. Approval was granted by the Ethics Committee of University B (Date.../No. ...).\newline 
%		In addition, for human participants, authors are required to include a statement that informed consent (to participate and/or to publish) was obtained from individual participants or parents/guardians if the participant is minor or incapable.\newline
%		If animals are studied, authors should make sure that the legal requirements or guidelines in the country and/or state or province for the care and use of animals have been followed or specify that no ethics approval was required.}
%
%\section*{Appendix}
%\addcontentsline{toc}{section}{Appendix}
%
%
%	When placed at the end of a chapter or contribution (as opposed to at the end of the book), the numbering of tables, figures, and equations in the appendix section continues on from that in the main text. Hence please \textit{do not} use the \verb|appendix| command when writing an appendix at the end of your chapter or contribution. If there is only one the appendix is designated ``Appendix'', or ``Appendix 1'', or ``Appendix 2'', etc. if there is more than one.
%
%	
%	\section{Styling of References}
%	References may be \textit{cited} in the text either by number (preferred) or by author/year.\footnote{Make sure that all references from the list are cited in the text. Those not cited should be moved to a separate \textit{Further Reading} section or chapter.} If the citatiion in the text is numbered, the reference list should be arranged in ascending order. If the citation in the text is author/year, the reference list should be \textit{sorted} alphabetically and if there are several works by the same author, the following order should be used:
%	\begin{enumerate}
%		\item all works by the author alone, ordered chronologically by year of publication
%		\item all works by the author with a coauthor, ordered alphabetically by coauthor
%		\item all works by the author with several coauthors, ordered chronologically by year of publication.
%	\end{enumerate}
%	The \textit{styling} of references\footnote{Always use the standard abbreviation of a journal's name according to the ISSN \textit{List of Title Word Abbreviations}, see \url{http://www.issn.org/en/node/344}} depends on the subject of your book:
%	\begin{itemize}
%		\item The \textit{two} recommended styles for references in books on \textit{mathematical, physical, statistical and computer sciences} are depicted in ~\cite{science-contrib, science-online, science-mono, science-journal, science-DOI} and ~\cite{phys-online, phys-mono, phys-journal, phys-DOI, phys-contrib}.
%		\item Examples of the most commonly used reference style in books on \textit{Psychology, Social Sciences} are~\cite{psysoc-mono, psysoc-online,psysoc-journal, psysoc-contrib, psysoc-DOI}.
%		\item Examples for references in books on \textit{Humanities, Linguistics, Philosophy} are~\cite{humlinphil-journal, humlinphil-contrib, humlinphil-mono, humlinphil-online, humlinphil-DOI}.
%		\item Examples of the basic Springer Nature style used in publications on a wide range of subjects such as \textit{Computer Science, Economics, Engineering, Geosciences, Life Sciences, Medicine, Biomedicine} are ~\cite{basic-contrib, basic-online, basic-journal, basic-DOI, basic-mono}. 
%	\end{itemize}
%	%}

%	\begin{thebibliography}{99.}%
%	% and use \bibitem to create references.
%	%
%	% Use the following syntax and markup for your references if 
%	% the subject of your book is from the field 
%	% "Mathematics, Physics, Statistics, Computer Science"
%	%
%	% Contribution 
%	\bibitem{science-contrib} Broy, M.: Software engineering --- from auxiliary to key technologies. In: Broy, M., Dener, E. (eds.) Software Pioneers, pp. 10-13. Springer, Heidelberg (2002)
%	%
%	% Online Document
%	\bibitem{science-online} Dod, J.: Effective substances. In: The Dictionary of Substances and Their Effects. Royal Society of Chemistry (1999) Available via DIALOG. \\
%	\url{http://www.rsc.org/dose/title of subordinate document. Cited 15 Jan 1999}
%	%
%	% Monograph
%	\bibitem{science-mono} Geddes, K.O., Czapor, S.R., Labahn, G.: Algorithms for Computer Algebra. Kluwer, Boston (1992) 
%	%
%	% Journal article
%	\bibitem{science-journal} Hamburger, C.: Quasimonotonicity, regularity and duality for nonlinear systems of partial differential equations. Ann. Mat. Pura. Appl. \textbf{169}, 321--354 (1995)
%	%
%	% Journal article by DOI
%	\bibitem{science-DOI} Slifka, M.K., Whitton, J.L.: Clinical implications of dysregulated cytokine production. J. Mol. Med. (2000) doi: 10.1007/s001090000086 
%	%
%	%\bigskip
%	
%	% Use the following (APS) syntax and markup for your references if 
%	% the subject of your book is from the field 
%	% "Mathematics, Physics, Statistics, Computer Science"
%	%
%	% Online Document
%	\bibitem{phys-online} J. Dod, in \textit{The Dictionary of Substances and Their Effects}, Royal Society of Chemistry. (Available via DIALOG, 1999), 
%	\url{http://www.rsc.org/dose/title of subordinate document. Cited 15 Jan 1999}
%	%
%	% Monograph
%	\bibitem{phys-mono} H. Ibach, H. L\"uth, \textit{Solid-State Physics}, 2nd edn. (Springer, New York, 1996), pp. 45-56 
%	%
%	% Journal article
%	\bibitem{phys-journal} S. Preuss, A. Demchuk Jr., M. Stuke, Appl. Phys. A \textbf{61}
%	%
%	% Journal article by DOI
%	\bibitem{phys-DOI} M.K. Slifka, J.L. Whitton, J. Mol. Med., doi: 10.1007/s001090000086
%	%
%	% Contribution 
%	\bibitem{phys-contrib} S.E. Smith, in \textit{Neuromuscular Junction}, ed. by E. Zaimis. Handbook of Experimental Pharmacology, vol 42 (Springer, Heidelberg, 1976), p. 593
%	%
%	%\bigskip
%	%
%	% Use the following syntax and markup for your references if 
%	% the subject of your book is from the field 
%	% "Psychology, Social Sciences"
%	%
%	%
%	% Monograph
%	\bibitem{psysoc-mono} Calfee, R.~C., \& Valencia, R.~R. (1991). \textit{APA guide to preparing manuscripts for journal publication.} Washington, DC: American Psychological Association.
%	%
%	% Online Document
%	\bibitem{psysoc-online} Dod, J. (1999). Effective substances. In: The dictionary of substances and their effects. Royal Society of Chemistry. Available via DIALOG. \\
%	\url{http://www.rsc.org/dose/Effective substances.} Cited 15 Jan 1999.
%	%
%	% Journal article
%	\bibitem{psysoc-journal} Harris, M., Karper, E., Stacks, G., Hoffman, D., DeNiro, R., Cruz, P., et al. (2001). Writing labs and the Hollywood connection. \textit{J Film} Writing, 44(3), 213--245.
%	%
%	% Contribution 
%	\bibitem{psysoc-contrib} O'Neil, J.~M., \& Egan, J. (1992). Men's and women's gender role journeys: Metaphor for healing, transition, and transformation. In B.~R. Wainrig (Ed.), \textit{Gender issues across the life cycle} (pp. 107--123). New York: Springer.
%	%
%	% Journal article by DOI
%	\bibitem{psysoc-DOI}Kreger, M., Brindis, C.D., Manuel, D.M., Sassoubre, L. (2007). Lessons learned in systems change initiatives: benchmarks and indicators. \textit{American Journal of Community Psychology}, doi: 10.1007/s10464-007-9108-14.
%	%
%	%
%	% Use the following syntax and markup for your references if 
%	% the subject of your book is from the field 
%	% "Humanities, Linguistics, Philosophy"
%	%
%	%\bigskip
%	%
%	% Journal article
%	\bibitem{humlinphil-journal} Alber John, Daniel C. O'Connell, and Sabine Kowal. 2002. Personal perspective in TV interviews. \textit{Pragmatics} 12:257--271
%	%
%	% Contribution 
%	\bibitem{humlinphil-contrib} Cameron, Deborah. 1997. Theoretical debates in feminist linguistics: Questions of sex and gender. In \textit{Gender and discourse}, ed. Ruth Wodak, 99--119. London: Sage Publications.
%	%
%	% Monograph
%	\bibitem{humlinphil-mono} Cameron, Deborah. 1985. \textit{Feminism and linguistic theory.} New York: St. Martin's Press.
%	%
%	% Online Document
%	\bibitem{humlinphil-online} Dod, Jake. 1999. Effective substances. In: The dictionary of substances and their effects. Royal Society of Chemistry. Available via DIALOG. \\
%	http://www.rsc.org/dose/title of subordinate document. Cited 15 Jan 1999
%	%
%	% Journal article by DOI
%	\bibitem{humlinphil-DOI} Suleiman, Camelia, Daniel C. O'Connell, and Sabine Kowal. 2002. `If you and I, if we, in this later day, lose that sacred fire...': Perspective in political interviews. \textit{Journal of Psycholinguistic Research}. doi: 10.1023/A:1015592129296.
%	%
%	%
%	%
%	%\bigskip
%	%
%	%
%	% Use the following syntax and markup for your references if 
%	% the subject of your book is from the field 
%	% "Computer Science, Economics, Engineering, Geosciences, Life Sciences"
%	%
%	%
%	% Contribution 
%	\bibitem{basic-contrib} Brown B, Aaron M (2001) The politics of nature. In: Smith J (ed) The rise of modern genomics, 3rd edn. Wiley, New York 
%	%
%	% Online Document
%	\bibitem{basic-online} Dod J (1999) Effective Substances. In: The dictionary of substances and their effects. Royal Society of Chemistry. Available via DIALOG. \\
%	\url{http://www.rsc.org/dose/title of subordinate document. Cited 15 Jan 1999}
%	%
%	% Journal article by DOI
%	\bibitem{basic-DOI} Slifka MK, Whitton JL (2000) Clinical implications of dysregulated cytokine production. J Mol Med, doi: 10.1007/s001090000086
%	%
%	% Journal article
%	\bibitem{basic-journal} Smith J, Jones M Jr, Houghton L et al (1999) Future of health insurance. N Engl J Med 965:325--329
%	%
%	% Monograph
%	\bibitem{basic-mono} South J, Blass B (2001) The future of modern genomics. Blackwell, London 
%	%
%	\end{thebibliography}


%In the case of delay differential equations with unbounded delay estimates on the rate of
%decay to a zero equilibrium have been obtained by many authors. Some representative papers
%include Krisztin [25], Kato [23], Diblík [14] and Diblík and R°užiˇcková [15], Cermak [12],
%and Haddock and Krisztin [18].
%An especially interesting equation which has received much attention is one with proportional
%delay, called the pantograph equation: fundamental work on the asymptotic behaviour
%dates back to Kato and McLeod [24], Fox et. al. [16], Ockendon and Tayler [35]. Complexvalued
%and finite dimensional treatments were considered by Carr and Dyson [10,11], while
%more recent treatments and generalisations include Iserles [19], Liu [29], and Makay and
%Terjeki [31].

\bibliographystyle{unsrt}
\begin{thebibliography}{10}

\bibitem{App2005} 
J. A. D. Appleby, Decay and growth rates of solutions of scalar stochastic delay differential equations with unbounded delay and state dependent noise, {\em Stoch. Dyn.} 5(2), 133--148, 2005.

\bibitem{AppBuck10}
J.~A.~D.~Appleby and E.~Buckwar, A constructive comparison technique for determining the asymptotic behaviour of linear functional differential equations with unbounded delay, {\em Differ. Equ. Dynam. Syst.} 18, No. 3, 271--301, 2010.

\bibitem{AppBuck2016} 
J.~A.~D.~Appleby and E.~Buckwar, Sufficient conditions for polynomial asymptotic behaviour of the stochastic pantograph equation. 
{\em Electronic Journal of Qualitative Theory of Differential Equations
Proc. 10th Coll. Qualitative Theory of Diff. Equ. }
2016, No. 2, 1--32. 

\bibitem{ACR:2011(DCDS)}
J.~A.~D. Appleby, J.~Cheng, and A.~Rodkina, 
Characterisation of the asymptotic behaviour of scalar linear differential equations with respect to a fading stochastic perturbation,
{\em Discrete Contin. Dyn. Syst.}, pages 79--90, 2011.

\bibitem{AL:2023(AppliedNumMath)}
J.~A.~D. Appleby and E.~Lawless, Mean square asymptotic stability characterisation of perturbed linear stochastic functional differential equations, {\em Applied Numerical Mathematics}, 200, 80--109, 2024. 

%\bibitem{AL:2023(AppliedMathLetters)}
%J.~A.~D. Appleby and E.~Lawless, Solution space characterisation of perturbed linear Volterra integrodifferential convolution equations: the ${L}^p$ case,
%\newblock {\em Applied Mathematics Letters}, 146:108825, 2023.

\bibitem{AL:2024Cesaro}
J.~A.~D. Appleby and E.~Lawless,
\newblock Solution space characterisation of perturbed linear functional and integrodifferential Volterra convolution equations: Ces\`aro limits, arXiv:2407.08416, 2024, 21pp. 

%\bibitem{AL:2024Lp}
%J.~A.~D. Appleby and E.~Lawless,
%Solution space characterisation of perturbed linear discrete and continuous stochastic Volterra convolution equations: the $\ell^p$ and $L^p$ cases, arXiv:2407.07767, 2024, 24pp.

\bibitem{AppLaw2024ODE}
J.~A.~D. Appleby and E.~Lawless, Weighted $L_\infty$ Asymptotic Characterisation of Perturbed Autonomous Linear Ordinary and Stochastic Differential Equations: Part I - ODEs, 2024, arXiv, 40pp.

%\bibitem{AppRey02}
%J.~A.~D. Appleby and D. W. Reynolds, Subexponential solutions of linear
%integro-differential equations and transient renewal equations,
%{\em Proc. R. Soc. Edinb.} 132A, 521--543, 2002.

\bibitem{BakBuck2000} 
C.~T.~H.~Baker and E.~Buckwar, Continuous $\theta$--methods for the stochastic pantograph equation, {\em Electron. Trans. Numer. Anal.} 11, 131--151, 2000.

\bibitem{BakBuck2005} 
C.~T.~H.~Baker and E.~Buckwar, Exponential stability in $p$-th mean of solutions, and of convergent Euler--type solutions, of stochastic delay differential equations. {\em J. Comput. Appl. Math.} 184, 404--427, 2005.

\bibitem{BGT}
N.~H.~Bingham, C.~M~.Goldie and J.~L.~Teugels, {\em Regular Variation},  Encyclopedia of Mathematics and its Applications. Cambridge University Press, 1991.

\bibitem{CarrD1} 
J.~Carr and J.~Dyson, The functional differential equation $y'(x) = ay(\lambda x) + by(x)$. {\em Proc. R. Soc. Edinb. A}, 74, 165--174 (1974/1975).

\bibitem{CarrDyson2} 
J.~Carr and J.~Dyson, The matrix functional differential equation $y'(x) = Ay(\lambda x) + By(x)$. {\em Proc. R. Soc. Edinb. A}, 75, 5--22 (1975/1976).

\bibitem{cermak} 
J.~\v Cerm\'ak, A change of variables in the asymptotic theory of differential equations with unbounded delay. {\em J. Comput. Appl. Math.} 143(1), 81--93, 2002.	

\bibitem{cer2022} 
\v Cerm\'ak, Jan and Nechv\'atal, Lud\v ek, On stability of linear differential equations with
commensurate delayed arguments, {\em Appl. Math. Lett.}, 125, 2022, Paper No. 107750, 8.

\bibitem{MR4278069} Caraballo, Tom\'as and Mchiri, Lassaad and Mohsen, Belfeki and
Rhaima, Mohamed, $p$-th moment exponential stability of neutral stochastic
pantograph differential equations with {M}arkovian switching, {\em Commun. Nonlinear Sci. Numer. Simul.}, 102, 2021, Paper No. 105916, 19.

\bibitem{ChanWilliams:1989}
T.~Chan and D.~Williams, An ``excursion'' approach to an annealing problem, {\em Math. Proc. Cambridge Philos. Soc.}, 105(1):169--176, 1989.

\bibitem{Diblik98}
J.~Dibl\'ik, Asymptotic representation of solutions of the equation $y'(t) = \beta(t)[y(t) - y(t - \tau(t))]$.
{\em J. Math. Anal. Appl.} 217, 200--215, 1998.

\bibitem{DibRuz} 
J.~Dibl\'ik and M. R\r{u}\v{z}i\v{c}kov\'a, Convergence of the solutions of the equation $y'(t) = \beta(t)[y(t -\delta)- y(t -\tau)]$
in the critical case. {\em J. Math. Anal. Appl.} 331, 1361--1370, 2007.

\bibitem{Fan} 
Z.~Fan, M.~Liu and W.~Cao, Existence and uniqueness of the solutions and convergence
of semi-implicit Euler methods for stochastic pantograph equations, {\em J. Math. Anal. Appl.} 325, No. 2, 1142--1159, 2007.  

\bibitem{Foxetal} 
L.~Fox, D.~F.~Mayers, J.~R.~Ockendon, and A.~B.~Tayler, On a functional differential equation. {\em J. Inst. Math. Appl.} 8, 271--307, 1971.

\bibitem{GLS}
G.~Gripenberg, S.-O. Londen, and O.~Staffans.
\newblock {\em Volterra Integral and Functional Equations}.
\newblock Encyclopedia of Mathematics and its Applications. Cambridge University Press, 1990.

\bibitem{MR3927238}
P.~Guo and C.-J.~Li, Almost sure stability with general decay rate of exact and
numerical solutions for stochastic pantograph differential
equations, {\em Numer. Algorithms}, 80, 4, 1391--1411, 2019.

\bibitem{HaddKris} 
J.~R.~Haddock and T.~Krisztin, On the rate of decay of solutions of functional differential equations with infinite delay. {\em Nonlinear Anal.} 10(8), 727--742, 1986

\bibitem{HW55}
P.~Hartman and A.~Wintner,  Asymptotic integration of ordinary nonlinear differential equations, {\em Amer. J. Math.} 77, 692--724, 1955.

\bibitem{Hart2002} Ph.~Hartman, {\em Ordinary differential equations}, second ed., SIAM, Philadelphia, 2002.

\bibitem{Huetal}
Y.~Hu, F.~Wu and C.~Huang, Stochastic stability of a class of unbounded delay neutral stochastic differential equations with general decay rate, {\em Internat. J. Systems Sci.} 43(2012), No. 2, 308--318, 2012

\bibitem{Huang}
Y.~Huang, Z.~Yan and H.~Zheng, Stochastic $H_\infty$ filtering for pantograph systems with state dependent noise, in: 2012 IEEE International Conference on Automation and Logistics (ICAL),
5–17 August 2012, Zhengzhou, 2012, pp. 601--605.

\bibitem{Iser} 
A.~Iserles, On the generalized pantograph functional--differential equation. 
{\em Eur. J. Appl. Math.} 4, 1--38, 1993.

\bibitem{Jiang}
F.~Jiang, H. Yang, and S.~Wang, Asymptotic stability of stochastic pantograph differential equations with Markovian switching, {\em J. Nonlinear Anal. Optim.} 1, No. 1, 9--16, 2010.

\bibitem{KS} 
I.~Karatzas and S.~E.~Shreve, {\em Brownian  Motion and Stochastic Calculus}, volume 113 of
Graduate Texts in Mathematics, Springer, New York, 1991.

%[8] D. Revuz and M. Yor, Continuous Martingales and Brownian Motion. Third edition,
%Springer, New York, 1999.

\bibitem{KatoMc} T.~Kato and J.~B.~ McLeod, The functional--differential equation 
$y'(x) = ay(\lambda x) + by(x)$. {\em Bull. Am. Math.
	Soc.} 77, 891--937,  1971.

\bibitem{Kris88} T.~Krisztin, Asymptotic estimation for functional differential equations via Lyapunov functions, 
{\em Colloquia Math. Soc. J\'anos Bolyai I} 109, 365--376, 1988.

\bibitem{MR4534554}
A.~Lachouri, M.~E.~Samei, and A.~Ardjouni, Existence and stability analysis for a class of fractional
pantograph {$q$}-difference equations with nonlocal boundary
conditions, {\em Bound. Value Probl.}, Paper No. 2, 20, 2023.

\bibitem{Liu} Y.~Liu, On functional differential equations with proportional delays. Ph.D. Thesis, Department of Mathematics, Cambridge University, Cambridge.

\bibitem{LiuChen}
X.~Liu and T.~Chen, Robust $\mu$--stability for uncertain stochastic neural networks with unbounded time--varying delays, {\em Phys. A.} 387, No. 12, 2952--2962, 2008. 

\bibitem{MakTerj}
G.~Makay and J.~Terj\'eki, On the asymptotic behavior of the pantograph equations. {\em Electron. J. Qual. Theory Differ. Equ.} 1998(2), 1--12, 1998.

\bibitem{MR4053553} 
W.~Mao, L.~Hu, and X.~Mao, The asymptotic stability of hybrid stochastic systems with
pantograph delay and non-Gaussian L\'evy noise, {\em J. Franklin Inst.}, 357, 2, 1174--1198, 2020.

\bibitem{Mao1992a} 
X.~Mao, Almost sure polynomial stability for a class of stochastic differential equations,
{\em Quart. J. Math. Oxford Ser.}, (2) 43, No. 171, 339--348, 1992.

\bibitem{Mao1992b} X.~Mao, Polynomial stability for perturbed stochastic differential equations with respect
to semimartingales, {\em Stochastic Process. Appl.},  41, No. 1, 101--116, 1992.

\bibitem{MaoBK}
X.~Mao, {\em Stochastic differential equations and their applications}, Horwood Publishing Limited, Chichester, 2008. 

\bibitem{Meng1} 
X.~Meng, S.~Hu and P.~Wu, Pathwise estimation of stochastic differential equations with unbounded delay and its application to stochastic pantograph equations, {\em Acta Appl. Math.} 113, No. 2, 231--246, 2011. 

\bibitem{Meng2}
X.~Meng and B.~Yin, On the general decay stability of stochastic differential equations with unbounded delay, {\em J. Korean Math. Soc.} 49, No. 3, 515--536, 2012. 

\bibitem{Milo}
M.~Milosevic, Existence, uniqueness, almost sure polynomial stability of solution to a class of highly
nonlinear pantograph stochastic differential equations and the Euler--Maruyama approximation. {\em Appl.
	Math. Comput.} 237, 672--685, 2014.

\bibitem{OckTay}
J.~R.~Ockendon and A.~B.~Tayler, he dynamics of a current collection system for an electric locomotive. {\em Proc. R. Soc. A} 322, 447--468, 1971.

\bibitem{PJ2012}
G.~Pavlovic and S.~Jankovic, Razumikhin--type theorems on general decay stability of stochastic
functional differential equations with infinite delay, {\em J. Comput. Appl. Math.} 236(7), 1679--1690, 2012.

\bibitem{Pituk2002}
M.~Pituk, The Hartman--Wintner theorem for functional differential equations, {\em J. Differential
	Equations}, 155, no. 1, 1--16, 1999. 

\bibitem{Pituk2006}
M.~Pituk, A Perron type theorem for functional differential equations, {\em J. Math. Anal. Appl.}, 316 (1), 24--41, 2006.

\bibitem{MR4585696}
A.~Siva Ranjani, M.~Suvinthra and K.~Balachandran, Large deviations for stochastic pantograph integrodifferential equation, {\em Filomat}, 37, 20, 6751--6766, 2023.
%	FJOURNAL = {Univerzitet u Ni\v su. Prirodno-Matemati\v cki Fakultet.
	%	Filomat},
%	VOLUME = {37},
%	YEAR = {2023},
%	NUMBER = {20},
%	PAGES = {6751--6766}
%	ISSN = {0354-5180,2406-0933},
%	MRCLASS = {60H10 (45J05 60F10)},
%	MRNUMBER = {4585696},
%}

%\bibitem{MR4534554}
%	Lachouri, Adel and Samei, Mohammad Esmael and Ardjouni,
%		Abdelouaheb, Existence and stability analysis for a class of fractional
%		pantograph {$q$}-difference equations with nonlocal boundary
%		conditions, Bound. Value Probl., Paper No. 2, 20, 2023.
%	FJOURNAL = {Boundary Value Problems},
%	YEAR = {2023},
%	PAGES = {Paper No. 2, 20},
%	ISSN = {1687-2762,1687-2770},
%	MRCLASS = {39A13 (26A33 34A08 34B10)},
%	MRNUMBER = {4534554},
%	MRREVIEWER = {Richard\ L.\ Rubin},
%	DOI = {10.1186/s13661-022-01691-1},
%	URL = {https://doi.org/10.1186/s13661-022-01691-1},
%}

%\bibitem{RevYor}
%D.~Revuz and M.~Yor, {\em Continuous Martingales and Brownian Motion}, Third edition, Springer, New York, 1999.

\bibitem{Shiryaev}
A.~N.~Shiryaev, {\em Probability} (2nd edition), Springer, Berlin, 1996.

\bibitem{SY67a}
A.~Strauss and J.~A.~Yorke, Perturbation theorems for ordinary differential equations, {\em Journal of Differential Equations}, 3:15--30, 1967.

\bibitem{SY67b}
A.~Strauss and J.~A.~Yorke, On asymptotically autonomous differential equations,
{\em Mathematical Systems Theory}, 1:175--182, 1967.

\bibitem{WuHu2012} F.~Wu and S.~Hu, Razumikhin--type theorems on general decay stability and robustness for stochastic
functional differential equations. {\em Int. J. Robust Nonlinear Control} 22(7), 763--777, 2012.

\bibitem{MR4439350}
H.~Wu, J.~Hu, and C.~Yuan, Stability of numerical solution to pantograph stochastic functional differential equations, {\em Appl. Math. Comput.}, 431, Paper No. 127326, 13, 2022.

\bibitem{YanHuang}
Z.~Yan and Y.~Huang, Robust $H_\infty$ filter design for It\^o stochastic pantograph systems, {\em Mathe. Probl. Eng.} (2013), Art. ID 747890, 8 pp., 2013.

\bibitem{Yangetc}
H.~Yang, F.~Jiang, and Y.~Jiang, Robust stability of stochastic pantograph differential equations Markovian switching, in: International Conference on Computer Application and System Modelling (ICCASM), October 22–24, 2010, Shanxi, Taiyuan, Vol. 5, pp. 541--545, 2010. 
\end{thebibliography}


%\bibitem{MR4585696,
%	AUTHOR = {Siva Ranjani, A. and Suvinthra, M. and Balachandran, K.},
%	TITLE = {Large deviations for stochastic pantograph integrodifferential
%		equation},
%	JOURNAL = {Filomat},
%%	FJOURNAL = {Univerzitet u Ni\v su. Prirodno-Matemati\v cki Fakultet.
%		Filomat},
%	VOLUME = {37},
%	YEAR = {2023},
%	NUMBER = {20},
%	PAGES = {6751--6766},
%%	ISSN = {0354-5180,2406-0933},
%%	MRCLASS = {60H10 (45J05 60F10)},
%%	MRNUMBER = {4585696},
%}

\end{document}

